\documentclass[reqno,final,11pt]{amsart}
%% Packages
\usepackage{natbib}  %Nature-like bibliography
\usepackage{fancyhdr} %Headers
\usepackage{color} %Color definition
\usepackage{hyperref} %Internal and external links
\usepackage{graphicx} %Graphics inclusion
\usepackage{geometry}
\geometry{
  top=1in,
  inner=1in,
  outer=1in,
  bottom=1in,
  headheight=3ex,
  headsep=2ex,
}
\usepackage[utf8]{inputenc}
\usepackage[T1]{fontenc}

%%%%%%%%%%%%%%%%%%%%%%%%%%%%%%%%%%%%%
%% Author packages
%%%%%%%%%%%%%%%%%%%%%%%%%%%%%%%%%%%%%
\usepackage{amssymb}
\usepackage{cleveref}
\emergencystretch=3em

%%%%%%%%%%%%%%%%%%%%%%%%%%%%%%%%%%%%%


%% Definition of colors for links
\definecolor{aleacolor}{rgb}{0.16,0.59,0.78}

%% Settings for hyperref package
\hypersetup{
breaklinks,
colorlinks=true,
linkcolor=aleacolor,
urlcolor=aleacolor,
citecolor=aleacolor}

%% Settings for fancyhdr package
%\renewcommand{\headheight}{11pt}
\pagestyle{fancy} \fancyhf{} \fancyhead[RO,LE]{\small\thepage}
\fancyhead[RE]{\small\shortauthors} \fancyhead[LO]{\small\shorttitle}

%% Settings for natbib package
\renewcommand{\cite}{\citet}

%% Setting the theorem-like environments
\theoremstyle{plain}
\newtheorem{theorem}{Theorem}[section]
\newtheorem{proposition}[theorem]{Proposition}
\newtheorem{lemma}[theorem]{Lemma}
\newtheorem{corollary}[theorem]{Corollary}
\newtheorem{conjecture}[theorem]{Conjecture}
\theoremstyle{definition}
\newtheorem{definition}[theorem]{Definition}
\newtheorem{assumption}[theorem]{Assumption}
\newtheorem{convention}[theorem]{Convention}
\newtheorem{hypothesis}[theorem]{Condition}
\newtheorem{problem}[theorem]{Open problem}
\theoremstyle{remark}
\newtheorem{remark}[theorem]{Remark}
\newtheorem{example}[theorem]{Example}

%% Numbering
\makeatletter \@addtoreset{equation}{section} \makeatother
\renewcommand\theequation{\thesection.\arabic{equation}}
\renewcommand\thefigure{\thesection.\arabic{figure}}
\renewcommand\thetable{\thesection.\arabic{table}}

%% First page header, modify the link w.r.t. volume number
\newcommand{\aleaIndex}[1]{\href{http://alea.impa.br/english/index_v23.htm}{\bf 23}}
\newcommand{\aleaDOI}[1]{\href{https://doi.org/10.30757/ALEA.v23-}{DOI: 10.30757/ALEA.v23-}}
\eheader{ALEA, Lat. Am. J. Probab. Math. Stat.}{\aleaIndex{23}}{2026}{1}{4}{\aleaDOI}
%% Uncomment the following line to include Alea logo
\elogo{\parbox[c]{180pt}{\includegraphics[width=170pt]{logo.pdf}}}

%%%%%%%%%%%%%%%%%%%%%%%%%%%%%%%%%%%%%
%% Author commands and definitions
%%%%%%%%%%%%%%%%%%%%%%%%%%%%%%%%%%%%%
\newcommand{\R}{\mathbb R}
\newcommand{\N}{\mathbb N}
\newcommand{\E}{\mathbb E}
\newcommand{\Prob}{\mathbb P}
\newcommand{\Dom}{\mathrm{Dom}}
\newcommand{\cA}{\mathcal A}
\newcommand{\cD}{\mathcal D}
\newcommand{\cE}{\mathcal E}
\newcommand{\cF}{\mathcal F}
\newcommand{\cN}{\mathcal N}
\newcommand{\cU}{\mathcal U}
\newcommand{\Gam}{\Gamma}
\newcommand{\GamT}{\Gamma^{(2)}_{\theta}}
\newcommand{\GamN}{\Gamma^{(2)}_{\cN_\theta}}
\newcommand{\Lth}[1]{L^{(#1)}_{\theta}}
\newcommand{\DN}{D_{\cN_\theta}}
\newcommand{\divN}{\delta_{\cN_\theta}}
\newcommand{\Pth}{\Prob_{\theta}}
\newcommand{\Hone}{\mathcal H_{1}}
\newcommand{\Aone}{\cA^{\mathrm{1t}}_{\mathrm{cyl}}}
\newcommand{\Apair}{\cA^{\le 2}_{\mathrm{cyl}}}
\newcommand{\Apairflat}{\cA^{\le 2,\flat}_{\mathrm{cyl}}}
\newcommand{\muD}{\mu_{\Delta}}
\newcommand{\cth}{c_{\theta}}
\newcommand{\Ecyl}{\cE^{\mathrm{cyl}}_{\theta}}
\newcommand{\CKC}{(\mathrm{CKC}_{2})}
\newcommand{\PiN}{\Pi_{\cN_\theta}}
\newcommand{\COcap}{(\mathrm{CO\text{-}Cap})}
\newcommand{\COmix}{(\mathrm{CO\text{-}Mix})}
\newcommand{\Ltwoz}{L^2_0}

%%%%%%%%%%%%%%%%%%%%%%%%%%%%%%%%%%%%%
\date{DATE; accepted DATE; published under license
\href{https://creativecommons.org/licenses/by/4.0/}{\textup{CC BY 4.0}}} %%% to be modified by the editor

\begin{document}

\title[Black Clark--Ocone and Operator-It\^o for Sticky Flows]{Pairwise Black
Clark--Ocone and Operator-It\^o Formulae for Le Jan--Raimond Sticky Flows,
and the Sticky--Coalescing Dichotomy}

\author{Ramiro Fontes}
\address{Quijotic Research}
\email{ramirofontes@gmail.com}
\urladdr{\href{https://www.quijotic.com}{www.quijotic.com}}

\thanks{The author declares that no funds, grants, or other support were
received during the preparation of this manuscript.}
\subjclass[2020]{Primary 60H07; Secondary 60J46, 60J55, 60G20}
\keywords{Black noise, sticky flow, stochastic flow of kernels,
Clark--Ocone representation, Dirichlet form, operator It\^o formula,
coalescence}

\begin{abstract}
The noise of the Le Jan--Raimond Brownian sticky flow is black: its first
chaos is trivial, so functionals of the flow admit no martingale
representation and no stochastic-integral It\^o formula. We develop the
calculus that survives. The operator derivative $D$ on a carr\'e-du-champ
energy space comes first, and the divergence is defined by adjointness,
\[
  \delta\;:=\;D^{*},
\]
on fixed one- and two-point sectors of the flow. Three results follow. First,
a boundary dichotomy: the deletion contraction that the construction requires
holds in the reflecting phase under a first-order condition
$(\mathrm{CKC}_2^{\flat})$ on the two-point law and fails at the coalescing
endpoint, so the boundary type rather than blackness decides the calculus; in
the reflecting phase the diagonal-flat pairwise sector then carries closed
operator pairs and a pairwise black integration-by-parts identity. Second, a
Clark--Ocone--Green representation
\[
  F-\E F\;=\;\delta_{\cN_\theta}\bigl(\PiN D_{\cN_\theta}F\bigr),
\]
in which the representer is a minimal-norm Green solution rather than a
predictable projection, there being no first chaos to project onto. The
two-point case rests on a sectorwise Poincar\'e inequality; on the
autonomous gap sector the representation holds with
explicit constants and fails at the coalescing endpoint, where the
coalescence indicator is a nonconstant energy-null functional orthogonal to
every divergence. Third, a partitionwise dynamic operator-It\^o formula: the
increment is an adapted sum of slicewise divergences plus the generator
drift, exact at every partition, with the Dynkin martingale recovered in the
mesh limit. Blackness forbids a driving noise but not adapted increments, so
the formula is driverless; on the gap sector it requires nothing beyond the
model itself.
\end{abstract}

\maketitle

%==============================================================
\section{Introduction}\label{sec:intro}
%==============================================================

A stochastic flow of kernels $K=(K_{s,t})_{s\le t}$ on the line (or
circle) is a random evolution of probability kernels; the Le Jan--Raimond
(LJR) construction~\citep{LJR-FCN} produces such a flow from a consistent
family of $n$-point Feller semigroups $(\Prob^{(n)}_t)_{n\ge1}$, the
$n$-point motion being ``the motion of $n$ unsecable points thrown into the
fluid''~\citep[\S1]{LJR-FCN}. The Brownian sticky flow is the flow
whose one-point motion is Brownian and whose two-point motion is sticky on
the diagonal: two points started apart can meet and, having met, can
separate again (the intermediate or reflecting phase), in
contrast with the coalescing (Arratia) flow where points that meet are
absorbed together forever.

The flow noise $\cN_\theta$, the continuous product of $\sigma$-fields
$(\cF^{\theta}_{s,t})_{s\le t}$ generated by $K$, is a black noise
in the sense of Tsirelson~\citep{Tsirelson-StFlour,Tsirelson-Survey}: its first
chaos (the space of decomposable, i.e.
additive-across-time, square-integrable functionals) is trivial. Le Jan and
Raimond prove exactly this~\citep{LJR-circle-noise}: the noise of a Brownian
sticky flow is black. This is what obstructs the classical
martingale-representation route to an It\^o formula, and it is a stronger failure
than the usual loss of semimartingale calculus: it removes the first-chaos
representation itself, not only the finite predictable bracket. There is no
internal martingale representation from which to read off a derivative, so the
compression-from-a-source route does not apply. What remains is the
reversal: the operator
derivative $D$ on a carr\'e-du-champ energy space comes first,
\begin{equation}\label{eq:reversal}
  \delta\;:=\;D^{*},
\end{equation}
and build the calculus downstream of \eqref{eq:reversal} from a
$\Gamma$-tuple extracted directly from the LJR two-point generator. Working
pairwise, at the level of one- and two-point observables, isolates the single
analytic question that the reversal leaves open: whether the pairwise energy
contracts under deletion of a coordinate. We phrase this as a
conditional-expectation energy contraction and reduce it to an explicit inequality
$\CKC$ (\cref{def:ckc}) on the structural scalars
$(q,q',p,\partial_\xi p,\cth,\muD)$ of the two-point time-$T$ law. The reduction
runs through an amended gradient identity for the deletion map
(\cref{prop:amended-star}): because the diagonal atom of the sticky kernel moves
with the retained coordinate, the naive score identity omits a term
$q(\xi)\,\partial_\eta\Phi(\xi,\xi)$, and carrying that term is what
makes the reduction correct.

\subsection*{Three results, and the road map}
The paper proves three things, in this order.
\begin{enumerate}
\item An energy-first construction over a black noise
(\cref{sec:main}): closed sectorwise pairs $D$, $\delta=D^{*}$, defined
directly since first-chaos representation is unavailable, and a pairwise
black integration-by-parts identity (\cref{thm:pairwise-ito}), with the
sticky--coalescing dichotomy locating exactly where the construction lives
(\cref{thm:unobstruction,thm:wall-endpoint,thm:discriminant}).
\item A pairwise Clark--Ocone--Green representation
(\cref{sec:clarkocone}): $F-\E F=\divN(\PiN\DN F)$
(\cref{thm:clarkocone}), together with the gap-sector dichotomy
(\cref{thm:gap-co}) and the coalescing obstruction
(\cref{prop:co-failure}).
\item A partitionwise driverless It\^o formula (\cref{sec:dynamic}):
the increment as an adapted sum of slicewise divergences
(\cref{thm:dynamic}), not an integral against a first-chaos driver
(\cref{def:driverless}).
\end{enumerate}
The dependency chain:
\begin{align*}
  (\mathrm{CKC}_2^{\flat})
  &\;\Longrightarrow\;
  \text{integration by parts (\cref{thm:pairwise-ito})},\\
  \text{sector Poincar\'e inequality}
  &\;\Longrightarrow\;
  \text{Clark--Ocone (\cref{thm:clarkocone})},\\
  \text{slicewise Clark--Ocone}
  &\;\Longrightarrow\;
  \text{dynamic formula (\cref{thm:dynamic})}.
\end{align*}
The remainder of this introduction describes the route through the two-point
structure; the reader oriented by this summary may skip to \cref{sec:flow}.

The inequality $\CKC$ is then read off the gap. The two-point gap $g=\eta-\xi$ is
a scalar sticky Brownian motion (\cref{lem:gap-sticky}); its time-$T$ law splits
into a diagonal atom and a diffuse density $p$, and the one number that controls
$\CKC$ is the near-diagonal limit of $p$. Tied to the LJR boundary classification,
this limit produces a sharp dichotomy (\cref{cor:dichotomy}), the core of
the first main result: at every
reflecting-regular boundary the value-trace, Poincar\'e and Fisher constants
$\tau_{\mathrm{wt}},C^p_P,I_p$ are finite and the obstruction to the contraction
is absent (\cref{thm:unobstruction}), whereas at the coalescing endpoint the
density vanishes, these constants diverge, and a difference-direction probe
defeats the contraction (\cref{thm:wall-endpoint}). Both phases carry black
noises, so what the boundary type discriminates is not blackness but exactly
this capacity obstruction (\cref{thm:discriminant}).

One element escapes the first-order analysis. The moving-atom term contributes
a derivative trace $\tau_{\mathrm{at}}$ whose finiteness is a second-order ($H^2$-type) condition, not a consequence of reflecting regularity.
On the diagonal-flat sub-core, where the difference-direction derivative vanishes
at coincidence, this term rotates into the center-of-mass direction and is
absorbed by the diagonal energy (\cref{cor:neumann-subcore}), so the trace drops
out. There, under the reduced contraction condition $(\mathrm{CKC}_2^{\flat})$,
the reflecting phase carries a single-time pairwise
black integration-by-parts identity (\cref{thm:pairwise-ito}), whose centered
divergence is the intrinsic energy-space object and never a first-chaos integral,
with pairing governed by the diagonal carr\'e $\GamT$. The closed-pair
construction it rests on is developed self-contained and applied on each fixed
one- or two-point sector (\cref{app:closedpair,app:pairclos}).

The second main result is the Clark--Ocone upgrade
(\cref{sec:clarkocone}). The integration-by-parts identity is adjointness, not
yet representation; the classical upgrade projects $DF$ onto predictable
integrands and is doubly unavailable here, blackness having removed the
first-chaos target and the single-time sector carrying no filtration. The
energy route survives: the upgrade is equivalent to a sectorwise Poincar\'e
inequality, which we prove from the paper's own trace bounds together with a
uniformization of the two-point value capacity ($\COcap$, or the mixed-mean
bound $\COmix$), with the deletion contraction $(\mathrm{CKC}_2^{\flat})$ paying
the two-point step and the classical Gaussian Poincar\'e inequality the
one-point step. The result (\cref{thm:clarkocone}) is a pairwise black
Clark--Ocone--Green representation $F-\E F=\divN(\PiN\DN F)$, the representer
being the minimal-norm Green solution in the energy module. Two statements
frame it from \cref{conv:LJR} alone: on the autonomous gap sector the
representation holds with explicit constants, in a sharp
reflecting/coalescing dichotomy (\cref{thm:gap-co}); at the coalescing
endpoint it fails outright, the coalescence indicator being a nonconstant
energy-null functional that no divergence represents (\cref{prop:co-failure}).
Whether the reflecting phase alone carries the full two-point representation
remains open.

The third main result is dynamic (\cref{sec:dynamic}). The Dynkin
comparison (\cref{prop:dynkin}) identifies what a dynamic formula must
disintegrate, and the adapted-telescope theorem (\cref{thm:dynamic}) carries
the disintegration out at the level of partitions, exactly: the increment
$f(Z_T)-f(Z_0)$ equals an adapted sum of conditional single-slice energy
divergences plus the generator drift, at every partition, with the Dynkin
martingale recovered in the mesh limit. Blackness rules out slice-additive
functionals, not adapted increments, so the formula is driverless rather than
impossible (\cref{def:driverless}): state-indexed conditional divergences
replace a driver, and an orthogonal-increment identity replaces the It\^o
isometry. On the gap sector the dynamic formula needs nothing further (\cref{cor:dynamic-gap}),
the Wentzell drift being clocked by the sticky occupation time, whose
near-total tail is computed in closed form (\cref{prop:occupation}).

Supporting results sharpen the boundaries of the three:
the center-of-mass reduction \cref{prop:com-half} compresses the two-point
Poincar\'e question into $\COmix$ and identifies the essential-supremum form
of $\COcap$ as the wrong uniformity (\cref{rem:two-well}); the moving-atom
trace condition becomes a theorem on the second-order sub-core by an
elementary Agmon bound, isolating the full reflecting case in a single
intertwining question; and the codimension-two polarity route to a
three-point failure is blocked by the one-dimensional capacity of the
pairwise strata (\cref{sec:sharpness}).

Beyond the flat sub-core the moving-atom trace is higher-order, and a
full reflecting theorem turns on a semigroup-intertwining question for the
time-$T$-law operators (\cref{rem:dtrace-gamma2}). That question, the three- and
higher-point contraction, and the multi-time lift are the natural next steps;
they are collected in \cref{sec:openproblems}.

Throughout, ``pairwise'' refers to observations of at most two points. This
sub-algebra is stable under the coordinate deletions $2\to1\to0$, which is what
keeps the pairwise sector self-contained and lets the reduction proceed without
reference to the higher-point laws.

%==============================================================
\section{The sticky flow, its noise, and the boundary trichotomy}\label{sec:flow}
%==============================================================

Fix a stickiness parameter $\theta\in[0,\infty]$. The LJR Brownian sticky
flow is the stochastic flow of kernels $K^{\theta}=(K^{\theta}_{s,t})_{s\le t}$
whose law is determined by a consistent family of $n$-point Feller
semigroups $(\Prob^{(n)}_t)_{n\ge1}$~\citep[Thm.~2.1, \S2]{LJR-FCN}; the
$n$-point motion $X^{(n)}=(X^{x_1}_t,\dots,X^{x_n}_t)_t$ is the Markov
process on $\R^n$ associated with $\Prob^{(n)}_t$. The one-point motion is
Brownian~\citep{LJR-circle-noise}. We drop the superscript $\theta$ when fixed
and write $\Pth$ for the law of the flow, $\E$ for $\E_{\Pth}$, and
$\cN_\theta=(\cF^{\theta}_{s,t})_{s\le t}$ for the flow noise.

The starting point is that this noise is black.

\begin{theorem}[Le Jan--Raimond {\citep{LJR-circle-noise}}]\label{thm:black}
The noise $\cN_\theta$ of a Brownian sticky flow is black: its first chaos
(the space of centered linear, i.e. additive-across-time, $L^2$
representations) is trivial.
\end{theorem}
We use \cref{thm:black} as an input. Its mechanism is the
same-law-until-$\tau$ estimate~\citep{LJR-circle-noise}: for the $2d$-point
joint motion and two independent $d$-point motions, the processes have the
same law up to the first meeting time
$\tau=\inf_{i,j}\inf\{t:X^i_t=Y^j_t\}$, and
$\Prob[\tau<\varepsilon]\le C\sqrt{\varepsilon}$ (one Brownian gap hitting
zero, by reflection), so the joint-minus-independent difference of any
second moment of a difference functional $(g(X)-g(Y))^2$ is
$O(\varepsilon^{3/2})=o(\varepsilon)$. Summed over a dyadic time partition
this gives zero quadratic variation of the first-chaos projection, hence
triviality of the first chaos (Tsirelson's
criterion~\citep{Tsirelson-StFlour,Tsirelson-Survey}). One domain convention
should be flagged once: \citep{LJR-circle-noise} states the theorem for the
sticky flow on the circle, whereas \cref{conv:LJR} and the analysis below live
on the line. The mechanism just described transfers verbatim: it compares the
joint and the independent $n$-point motions only up to the first meeting time
and uses a single Brownian-gap hitting estimate, both local in space and time
and insensitive to the global geometry (the line being the universal cover of
the circle, with the $n$-point sticky semigroups those of the consistent
family~\citep{LJR-FCN}); we use the line statement throughout. Two features of this
mechanism return below: it lives on difference functionals at
coincidence, and it integrates against an absolutely continuous time
marginal that does not charge the diagonal.

We now fix the two-point structure of the flow as a convention.

\begin{convention}[LJR two-point sticky generator]
\label{conv:LJR}
The two-point generator $\Lth{2}$ on $L^2(\R^2,dx\,dy)$ is characterized by
four data inherited from~\citep{LJR-FCN}:
\textnormal{(a)} a stickiness parameter $\theta\in[0,\infty]$, one endpoint
the independent-Brownian flow, the other the coalescing flow;
\textnormal{(b)} a diagonal coefficient $\cth\ge0$ with
$\cth=0$ at the independent endpoint and $\cth=c_{\mathrm{coal}}>0$ at the
coalescing endpoint;
\textnormal{(c)} a Radon diagonal reference measure $\muD$ on
$\Delta\simeq\R$ (Lebesgue, or the local-time-at-$0$ measure of $Y-X$);
\textnormal{(d)} a Wentzell-type boundary condition on $\Delta$
characterized by $\theta$, with off-diagonal action $\tfrac12\Delta$, the
exact statement delegated to~\citep{LJR-FCN}.
\end{convention}

The carr\'e seed lives on the diagonal. The two-point product defect
$\mathfrak D^{(2)}_\theta(f,g):=\Lth{2}(f\otimes g)-(\tfrac12 f'')\otimes g
-f\otimes(\tfrac12 g'')$ is supported on $\Delta$ and equals, on the diagonal,
$\cth f'(x)g'(x)\,\muD(dx)$ (the off-diagonal part of $\Lth{2}$ is
$\tfrac12\Delta$, so the defect is diagonal-supported; the diagonal density is
fixed by the Wentzell condition of \cref{conv:LJR}(d)). The diagonal
carr\'e-du-champ seed is
\begin{equation}\label{eq:gamma-seed}
  \GamT(f,g)\;=\;\tfrac12\,\mathfrak D^{(2)}_\theta(f,g)
  \;=\;\tfrac{\cth}{2}\,f'(x)g'(x)\,\muD(dx)\quad\text{on }\Delta,
\end{equation}
the diagonal seed, supported on the coincidence stratum.

Everything that follows turns on the Feller boundary type of the two-point
distance process, so we name it now and quote the LJR classification verbatim. Let
$d_t=d(X_t,Y_t)$ be the two-point distance process and
$\Prob_r,H_r$ as in~\citep[\S9]{LJR-FCN}.

\begin{definition}[Boundary trichotomy]\label{def:trichotomy}
The flow is in the
\textnormal{(N)} non-coalescing phase, the
\textnormal{(R)} reflecting (intermediate, sticky) phase, or the
\textnormal{(C)} coalescing phase, according to the boundary type of
$0$ for $d_t$.
\end{definition}

\noindent The two cases that matter here are the two-point boundary
classification of~\citep[Theorem~1.1 and~\S9]{LJR-FCN} (we cite the statement by
its place in the classification theorem; section numbers refer to the published
Ann.\ Probab.\ version):
\begin{quote}
\textnormal{(C, exit/absorbing, coalescing case)} ``\,$\nu^{s}$ defines a
coalescing flow of maps if and only if $0$ is a closed exit boundary point,
i.e. if $\exists\, r>0,\ \Prob_r[H_0<\infty]>0$ and $\forall\, r>0,\
\Prob_0[H_r<\infty]=0$.''

\textnormal{(R, reflecting regular, intermediate case)} ``\,$\nu^{s}$
defines a turbulent flow with hitting (i.e. such that the 2-point motion
starting outside of the diagonal hits the diagonal with a positive
probability) if and only if $0$ is a reflecting regular boundary point, i.e. if
$\exists\, r>0,\ \Prob_r[H_0<\infty]>0$ and $\exists\, r>0,\
\Prob_0[H_r<\infty]>0$.''
\end{quote}
The distinguishing clause is $\Prob_0[H_r<\infty]$: in phase (R) the process
leaves the diagonal (positive probability of separating after
meeting); in phase (C) it does not. This single clause is what discriminates the
identified capacity obstruction (\cref{thm:discriminant}).

%==============================================================
\section{The pairwise lifted carr\'e-du-champ structure}\label{sec:recap}
%==============================================================

We introduce the pairwise lifted carr\'e-du-champ structure used in the sequel;
the closed-pair and closability machinery is established in
\cref{app:closedpair,app:pairclos}. Everything but the closability axiom
(FormClos) is algebraic; per fixed $\le2$-point sector that axiom is the settled
first-order closability (P1) (\cref{thm:sectorclos}), and the remaining pairwise
content is the $2\to1$ deletion-compatibility (P3), which the present paper
attacks pairwise (\cref{thm:deletion}).

Fix a horizon $T>0$. For $\mathbf x=(x_1,\dots,x_m)\in\R^m$ and
$\Phi\in C_b^\infty(\R^m)$, the single-time cylindrical observable is
$F_{\mathbf x,T,\Phi}=\Phi(X^{x_1}_T,\dots,X^{x_m}_T)$. The single-time
cylindrical algebra $\Aone$ is the unital algebra they generate (products close
by tuple concatenation, with the number of points unrestricted); it is dense in
$\Hone:=L^2(\Omega,\cF_T^{\mathrm{1t}},\Pth)$, the time-$T$ sub-Hilbert space.
The pairwise object $\Apair$ used below ($\le2$ points) is, by contrast, only a
vector sub-core, not a subalgebra (a product of two-point functionals
leaves the pairwise sector); the unital closed-pair construction is therefore
applied per fixed sector (\cref{rem:sectorwise}), not to $\Apair$ as an algebra. The lifted carr\'e-du-champ is the pullback
\begin{equation}\label{eq:lifted-gamma}
  \GamN(F,G)\;=\;\Gamma^{(m)}_\theta(\Phi,\Psi)\bigl(X^{x_1}_T,\dots,X^{x_m}_T\bigr),
  \qquad F=\Phi(\cdot),\ G=\Psi(\cdot),
\end{equation}
where $\Gamma^{(m)}_\theta$ is the spatial $m$-point carr\'e of $\Lth{m}$.
Its structural form (square-of-sum), abstracted from the LJR convention, is:
\begin{equation}\label{eq:sos}
  \Gamma^{(m)}_\theta(\Phi,\Phi)(\mathbf z)
  =\sum_{\iota}\alpha^{(m)}_\iota(\mathbf z)\,(\partial^{(m)}_\iota\Phi)^2,
  \quad
  \begin{cases}
   \alpha=\tfrac12,\ \partial_\iota=\partial_i & \text{open stratum},\\
   \alpha=\tfrac{c_{\theta,|S|}}{2},\ \partial_\iota=\sum_{j\in S}\partial_j
     & \text{collision stratum }\Delta_S.
  \end{cases}
\end{equation}
On a collision stratum $\Delta_S$ only the center-of-mass direction
$\sum_{j\in S}\partial_j$ survives; the difference directions
$\partial_i-\partial_j$ ($i,j\in S$) are dropped. This degeneracy, the blindness
of $\Gamma$ to intra-cluster differences on the diagonal, is the central analytic
fact of the paper.

The lifted single-time energy form is $\Ecyl(F,G)=\E[\GamN(F,G)]$ on $\Aone$.
The energy-first closed-pair construction (reproduced self-contained in
\cref{app:closedpair}, \cref{thm:foldA,thm:foldB}) produces a closed
operator-derivative pair $(\DN,\divN=\DN^{*})$ with integration by parts from a
$\Gamma$-tuple satisfying five axioms (\cref{ass:gamma-axioms}); the four
algebraic ones (symmetry, constancy-vanishing, matrix-positivity, weighted
matrix-positivity) are immediate for $\GamN$ on $\Aone$ from the square-of-sum
structure \eqref{eq:sos}. The fifth is fixed-sector form closability of $\Ecyl$
(\textnormal{(P1)}, on each fixed $\le2$-point sector; \cref{thm:sectorclos}). The
corresponding global form closability, the abstract target of the
projective criterion \eqref{eq:criterion} (established here only sectorwise), is the following.

\begin{hypothesis}[\textnormal{(FormClos)}, global target; $\beta$-1]
\label{hyp:beta1}
For every $(F_n)\subset\Aone$ with $F_n\to0$ in $\Hone$ and
$\Ecyl(F_n-F_m,F_n-F_m)\to0$, one has $\Ecyl(F_n,F_n)\to0$.
\end{hypothesis}

The route to this condition is a projective-limit criterion.
For a finite tuple $\mathbf x$ write $\mu_{\mathbf x}=$ law of
$(X^{x_1}_T,\dots,X^{x_m}_T)$ (the time-$T$ law, a probability on
$\R^m$), $\cE_{\mathbf x}(\Phi)=\int\Gamma^{(m)}_\theta(\Phi,\Phi)\,d\mu_{\mathbf x}$,
and for $\mathbf x\subseteq\mathbf y$ let $P^{\mathbf y}_{\mathbf x}=
\E[\,\cdot\mid\sigma(\text{$\mathbf x$-coords})]$. The criterion is:
\begin{equation}\label{eq:criterion}
  \text{(P1)}+\text{(P2)}+\text{(P3)}\;\Longrightarrow\;\text{\cref{hyp:beta1}},
\end{equation}
where \textnormal{(P1)} is finite-level closability of $\cE_{\mathbf x}$ on
$L^2(\mu_{\mathbf x})$, \textnormal{(P2)} is pullback compatibility
(automatic from LJR $n$-point consistency), and
\textnormal{(P3)} is the conditional-expectation energy contraction
\begin{equation}\label{eq:p3}
  \cE_{\mathbf x}\bigl(P^{\mathbf y}_{\mathbf x}\Phi,P^{\mathbf y}_{\mathbf x}\Phi\bigr)
  \;\le\;\cE_{\mathbf y}(\Phi,\Phi),\qquad \mathbf x\subseteq\mathbf y .
\end{equation}
(P1) at $m\le2$ is settled: it reduces to the standard weighted-Dirichlet
closability of $\nabla$ on $L^2(\muD)$~\citep{FOT}, so each fixed $\le2$-point
sector closes (\cref{thm:sectorclos}). The remaining pairwise structural question,
compatibility of the sectorwise pairs under the $2\to1$ deletion, is
\eqref{eq:p3} at $m=2$, supplied by $(\mathrm{CKC}_2^{\flat})$
(\cref{thm:deletion}); this is deletion-compatibility, not closability. Both
fixed-sector instances are reproduced and proved self-contained in
\cref{app:pairclos} (\cref{thm:sectorclos,thm:deletion}); no full-span closability enters.

\begin{remark}[The reference-measure subtlety]\label{rem:reference-measure}
$\cE_{\mathbf x}$ is integrated against the time-$T$ law
$\mu_{\mathbf x}$, not against the $\Lth{m}$-symmetrizing measure. This is
forced by $\Ecyl=\E_{\Pth}[\GamN]$ and is the source of all subtlety:
$\mu_{\mathbf x}$ charges the coincidence strata, exactly where
\eqref{eq:sos} is blind. We make this quantitative pairwise below.
\end{remark}

%==============================================================
\section{The pairwise reduction}\label{sec:reduction}
%==============================================================

With the lifted structure in place, the one open input is the deletion
inequality \eqref{eq:p3}, and this section reduces it to scalar data of the
two-point law. We work in the pairwise sub-core
$\Apair=\{F_{\mathbf x,T,\Phi}:|\mathbf x|\le2\}\subset\Aone$, stable under
coordinate deletion $2\to1\to0$. We reduce \eqref{eq:p3} for the single
inclusion $\{1\}\subset\{1,2\}$: retain one coordinate $\xi$, delete one
$\eta$.

By \cref{conv:LJR}(d) the conditional law of $\eta$ given $\xi$ is, with the
atom at the coincidence event $\eta=\xi$,
\begin{equation}\label{eq:kernel}
  \kappa(\xi,d\eta)=q(\xi)\,\delta_\xi(d\eta)+(1-q(\xi))\,p(\xi,\eta)\,d\eta,
  \qquad 0\le q\le1,\ \textstyle\int p(\xi,\cdot)=1,
\end{equation}
where $q(\xi)=\Prob[\eta=\xi\mid\xi]$ is the coalescence probability and
$p$ is the diffuse conditional density. In the reflecting-regular phase (R),
reflecting-regularity gives $q(\xi)<1$ for $\mu_x$-a.e.\ $\xi$ (the gap leaves
the diagonal, so the time-$T$ atom never captures full mass), hence
$1-q(\xi)>0$ and the prefactor $1/(1-q(\xi))$ is finite; both
$(\mathrm{CKC}_2)$ and $(\mathrm{CKC}_2^{\flat})$ below are understood on
$\{1-q>0\}$, which is $\mu_x$-co-null in phase (R). Write
$\langle\Phi\rangle_p=\int\Phi p\,d\eta$ and
$\langle\Phi\rangle_\xi=q\Phi(\xi,\xi)+(1-q)\langle\Phi\rangle_p$. The
energy decomposes, using \eqref{eq:sos} at $m=2$ (open stratum:
$\tfrac12[(\partial_\xi\Phi)^2+(\partial_\eta\Phi)^2]$; diagonal: keep
$\partial_\xi+\partial_\eta$, drop $\partial_\eta-\partial_\xi$), as
\begin{align}\label{eq:Ey}
  \cE_{\mathbf y}(\Phi)
  &=\tfrac12\!\int\!\mu_\xi(d\xi)(1-q)\!\int p\,[(\partial_\xi\Phi)^2+(\partial_\eta\Phi)^2]\,d\eta
   +\tfrac{\cth}{2}\!\int\!((\partial_\xi+\partial_\eta)\Phi)^2(\xi,\xi)\,\omega_\Delta(d\xi),\\
  \label{eq:Ex}
  \cE_{\mathbf x}(P\Phi)&=\tfrac12\!\int\!\mu_\xi(d\xi)\,(\partial_\xi(P\Phi))^2,
\end{align}
where $P\Phi(\xi)=\langle\Phi\rangle_\xi$, $\mu_\xi$ is the one-point
time-$T$ law (Brownian), and $\omega_\Delta$ is the diagonal weight (the
singular part of $\mu_{\mathbf y}$, $\omega_\Delta\propto q(\xi)\mu_\xi$,
tied to $\muD$ by \cref{conv:LJR}(c)). (P3) at $m=2$ is
$D(\Phi):=\cE_{\mathbf y}(\Phi)-\cE_{\mathbf x}(P\Phi)\ge0$.

Differentiating $P\Phi$ requires care, since the atom $\delta_\xi$ moves
with $\xi$ and the naive score identity is therefore wrong.

\begin{proposition}[Amended $(\star)$; derived in \cref{app:amended}]\label{prop:amended-star}
For $\kappa$ as in \eqref{eq:kernel} and $\Phi\in C^1_b$,
\begin{equation}\label{eq:amended-star}
  \partial_\xi(P\Phi)(\xi)
  =\underbrace{\langle\partial_\xi\Phi\rangle_\xi}_{\textnormal{(A) Jensen}}
   +\underbrace{q'(\xi)\,[\Phi(\xi,\xi)-\langle\Phi\rangle_p]}_{\textnormal{(B) weight-derivative}}
   +\underbrace{(1-q)\!\int\!\Phi\,\partial_\xi p\,d\eta}_{\textnormal{(C) diffuse score}}
   +\underbrace{q(\xi)\,(\partial_\eta\Phi)(\xi,\xi)}_{\textnormal{(D) MOVING-ATOM}}.
\end{equation}
Terms \textnormal{(B)+(C)} form the score covariance of the naive identity, while
term \textnormal{(D)}, a deleted-direction derivative at the atom, is
extra and absent from the naive identity. All four vanish
at $\theta=0$ ($q\equiv0$, $p$ independent of $\xi$).
\end{proposition}

Write $R:=\text{(B)}+\text{(C)}+\text{(D)}$ for the total obstruction
functional, so $\partial_\xi(P\Phi)=\langle\partial_\xi\Phi\rangle_\xi+R$.
From \eqref{eq:Ex} and \cref{prop:amended-star},
$\cE_{\mathbf x}(P\Phi)=\tfrac12\int\mu_\xi[\langle\partial_\xi\Phi\rangle_\xi^2
+2\langle\partial_\xi\Phi\rangle_\xi R+R^2]$. The term
$\langle\partial_\xi\Phi\rangle_\xi^2$ is dominated by the Jensen reference
inside $\cE_{\mathbf y}$; the sign-negative obstruction is
\begin{equation}\label{eq:obstruction}
  -\tfrac12\!\int\!\mu_\xi\,[\,2\langle\partial_\xi\Phi\rangle_\xi R+R^2\,].
\end{equation}
Expanding $R^2=((\mathrm B)+(\mathrm C)+(\mathrm D))^2$ exhibits a
moving-atom square $q^2(\partial_\eta\Phi(\xi,\xi))^2$, a weight-derivative
square $q'^2(\Phi(\xi,\xi)-\langle\Phi\rangle_p)^2$, a diffuse-score square,
and cross terms. The moving-atom and weight-derivative pieces live in the
difference direction at the diagonal, where $\Gamma$ is blind, so they
cannot be paid by the diagonal energy in \eqref{eq:Ey}; they must be
paid by the diffuse $\eta$-energy $G_\eta(\xi):=\langle(\partial_\eta\Phi)^2\rangle_p$.

Each squared obstruction piece is bounded by $G_\eta$ times an explicit
functional of the gap law (the trace inequalities of \cref{app:assembly}):
\begin{equation}\label{eq:traces}
\begin{aligned}
  |\Phi(\xi,\xi)-\langle\Phi\rangle_p|^2 &\le \tau_{\mathrm{wt}}(\xi)\,G_\eta(\xi),\\
  |\partial_\eta\Phi(\xi,\xi)|^2 &\le \tau_{\mathrm{at}}(\xi)\,\mathcal E_{\eta,\xi}(\Phi),\\
  |\mathrm{Cov}_p(\Phi,\partial_\xi\log p)|^2 &\le I_p(\xi)\,C^p_P(\xi)\,G_\eta(\xi),
\end{aligned}
\end{equation}
where the elements are
\begin{itemize}
\item the value trace
  $\tau_{\mathrm{wt}}(\xi)=\bigl(\int\sqrt{r(\xi,\eta)}\,p\,d\eta\bigr)^2$, with
  resistance $r(\xi,\eta)=\bigl|\int_\xi^\eta ds/p(\xi,s)\bigr|$;
\item the point-derivative trace $\tau_{\mathrm{at}}(\xi)$, finite only
  under \cref{hyp:dtrace} (not implied by reflecting regularity), bounded
  against the enhanced norm $\mathcal E_{\eta,\xi}$;
\item the diffuse Fisher information
  $I_p(\xi)=\langle(\partial_\xi\log p)^2\rangle_p$ and the Poincar\'e
  constant $C^p_P(\xi)$ of $p(\xi,\cdot)$.
\end{itemize}

The value-trace, Poincar\'e, and Fisher lines of \eqref{eq:traces} are
controlled by the first-order energy $G_\eta$ and are finite under reflecting
regularity (\cref{cor:dichotomy}(R)). The derivative line is not: point
evaluation of $\partial_\eta\Phi$ is not continuous for the first-order
Dirichlet energy $G_\eta$, so $\tau_{\mathrm{at}}$ finiteness is a genuine
second-order requirement, which we name.

\begin{hypothesis}[Moving-atom derivative trace]\label{hyp:dtrace}
On the pairwise core, equipped with an enhanced graph norm $\mathcal
E_{\eta,\xi}(\Phi)$ that includes second-order ($H^2$/Agmon) control of $\Phi$
in $\eta$ near the diagonal, there is $C_{\mathrm{tr}}(\xi)$, finite for
$\mu_\xi$-a.e.\ $\xi$, with
\[
  |\partial_\eta\Phi(\xi,\xi)|^2\;\le\;C_{\mathrm{tr}}(\xi)\,\mathcal E_{\eta,\xi}(\Phi);
\]
equivalently $\tau_{\mathrm{at}}(\xi)<\infty$.
\end{hypothesis}

The condition does not follow from reflecting regularity, which controls only
the value trace $\tau_{\mathrm{wt}}$ and the first-order quantities
$C^p_P,I_p$; it is a second-order, $H^2$-type requirement, a second-order
analogue of a negative-Sobolev dual obstruction, and plausibly corresponds to
an iterated-carr\'e ($\Gamma_2$) level of the carr\'e
hierarchy~\citep{BGL}; this correspondence is left as a question.

\begin{remark}[Boundary sign and $H^2$-characterization of \cref{hyp:dtrace}]\label{rem:dtrace-gamma2}
The $\Gamma_2$ level of the carr\'e hierarchy de-risks \cref{hyp:dtrace} without
discharging it. In gap/center-of-mass coordinates the Bakry
identity~\citep{BGL}
$\int\Gamma_2(f,f)\,dm=\int(Lf)^2\,dm$ for the $m$-symmetric sticky gap generator
($Lf=f''$ off $0$, Wentzell value $Lf(0)=\gamma^{-1}[f'(0{+})-f'(0{-})]$,
$m=dg+\gamma\,\delta_0$) gives
\[
  \int\Gamma_2(f,f)\,dm=\int_{g\neq0}(f'')^2\,dg+\gamma^{-1}[f'(0{+})-f'(0{-})]^2\ \ge\ 0 ,
\]
so the Wentzell-diagonal term is nonnegative for every finite stickiness and the
boundary does not deplete. With the diffuse density bounded below near the
diagonal in phase (R) (\cref{lem:checkB}), the one-dimensional Agmon trace
$|\partial_g\Phi(0)|^2\le C(\|\partial_g\Phi\|_{L^2(p)}^2+\|\partial_g^2\Phi\|_{L^2(p)}^2)$
controls the second-order term by the graph norm $\int(L\Phi)^2\,dm$, so
\cref{hyp:dtrace} holds on the $H^2$ (graph-norm) domain encoded by
$\mathcal E_{\eta,\xi}$, with $C_{\mathrm{tr}}\sim1/\inf_{\text{near }0}p$; it is
an $H^2$/Wentzell-membership condition, not a curvature one. On the diagonal-flat
sub-core $\partial_g\Phi(0)=0$ the trace is the center-of-mass derivative
controlled by the diagonal energy, recovering \cref{cor:neumann-subcore}
directly, while $C_{\mathrm{tr}}\sim1/\inf p$ degrades as $p\to0$ at the
coalescing endpoint, consistent with \cref{thm:wall-endpoint}.
\end{remark}

\begin{definition}[$\CKC$ at $m=2$]\label{def:ckc}
With the effective score strength
\begin{equation}\label{eq:S}
  S_A(\xi)\;=\;q'(\xi)^2\,\tau_{\mathrm{wt}}(\xi)
   +\underbrace{q(\xi)^2\,\tau_{\mathrm{at}}(\xi)}_{\textnormal{needs \cref{hyp:dtrace}}}
   +(1-q(\xi))^2 I_p(\xi)\,C^p_P(\xi)
   \;+\;\underbrace{A\,q(\xi)\,\tau_{\mathrm{at}}(\xi)}_{\substack{\textnormal{needs \cref{hyp:dtrace};}\\ \textnormal{universal }A>0}},
\end{equation}
for a universal constant $A>0$, and the cross-term Young factor
$1+1/\kappa_\xi(\xi)$ (\cref{app:assembly}, $\kappa_\xi$ the center-of-mass
margin from the diagonal channel), we denote the following sufficient condition
for $D(\Phi)\ge0$ for all $\Phi$ by $\CKC$:
\begin{equation}\label{eq:ckc}
  c(\xi)\;:=\;\frac{6\,(1+1/\kappa_\xi(\xi))}{1-q(\xi)}\,S_A(\xi)\;\le\;1
  \qquad\text{for }\mu_\xi\text{-a.e. }\xi
\end{equation}
(well-defined since $1-q>0$ a.e.\ in (R); see the safeguard after
\eqref{eq:kernel}).
\end{definition}

The numeric constants $6$ and $1+1/\kappa_\xi$, which are not optimized, come from
$(a+b+c)^2\le3(a^2+b^2+c^2)$ and one Young split; the content of $\CKC$ is the
functional form, that the obstruction strength $S$ be dominated by the diffuse
$\eta$-capacity in the trace constants. The patch summand $A\,q\,\tau_{\mathrm{at}}$
in \eqref{eq:S} carries the difference-direction half of the atom's
$\xi$-energy deficit, which the Jensen reference, a center-of-mass quantity in
\eqref{eq:Ey}, does not cover; the assembly (\cref{app:assembly}) fixes $A$ as
some finite universal constant, immaterial to the structural dichotomy of
\cref{sec:gap}. Like $q^2\tau_{\mathrm{at}}$, this term is finite if and only if
\cref{hyp:dtrace} holds, a second-order trace not implied by reflecting
regularity; both it and $A$ are (b)-level, and the structural verdict (a) runs on
the value pieces alone.

%==============================================================
\section{Evaluation: the two-point gap as a sticky Brownian motion}\label{sec:gap}
%==============================================================

The reduction has left explicit elements, all of them functions of the law
of the two-point gap. This section identifies that law and reads the
elements off the LJR boundary classification.

\begin{lemma}[Gap process]\label{lem:gap-sticky}
The gap $g=\eta-\xi$ is an autonomous scalar diffusion: off the diagonal it
is Brownian (diffusivity $2$, two independent unit Brownian motions), and at
$g=0$ it satisfies the Wentzell boundary condition of \cref{conv:LJR}(d),
i.e. it is a sticky Brownian motion with stickiness parameter
$\gamma=\gamma(\cth,\muD)$, monotone in $\cth$: $\gamma=0$ at the
independent endpoint, $\gamma\to$ the absorbing (exit) limit at the
coalescing endpoint.
\end{lemma}
\begin{proof}[Sketch]
Translation invariance of the flow decouples the gap from the center of
mass; the off-diagonal generator is $\partial_g^2$ and \cref{conv:LJR}(d)
gives the sticky boundary at $0$ (\cref{prop:sticky-bm}). The numeric map
$\gamma(\cth,\muD)$ is delegated by \cref{conv:LJR}; only the endpoints and
monotonicity are used.
\end{proof}

By \cref{prop:sticky-bm}(i)--(ii) the time-$T$ law of $g$ is
$\mathcal L(g_T)=q\,\delta_0+p(\cdot)\,dg$, with $q>0$ for $\gamma\in(0,\infty)$
and $p$ continuous; the qualitative facts used are recorded in
\cref{app:stickyBM}.

Two elements of $\CKC$ now have to be controlled. The first, the weight
derivative $q'$, turns out to be harmless.

\begin{lemma}\label{lem:checkA}
In the translation-invariant (homogeneous) reference, $q(\xi)\equiv q$ is
constant, so $q'\equiv0$ and the leading obstruction term
$q'^2\tau_{\mathrm{wt}}$ vanishes. In the point-started reference of
\cref{rem:reference-measure}, $q(\xi)$ is a smooth heat-kernel ratio with
$|q'|=O(T^{-1/2})$ bounded (\cref{prop:sticky-bm}), so $q'^2\tau_{\mathrm{wt}}$
is finite. Either way the $q'$-term does not drive failure.
\end{lemma}
\begin{proof}
The gap is autonomous (\cref{lem:gap-sticky}); the coincidence event
$\{g_T=0\}$ and the diagonal atom measure are translation-invariant, so in
the invariant reference their ratio $q(\xi)$ is constant. The point-started
bound is the gradient of the heat-kernel ratio.
\end{proof}

The point-versus-invariant choice is a modelling tension
(cf.~\cref{rem:reference-measure}), but \cref{lem:checkA} shows it is harmless at
$m=2$. The decisive element is instead the near-diagonal behaviour of the
diffuse density.

\begin{lemma}[Boundary dichotomy of $p$]\label{lem:checkB}
By \cref{prop:sticky-bm}(ii)--(iii), tied to \cref{def:trichotomy}:
\begin{itemize}
\item[\textnormal{(R)}] In the reflecting-regular phase, the diffuse density
is continuous and strictly positive at the boundary:
$p(\xi,\eta)\to P_0>0$ as $\eta\to\xi$.
\item[\textnormal{(C)}] At the coalescing/absorbing endpoint, $p$ vanishes:
$p(\xi,\eta)\sim C|\eta-\xi|$ as $\eta\to\xi$ (Dirichlet boundary).
\end{itemize}
\end{lemma}
\begin{proof}
A reflecting-regular boundary point is left by the process
(\cref{def:trichotomy}, clause $\Prob_0[H_r<\infty]>0$; the reflecting-regular
(intermediate) case of the two-point boundary classification of~\citep{LJR-FCN}),
so the diffuse transition density is continuous and positive at $0$; an
exit/absorbing boundary (the exit/absorbing (coalescing) case
of~\citep{LJR-FCN}) is a Dirichlet boundary, so the diffuse density vanishes
there. Both are standard boundary theory once the boundary type is fixed by the
LJR classification.
\end{proof}

\begin{corollary}[Finite/infinite dichotomy of the $\CKC$ elements]\label{cor:dichotomy}
With $r(\xi,\eta)=|\int_\xi^\eta ds/p(\xi,s)|$:
\begin{itemize}
\item[\textnormal{(R)}] $p\ge P_0>0$ near the diagonal $\Rightarrow$
$r\le|\eta-\xi|/P_0<\infty$ $\Rightarrow$ the value-trace, Poincar\'e and
Fisher constants $\tau_{\mathrm{wt}},\,C^p_P,\,I_p$ are finite
(positive codimension-one value-capacity, from reflecting regularity).
The derivative trace $\tau_{\mathrm{at}}$ is finite under \cref{hyp:dtrace} (a second-order requirement not supplied by
reflecting regularity).
\item[\textnormal{(C)}] $p\sim C|g|$ $\Rightarrow$ $r=\int ds/(Cs)=+\infty$
$\Rightarrow$ the value pieces blow up, $\tau_{\mathrm{wt}},\,C^p_P\to+\infty$
(value-capacity collapse), independently of any derivative-trace
consideration (the derivative trace $\tau_{\mathrm{at}}$ plays no role in phase
(C)).
\end{itemize}
\end{corollary}
\begin{proof}
Substitute \cref{lem:checkB} into \eqref{eq:traces}; the resistance integral
$r$ converges iff $p$ is bounded below, giving the value pieces. The derivative
trace $\tau_{\mathrm{at}}$ is governed by \cref{hyp:dtrace} in phase (R). The
coalescing failure runs on the value-capacity collapse alone and does not rely
on $\tau_{\mathrm{at}}$.
\end{proof}

This is the codimension-one positive-versus-zero value-capacity
dichotomy, now pinned to the boundary type; the derivative trace sits one order
higher (\cref{hyp:dtrace}).

%==============================================================
\section{Main results I: the sticky--coalescing dichotomy and the pairwise identity}\label{sec:main}
%==============================================================

The element-level dichotomy of \cref{cor:dichotomy} now assembles into the
first main result: absence of the obstruction in the reflecting phase
(\cref{thm:unobstruction}), the pairwise identity it carries under $(\mathrm{CKC}_2^{\flat})$
(\cref{thm:pairwise-ito}), failure at the coalescing endpoint
(\cref{thm:wall-endpoint}), and the discriminant statement
(\cref{thm:discriminant}). In the reflecting phase the elements of $\CKC$
are finite, and the apex obstruction has no witness.

For reference across the three main sections, the named conditions, each of
which appears in the title of every statement that uses it, are:
$(\mathrm{CKC}_2^{\flat})$ \eqref{eq:ckc-flat}, the first-order deletion
contraction on the diagonal-flat pair sector, which gives the deletion
inequality \textnormal{(P3)} \eqref{eq:p3} and the integration-by-parts
identity (\cref{thm:pairwise-ito}); $\COcap$ (\cref{def:cocap}) and the
mixed-mean bound $\COmix$ (\cref{def:comix}), uniformizations of the
two-point value capacity, either of which yields the two-point sector
Poincar\'e inequality and the Clark--Ocone--Green representation
(\cref{thm:clarkocone}), with $\COmix$ the operative form
(\cref{rem:two-well,rem:comix-instantiation}); and
$(\mathrm{Poin}^{\mathrm{dyn}})$ (\cref{def:dynamic-div}), the slicewise
Poincar\'e/Clark--Ocone input of the dynamic formula (\cref{thm:dynamic}),
automatic at $m=1$ and on the gap sector. \Cref{hyp:dtrace}, the moving-atom
trace, is a theorem on the second-order sub-core (\cref{prop:agmon}); the
three-point stratum regularity $(\mathrm R_3)$ (\cref{hyp:r3}) enters only
\cref{prop:m3-blocked}. The gap-sector dichotomy (\cref{thm:gap-co}), the
coalescence-indicator obstruction (\cref{prop:co-failure}), the occupation
law (\cref{prop:occupation}), and the dynamic formula on the gap and
one-point sectors (\cref{cor:dynamic-gap}) use only \cref{conv:LJR}.

\begin{theorem}[Absence of the apex value-capacity obstruction]\label{thm:unobstruction}
Assume \cref{conv:LJR} and the reflecting-regular phase \textnormal{(R)}. Then
the value-trace, Poincar\'e and Fisher constants
$\tau_{\mathrm{wt}},C^p_P,I_p$ are finite (positive codimension-one
value-capacity), so the apex value-capacity obstruction (the witness of
\cref{thm:wall-endpoint}) is absent in phase \textnormal{(R)}. The
moving-atom derivative trace $\tau_{\mathrm{at}}$ is finite under \cref{hyp:dtrace}.
\end{theorem}
\begin{proof}
\Cref{cor:dichotomy}(R) gives finiteness of the value pieces
$\tau_{\mathrm{wt}},C^p_P,I_p$; under \cref{hyp:dtrace} it also gives
$\tau_{\mathrm{at}}<\infty$. The only mechanism that forces $D(\Phi)<0$ for a
normalized difference-direction probe is a polar apex (zero
value-capacity), which the witness construction of \cref{thm:wall-endpoint}
requires; in codimension one polarity
needs $p\to0$ (\cref{cor:dichotomy}(C)), absent in phase (R) where
$p\ge P_0>0$. Hence the value witness cannot close.
\end{proof}

Removing the apex value-capacity witness is necessary for \eqref{eq:p3}, not
sufficient; the certified inequality $c(\xi)\le1$ is the separate statement
discussed next.

\begin{remark}[Structural vs.\ certified]\label{rem:struct-vs-cert}
\Cref{thm:unobstruction} asserts finiteness of the $\CKC$ constants, the
structural statement (a), not the certified inequality $c(\xi)\le1$, the
quantitative statement (b). The constants $6$ and $1+1/\kappa_\xi$ and the
delegated map $\gamma(\cth,\muD)$ leave the sharp stickiness threshold for (b)
unlocated, and the status of (b) is reference-dependent in a way worth
recording. In the translation-invariant reference, where $q'\equiv0$ and the
score reduces to $(1-q)^2I_pC^p_P$, no smallness is available: for every
diffuse density, testing the Poincar\'e inequality with $f(x)=x$ gives
$C^p_P\ge\mathrm{Var}_p$, while integrating the score by parts gives
$\mathrm{Var}_p\cdot I_p\ge1$ (Cram\'er--Rao), so $I_pC^p_P\ge1$ and the score
has an order-one floor; certification there must come from the margin
constants $A_\Delta,\kappa_\xi$, not from smallness of $S$. In the
point-started reference the score does vanish at the independent endpoint
($I_p\to0$ and $q'\to0$ as $\theta\to0$, the conditional density losing its
$\xi$-dependence), but the diagonal margin degenerates simultaneously
($c_\theta\to0$), and we do not resolve which effect prevails; exhibiting one
explicit $\theta$ with $c^{\flat}\le1$ is the most concrete open instantiation
(\cref{prob:cocap}). The prefactor $1/(1-q)$ in \eqref{eq:ckc}
is likewise a (b)-level quantity: the structural dichotomy of \cref{cor:dichotomy}
uses no non-collapse of the diffuse budget $1-q$, and at finite $T$ the atom mass
stays bounded away from $1$ uniformly in $\gamma$, so the coalescing failure
(\cref{thm:wall-endpoint}) runs solely through the value-capacity channel
$\tau_{\mathrm{wt}},C^p_P\to\infty$. The robust, $\theta$-uniform statement on
phase (R) is (a). The next proposition settles half of the certification
question: in the translation-invariant reference the margin constants cannot
rescue the inequality, because the contraction itself fails there.
\end{remark}

\begin{proposition}[The invariant reference fails outright]
\label{prop:invariant-fails}
In the translation-invariant reference with symmetric gap law ($q'\equiv0$,
$\langle g\rangle_p=0$), the deletion contraction \textnormal{(P3)} fails on
the diagonal-flat sector for every $\gamma\in(0,\infty)$. Explicitly, for
the flat observables $\Phi_R(\xi,\eta)=\chi_R(\xi+\eta)$, with $\chi_R$ the
identity on $[-R,R]$ smoothly truncated outside $[-2R,2R]$,
\[
  \cE_{\mathbf y}(\Phi_R)\;=\;(1+q)\,R+O(1)
  \;<\;2R+O(1)\;=\;\cE_{\mathbf x}(P\Phi_R)
  \qquad(R\to\infty),
\]
where $q\in(0,1)$ is the time-$T$ atom mass. Consequently
$(\mathrm{CKC}_2^{\flat})$ is unsatisfiable in the invariant reference.
\end{proposition}
\begin{proof}
On $\{\chi_R'=1\}$ the off-diagonal carr\'e of $\Phi_R$ is
$\tfrac12(1^2+1^2)=1$; at coincidence the Howitt--Warren bracket
$d\langle X^i,X^j\rangle_t=\mathbf 1_{\{X^i_t=X^j_t\}}\,dt$~\citep{HowittWarren}
makes the stuck pair a single Brownian point, so the diagonal carr\'e is
$\tfrac12\bigl((\partial_\xi+\partial_\eta)\Phi_R\bigr)^2=2$; integrating
against the time-$T$ law, whose diffuse and atomic masses are $1-q$ and $q$
per unit $\mu_\xi$-length, gives
$\cE_{\mathbf y}(\Phi_R)=\bigl[(1-q)+2q\bigr]R+O(1)$, the bulk being the
$\xi$-region $|2\xi|\le R$ of measure $R$ and the $O(1)$ coming from the
truncation collars and the Gaussian gap tails. For the deletion, symmetry
of the gap law gives $\langle\eta\rangle_p=\xi$, so $P\Phi_R(\xi)=\chi_R(2\xi)$
up to collar corrections, and
$\cE_{\mathbf x}(P\Phi_R)=\tfrac12\cdot4\cdot R+O(1)$. Since $q<1$ at
finite $T$ (\cref{prop:sticky-bm}), the deficit $(1-q)\cdot2R$ dominates
the collars for $R$ large, and \textnormal{(P3)} fails. The center-of-mass
channel is in deficit: the conditional expectation doubles the
center-of-mass derivative ($P\Phi$ rides both coordinates) while the
energy pays the stuck pair only once; no smallness of the score and no
margin constant can repair a failed inequality, so
$(\mathrm{CKC}_2^{\flat})$, which implies \textnormal{(P3)}, is
unsatisfiable.
\end{proof}

The certification question of \cref{rem:struct-vs-cert} therefore runs
through the point-started reference at small $\theta$ (\cref{prob:cocap}).

\begin{remark}[Numerical instantiation in the point-started reference]
\label{rem:instantiation}
For the point-started
pair with initial gap $g_0=1$ at horizon $T=1$, the joint time-$T$ law is
computable in closed form from Warren's time-change
construction~\citep{Warren97} and the Brownian local-time trivariate
density: with $\beta$ the diffusivity-$2$ Brownian gap and $\ell$ its
local time, the no-hit part of the gap law is the reflection difference,
the diffuse part is $\gamma^{-1}\psi_2(T-o;\,y,\,o/\gamma)$ in the
gap--occupation variables, with $\psi_2$ the trivariate density, the
atom--occupation density collapses to $\gamma\,\rho_{\tau_l}(T-\gamma l)$
with $\rho_{\tau_l}$ the inverse-local-time density, and the center of
mass, whose clock the Howitt--Warren bracket~\citep{HowittWarren}
accelerates while the pair is stuck, is conditionally Gaussian of
variance $(T+\mathrm{Occ}_T)/2$. Against this law the flat-sector
contraction was tested directly, as the generalized eigenvalue problem
$\sigma:=\max\,\cE_{\mathbf x}(P\Phi)/\cE_{\mathbf y}(\Phi)$ over exactly
flat Galerkin families (symmetric Hermite tensors plus cubic-antisymmetric
corrections, dimensions $36$ to $64$, grids to $260^2$), so that
\textnormal{(P3)} on the family is the statement $\sigma\le1$. The
results: $\sigma\in[0.94,\,0.97]$ across $\gamma\in\{0.05,\,0.2,\,0.5\}$
and across refinements, with the law normalization reproduced to between
$0.13$ and $1.3$ percent; and the center-of-mass direction, the witness of
\cref{prop:invariant-fails}, gives
$\cE_{\mathbf x}/\cE_{\mathbf y}\approx0.52$ here against the invariant
floor $4/(2+2q)>1.8$: the regression of the conditional mean toward the
starting configuration damps exactly the channel that defeats the
invariant reference. This is a Galerkin verification at stated tolerances,
not an interval-arithmetic proof, and it tests the contraction itself
rather than the sufficient form $c^{\flat}\le1$; within those
limits it answers the certification question of \cref{rem:struct-vs-cert}
for the point-started reference at explicit parameters, and the certified
upgrade, interval arithmetic over the same closed forms, is mechanical.
A second pass with exact dense eigensolvers and an explicit noise-floor
deflation (\cref{rem:comix-instantiation}, where the same law also
instantiates the mixed-mean bound) sharpens the margin: on families whose
variance spectrum stays above the quadrature noise floor the contraction
ratio improves to $\sigma\approx0.89$--$0.93$, while admitting sub-floor
directions inflates it toward $0.97$--$0.98$ through directions dominated by
quadrature error; in every configuration the tested value remains below~$1$.
Code and raw output are available at
\url{https://github.com/QuijoticResearch/sticky-flow-certificates}.
\end{remark}

On a fixed $\le2$-point sector, let $U_T=u(Z_{\mathbf x})$ be a real-valued smooth
sector observable, $u\in C_b^\infty(\R^{|\mathbf x|})$, so that $f(U_T)$,
$\DN U_T$, and $\GamN(F,U_T)$ are well-defined in that sector.

\begin{theorem}[Diagonal-flat pairwise black integration-by-parts identity; under $(\mathrm{CKC}_2^{\flat})$]\label{thm:pairwise-ito}
Assume \cref{conv:LJR} and the reflecting phase \textnormal{(R)}. Fix a one- or
two-point coordinate tuple $\mathbf x$ ($|\mathbf x|\le2$) and work in the
diagonal-flat sector
$\cD^{\flat}_{\mathbf x}=\{\Phi\in C_b^\infty(\R^{|\mathbf x|}):(\partial_\eta-\partial_\xi)\Phi(\xi,\xi)=0\}$
($|\mathbf x|=2$; the condition is vacuous for $|\mathbf x|=1$). Assume also the
reduced first-order contraction condition $(\mathrm{CKC}_2^{\flat})$
\eqref{eq:ckc-flat} on this fixed sector. This fixed sector
closes by \textnormal{(P1)} (\cref{thm:sectorclos}), so the fixed-sector
closed-pair construction (\cref{thm:foldA,thm:foldB}) produces a closed
densely-defined operator-derivative pair on the sector Hilbert space
$L^2(\mu_{\mathbf x})$,
\[
  \DN\colon\Dom(\DN)\subset L^2(\mu_{\mathbf x})\to\cU_{\GamN},
  \qquad \divN=\DN^{*}.
\]
On the diagonal-flat sector the moving-atom term is absorbed by the diagonal
(center-of-mass) energy (\cref{cor:neumann-subcore}), removing the
$\tau_{\mathrm{at}}$ terms, so $(\mathrm{CKC}_2^{\flat})$ is the relevant
first-order sufficient condition.
Under $(\mathrm{CKC}_2^{\flat})$, by \cref{thm:deletion} the
$2\to1$ deletion energy-contraction \textnormal{(P3)} holds, making the one- and
two-point sectorwise closed pairs compatible under coordinate deletion.
For the gradient one-form integrand $B_f=f'(U_T)\,\DN U_T$ ($f\in C_b^2$;
well-defined in the module by \cref{thm:foldA}$(\mathrm{iv})$), the
divergence functional
\begin{equation}\label{eq:black-ibp}
  \Lambda_f(F)\;:=\;\E\bigl[f'(U_T)\,\GamN(F,U_T)\bigr]
  \;=\;\langle\DN F,\,B_f\rangle_{\cU_{\GamN}},
  \qquad F\in\Dom(\DN),
\end{equation}
is well defined as it stands (since $B_f\in\cU_{\GamN}$); the pairing is
governed by the diagonal carr\'e $\GamT$ of~\eqref{eq:gamma-seed}, concentrated on
the coincidence stratum with weight $\cth\muD$. If moreover $\Lambda_f$ is
$L^2(\Omega)$-continuous, equivalently $B_f\in\Dom(\divN)$, then
$\divN(B_f)\in L^2(\Omega)$ exists and satisfies the weak identity
$\E[F\,\divN(B_f)]=\Lambda_f(F)$ for all $F\in\Dom(\DN)$. If, in addition,
$U_T\in\Dom(L_{\cN})$, where $L_{\cN}:=-\divN\DN$ is the sector Dirichlet
operator, and the corresponding chain-rule / domain-stability
assumptions hold for $f\in C_b^2$ (a Wentzell-compatible second-order condition,
of the \cref{hyp:dtrace} family), then the divergence has the explicit form
\[
  \divN(B_f)\;=\;-f'(U_T)\,L_{\cN}U_T-f''(U_T)\,\GamN(U_T,U_T).
\]
By blackness (\cref{thm:black}), this $\divN(B_f)$, when it exists, admits
no representation as a first-chaos / martingale stochastic integral.
\end{theorem}
\begin{proof}
Fix the sector $\mathbf x$ ($|\mathbf x|\le2$). (P1) at level $|\mathbf x|$
reduces to the settled weighted-Dirichlet closability of $\nabla$ on
$L^2(\muD)$~\citep{FOT}, so by \cref{thm:sectorclos} the sector closes and
\cref{hyp:beta1} holds on it; the relevant algebra is unital
($C_b^\infty(\R)$ resp.\ $C_b^\infty(\R^2)$) with $\Gamma$ the lifted seed
restricted to that sector (\cref{rem:sectorwise}), so the closed pair
$(\DN,\divN)$ on $L^2(\mu_{\mathbf x})$ is \cref{thm:foldA,thm:foldB}. Separately, on
the diagonal-flat sector the moving-atom term is absorbed by the diagonal
(center-of-mass) energy (\cref{cor:neumann-subcore}), removing the
$\tau_{\mathrm{at}}$ terms with no appeal to \cref{hyp:dtrace}, so the requirement
reduces to the first-order condition $(\mathrm{CKC}_2^{\flat})$ (assumed); under
it, by \cref{thm:deletion}, the $2\to1$ deletion energy-contraction (P3) holds. The pairing identity \eqref{eq:black-ibp} is the seed identity
$\langle\DN F,B_f\rangle=\E[f'(U_T)\GamN(F,U_T)]$ of \cref{thm:foldA}$(\mathrm{iv})$.
Under the stated continuity of $\Lambda_f$ (equivalently $B_f\in\Dom(\divN)$; a
sufficient condition is $U_T\in\Dom(L_{\cN})$ with the chain-rule/domain
stability for $f\in C_b^2$), $\divN(B_f)\in L^2$ exists by definition of
the adjoint (\cref{thm:foldB}$(\mathrm{iii})$), and the formula
$\divN(B_f)=-f'(U_T)L_{\cN}U_T-f''(U_T)\GamN(U_T,U_T)$ is the divergence of the
gradient one-form (carr\'e Leibniz on the smooth sector, $L_{\cN}=-\divN\DN$); the
blackness statement is~\cref{thm:black}.
\end{proof}

\begin{remark}[Operator-It\^o form: a single-time rearrangement, not a dynamic formula]\label{rem:ito-form}
Formally rearranging the centered divergence term gives the operator-It\^o
form
\begin{equation}\label{eq:black-ito}
  f(U_T)-f(U_0)\;=\;\divN\bigl(B_f\bigr)\;+\;\mathrm{Corr}_f(T),
\end{equation}
with correction $\mathrm{Corr}_f(T)$ governed by the diagonal carr\'e $\GamT$
(the $f''$-against-$d\,\GamN(U,U)$ term, weight $\cth\muD$). This is the
single-time interpretation of \eqref{eq:black-ibp}, not a separately proved
equation. Two points: writing $\divN(B_f)$ as a standalone $L^2$ object requires
$\Lambda_f$ to be $L^2$-continuous (equivalently $B_f\in\Dom(\divN)$), for which
$U_T\in\Dom(L_{\cN})$ (Wentzell-compatible $H^2$, of the \cref{hyp:dtrace}
family) is a sufficient condition; and \eqref{eq:black-ito} is a single-time
rearrangement, not a
dynamic It\^o formula, since $\divN(B_f)$ is centered ($\E[\divN(B_f)]=0$ since
$\DN 1=0$), so the non-centered increment $f(U_T)-f(U_0)$ cannot be produced by
the single-time pair; a dynamic formula requires the multi-time family
$\{U_t\}_{t\le T}$ and its bracket. A partitionwise dynamic formula of exactly
this kind is constructed in \cref{thm:dynamic}, the increment being produced
by an adapted telescope of conditional single-slice divergences; the
continuum disintegration remains open (\cref{prob:multitime}).
\end{remark}

\begin{remark}[Sectorwise closed-pair construction]\label{rem:sectorwise}
The closed-pair construction of \cref{thm:foldA,thm:foldB} requires a unital
algebra (the multiplier action $U\!\cdot\!D_0F$ and weighted matrix positivity
against multipliers $Q\in\cA$). The pairwise core $\Apair$ is not an
algebra: a product of two two-point functionals depends on up to four points,
leaving the pairwise sector. We therefore apply \cref{thm:foldA,thm:foldB}
per fixed sector, where the assumptions hold: for a fixed tuple
$\mathbf x$ ($|\mathbf x|\le2$) the algebra is $C_b^\infty(\R^{|\mathbf x|})$
(unital), $\Gamma$ the lifted seed restricted to that sector, yielding a
closed pair on $L^2(\mu_{\mathbf x})$. The global pairwise object is the
sectorwise family of these constructions, indexed by the fixed tuples
$\mathbf x$. It is not a direct sum and not a subspace of the flow
$L^2$ space: distinct fixed-pair sectors are non-orthogonal, and the uncountable
family of fixed tuples does not assemble into a single Hilbert space or a single
global operator. There is no single global operator on a ``$\Hone^{\le2}$'' here, and no $m\ge3$ closability enters: each member of the family is a fixed $\le2$-point sector, and
conditioning one fixed sector on another's coordinates would invoke joint laws of
up to four points (the $m\ge3,4$ structure the paper does not treat). The fixed
single sector is the unit of the construction.
\end{remark}

\begin{corollary}[Diagonal-flat sub-core: the reduced condition]\label{cor:neumann-subcore}
On the sub-core of pairwise observables with
$(\partial_\eta-\partial_\xi)\Phi(\xi,\xi)=0$ (flat in the $\Gamma$-blind
difference direction at coincidence), the moving-atom term
\eqref{eq:amended-star}\textnormal{(D)} does not vanish but is rotated into the
center-of-mass direction: $\partial_\eta\Phi(\xi,\xi)
=\tfrac12(\partial_\xi+\partial_\eta)\Phi(\xi,\xi)$, so
\textnormal{(D)}$=\tfrac{q}{2}(\partial_\xi+\partial_\eta)\Phi(\xi,\xi)$ is
controlled by the diagonal (center-of-mass) energy, available in phase
\textnormal{(R)}. Hence \cref{hyp:dtrace} is not needed: the moving-atom
contribution is absorbed by the diagonal energy (in phase (R), via the
comparison below), and the $\tau_{\mathrm{at}}$ terms drop from the score,
leaving the first-order reduced score
$S^{\flat}(\xi)=q'(\xi)^2\tau_{\mathrm{wt}}(\xi)+(1-q(\xi))^2 I_p(\xi)C^p_P(\xi)$
and the reduced sufficient condition
\begin{equation}\label{eq:ckc-flat}
  (\mathrm{CKC}_2^{\flat}):\quad
  c^{\flat}(\xi):=\frac{6\,A_\Delta(\theta)\,(1+1/\kappa_\xi(\xi))}{1-q(\xi)}\,S^{\flat}(\xi)\le1
  \quad\text{for }\mu_x\text{-a.e. }\xi
\end{equation}
(well-defined since $1-q>0$ a.e.\ in (R); see the safeguard after
\eqref{eq:kernel}), where $A_\Delta(\theta)<\infty$ is the finite
diagonal-absorption constant defined by the comparison
$\tfrac{q^2}{4}\mu_x\le A_\Delta(\theta)\,c_\theta\,\omega_\Delta$
(phase (R), $c_\theta>0$).
Its elements $\tau_{\mathrm{wt}},C^p_P,I_p$ are finite in phase (R)
(\cref{thm:unobstruction}), so no derivative-trace input is required,
and the pairwise black integration-by-parts identity
\eqref{eq:black-ibp} holds on this sub-core whenever
$(\mathrm{CKC}_2^{\flat})$ holds, or whenever (P3) is verified directly on the
sub-core.
\end{corollary}
\begin{proof}
Under the flatness condition,
$\partial_\eta\Phi(\xi,\xi)=\tfrac12(\partial_\xi+\partial_\eta)\Phi(\xi,\xi)$,
so term \eqref{eq:amended-star}\textnormal{(D)} lies in the center-of-mass
direction controlled by the diagonal energy
$\tfrac{c_\theta}{2}\int((\partial_\xi+\partial_\eta)\Phi)^2\,\omega_\Delta$ of
\eqref{eq:Ey}. Since $\omega_\Delta\propto q\,\mu_x$
(\cref{rem:reference-measure}; \cref{conv:LJR}(c)), the comparison
$\tfrac{q^2}{4}\mu_x\le A_\Delta(\theta)\,c_\theta\,\omega_\Delta$ holds for some
finite $A_\Delta(\theta)<\infty$ in phase \textnormal{(R)} ($c_\theta>0$,
$q\le1$), so the contribution of \textnormal{(D)} to the obstruction is absorbed
by $\cE_{\mathbf y}$'s diagonal term with no uncontrolled difference-direction
trace; the $\tau_{\mathrm{at}}$ terms therefore drop from the score, leaving the
first-order reduced score $S^{\flat}$ and the reduced condition
$(\mathrm{CKC}_2^{\flat})$, with the diagonal-absorption constant
$A_\Delta(\theta)$ entering its prefactor \eqref{eq:ckc-flat}; no derivative-trace
input is required; the remaining requirement is exactly the reduced first-order
contraction inequality $(\mathrm{CKC}_2^{\flat})$. The value-trace, Poincar\'e and
Fisher pieces are finite in (R) by \cref{cor:dichotomy}(R).
\end{proof}

We turn to the opposite phase, where the contraction fails. The coalescence wall
rests on a decomposition of the coalesced two-point law into gap and
center-of-mass parts.

\begin{lemma}[Gap/center-of-mass decomposition of the coalesced density]\label{lem:gap-com}
In the coordinates $g=\eta-\xi$ (gap), $c=\tfrac12(\xi+\eta)$ (center of mass), the
Brownian sticky two-point motion of \cref{conv:LJR} decomposes as follows.
\begin{enumerate}
\item[(i)] The gap $g$ is an autonomous scalar sticky Brownian motion of diffusivity
$2$, sticky at $0$ with stickiness fixed by \cref{conv:LJR}(b); in phase (R) its time-$T$
law has a positive atom at $0$ (\cref{lem:checkB})~\citep{EngelbertPeskir,ItoMcKean}.
\item[(ii)] The center of mass $c$ is a driftless martingale started at
$\bar x=\tfrac12(x_1+x_2)$, with driving noise orthogonal to the gap's throughout the
motion (including across the sticky diagonal: by construction of the Brownian sticky
flow the gap and center-of-mass noises are orthogonal). Its instantaneous variance rate
is $\tfrac12$ off the diagonal (independent particles) and $1$ on the diagonal
(coalesced, single Brownian particle). Hence, conditionally on the gap path $\mathcal G$,
\[
  c_T \mid \mathcal G \;\sim\; N\!\bigl(\bar x,\; V(\mathcal G)\bigr),\qquad
  V(\mathcal G)=\tfrac12 T+\tfrac12\,\mathrm{Occ},\quad
  \mathrm{Occ}=\bigl|\{s\le T: g_s=0\}\bigr|\in[0,T],
\]
so $V\in[\tfrac12 T,\,T]$.
\item[(iii)] Consequently the coalesced common-position density factors as a centered
Gaussian mixture,
\[
  \rho_{\mathrm{coal}}(\xi)=\int_{[\frac12 T,\,T]}\varphi_V(\xi-\bar x)\,\tilde\nu(dV),
\]
where $\varphi_V$ is the centered Gaussian of variance $V$ and $\tilde\nu$ is the
(sub-probability) law of $V$ on the coalescence event. In particular
$\rho_{\mathrm{coal}}$ is symmetric about $\bar x$, and since every component has
variance $V\le T$,
\[
  \rho_{\mathrm{coal}}(\xi)\le C\,e^{-(\xi-\bar x)^2/2T},\qquad
  |\rho_{\mathrm{coal}}'(\xi)|\le C(1+|\xi-\bar x|)\,e^{-(\xi-\bar x)^2/2T}.
\]
\end{enumerate}
\end{lemma}
\begin{proof}
(i)-(ii): the It\^o brackets $d\langle g\rangle=2\,dt$, $d\langle c\rangle=\tfrac12 dt$,
$d\langle g,c\rangle=0$ off the diagonal are immediate from \cref{conv:LJR} (off-diagonal
action $\tfrac12\Delta$, particles independent); the diagonal $\{X=Y\}$ is exactly
$\{g=0\}$, on which the Wentzell condition acts on $g$ alone, making $g$ autonomous, and
on which the coalesced pair moves as one Brownian particle so $c$ has variance rate $1$.
The orthogonality of the gap and center-of-mass noises holds throughout,
including across the sticky diagonal, and is a property of the two-point carr\'e,
by polarization of the square-of-sum structure \eqref{eq:sos}: off the diagonal,
\[
  \Gamma^{(2)}_\theta(g,c)
  =\tfrac12\bigl[\partial_\xi g\,\partial_\xi c+\partial_\eta g\,\partial_\eta c\bigr]
  =\tfrac12\bigl[(-1)\cdot\tfrac12+1\cdot\tfrac12\bigr]=0,
\]
while on the diagonal stratum only the cluster (center-of-mass) direction
survives and
$\Gamma^{(2)}_{\theta,\Delta}(g,c)\propto\bigl((\partial_\xi+\partial_\eta)g\bigr)
\bigl((\partial_\xi+\partial_\eta)c\bigr)=0\cdot1=0$, the gap being annihilated
by the center-of-mass derivative. Hence
$d\langle g,c\rangle=2\,\Gamma^{(2)}_\theta(g,c)\,dt=0$ in both strata, with no
appeal to the construction beyond \eqref{eq:sos}.

From the zero bracket to the conditional law: $g-g_0$ and $c-\bar x$ are
continuous $\Pth$-martingales, since both coordinate functions are
$\Lth2$-harmonic, off the diagonal $\tfrac12\Delta$ annihilates linear
functions, and on it the Wentzell value vanishes for $g$ (the one-sided
derivatives of the identity agree, so the jump is zero) while the coalesced
cluster moves as the driftless one-point motion, with
$d\langle c\rangle_t=\sigma(g_t)^2\,dt$, $\sigma^2=\tfrac12$ off the diagonal
and $1$ on it, so $\langle c\rangle$ is $\mathcal G$-measurable. By Knight's
theorem~\citep[Ch.~V]{RevuzYor} the orthogonal pair $(g-g_0,\,c-\bar x)$ has
independent Dambis--Dubins--Schwarz Brownian motions
$\beta^{g},\beta^{c}$ (after the usual extension beyond the horizon), with
$c_T=\bar x+\beta^{c}_{\langle c\rangle_T}$. Moreover
$\mathcal G=\sigma(g)\subseteq\sigma(\beta^{g})$: writing
$W_u:=g_0+\beta^{g}_{2u}$, a Brownian motion of diffusivity $2$, the sticky
path is recovered from $W$ by the time change of \cref{prop:occupation}
(Step~0), whose clock is a measurable functional of $W$. Hence $\beta^{c}$ is
independent of $\mathcal G$, and conditionally on $\mathcal G$ the COM is a
driftless Gaussian with deterministic variance
$\langle c\rangle_T=\int_0^T \sigma(g_s)^2\,ds=\tfrac12(T-\mathrm{Occ})+\mathrm{Occ}
=\tfrac12 T+\tfrac12\mathrm{Occ}$, giving $c_T\mid\mathcal G\sim N(\bar x,V)$ with
$V\in[\tfrac12T,T]$ since $\mathrm{Occ}\in[0,T]$.
(iii): on the coalescence event the common position equals $c_T$, so
$\rho_{\mathrm{coal}}(\xi)=\mathbb E[\mathbf 1_{g_T=0}\,\varphi_{V(\mathcal G)}(\xi-\bar x)]
=\int\varphi_V(\xi-\bar x)\,\tilde\nu(dV)$. Each component is centered at $\bar x$, giving
symmetry; and for $V\le T$, $\varphi_V(\xi-\bar x)\le(2\pi\cdot\tfrac12 T)^{-1/2}
e^{-(\xi-\bar x)^2/2T}$, so integrating $\tilde\nu$ gives the stated bound, and
differentiating gives the bound on $\rho_{\mathrm{coal}}'$ (the factor $(\xi-\bar x)/V$
is absorbed into $C(1+|\xi-\bar x|)$ using $V\ge\tfrac12 T$). The variance-$T$
bound is consistent with a Gaussian heat-kernel upper bound, stickiness only
delaying the spreading; a direct Aronson bound is obstructed by the
codimension-one Wentzell boundary, so this explicit decomposition is the rigorous
route.
\end{proof}

\begin{lemma}[Nonconstancy of the conditional atom weight]\label{lem:q-nonconstant}
Under the point-started reference $\mu_x=m_{\theta,T,x}$ from an initial pair
$(x_1,x_2)$ with $x_1\neq x_2$, the conditional coalescence weight $q(\xi)$ is
nonconstant in $\xi$ and lies in $\Dom(\cE_{\mathbf x})$ with
$\cE_{\mathbf x}(q)=\tfrac12\int q'(\xi)^2\,d\mu_x>0$. In the translation-invariant
reference, $q\equiv\mathrm{const}$ and $q'\equiv0$.
\end{lemma}
\begin{proof}
By \cref{lem:gap-com}(iii), $\rho_{\mathrm{coal}}$ is symmetric about $\bar x$, while the
retained one-point marginal $\rho_1=N(x_1,T)$ (Brownian one-point motion, \cref{conv:LJR})
is centered at $x_1$. For $x_1\neq x_2$ ($\bar x\neq x_1$) the ratio
$q=\rho_{\mathrm{coal}}/\rho_1$ is therefore nonconstant: a function symmetric about
$\bar x$ is not a constant multiple of a Gaussian centered at $x_1\neq\bar x$. For $T>0$
the densities are smooth and strictly positive, so $q$ is smooth and $q'\not\equiv0$ on a
set of positive $\mu_x$-measure.
For the energy: by \cref{lem:gap-com}(iii), $\rho_{\mathrm{coal}}(\xi)\le
C e^{-(\xi-\bar x)^2/2T}$ with the same quadratic rate $1/2T$ as $\rho_1$ (variance $T$),
so the quadratic exponents cancel in $q=\rho_{\mathrm{coal}}/\rho_1$ and $q(\xi),q'(\xi)$
grow at most exponentially, $|q'(\xi)|\le C(1+|\xi|)e^{a|\xi|}$ with $a=|x_2-x_1|/2T$.
Integrating against $d\mu_x=\rho_1\,d\xi$ (variance-$T$ Gaussian), the quadratic Gaussian
decay dominates the linear exponent:
$\cE_{\mathbf x}(q)=\tfrac12\int q'^2\,d\mu_x\le C\int (1+|\xi|)^2 e^{2a|\xi|}
e^{-(\xi-x_1)^2/2T}\,d\xi<\infty$. Hence $q\in\Dom(\cE_{\mathbf x})$, and
$\cE_{\mathbf x}(q)=\tfrac12\int q'^2\,d\mu_x>0$ by nonconstancy. In the
translation-invariant reference the COM-mixture and $\rho_1$ co-center and $q\equiv$const.
\end{proof}

\begin{theorem}[Failure of pairwise deletion-compatibility at the coalescing endpoint]\label{thm:wall-endpoint}
At the coalescing/absorbing endpoint (phase \textnormal{(C)}), under the
point-started reference $\mu_x=m_{\theta,T,x}$, the diffuse
density vanishes at the diagonal (\cref{lem:checkB}(C)), so the value-trace and
Poincar\'e constants $\tau_{\mathrm{wt}}$ and $C^p_P$ diverge, and the
difference-direction apex is polar for the weighted gap energy. This
failure is independent of the derivative-trace issue
($\tau_{\mathrm{at}}$/\cref{hyp:dtrace} play no role here). A
difference-direction probe
$\Phi_n(\xi,\eta)=g_n(\eta-\xi)$ with $g_n(0)=1$, $g_n\to0$ in $L^2(\bar p)$, and
weighted gap energy $\int|g_n'(g)|^2\,\bar p(g)\,dg\to0$, where $\bar p$ is the
diffuse density of the autonomous gap (such probes are supplied explicitly by
the logarithmic cutoffs of \cref{prop:co-failure}, Step~3), gives
$\cE_{\mathbf y}(\Phi_n)\to0$ while
$\cE_{\mathbf x}(P\Phi_n)\not\to0$, so $D(\Phi_n)<0$ for large $n$.
The deletion energy-contraction \textnormal{(P3)} \eqref{eq:p3} thus
fails, the inequality of \cref{thm:deletion} being violated, so the
sectorwise closed pairs are not deletion-compatible at the coalescing
endpoint (each fixed sector still closes by (P1)/\cref{thm:sectorclos}; it is the
$2\to1$ compatibility that breaks). The coalescing flow needs a strict sub-core
even at the pairwise level.
\end{theorem}
\begin{proof}
The argument needs both polarity (to send $\cE_{\mathbf y}(\Phi_n)\to0$)
and $q'\not\equiv0$ (to keep $\cE_{\mathbf x}(P\Phi_n)\not\to0$); neither alone
suffices.

Polarity, hence $\cE_{\mathbf y}(\Phi_n)\to0$. In phase \textnormal{(C)},
$p(g)\sim C|g|$ as $g\to0$ (\cref{lem:checkB}(C)), so the weighted resistance
$\int dg/p=\int dg/(C|g|)=\infty$ and the weighted point capacity of the apex
$\{g=0\}$ is $[\int dg/p]^{-1}=0$: the apex is polar for the weighted gap
energy (not for the ordinary $H^1$ energy, in which a point has positive
capacity; it is the weight $p\sim C|g|$, with $\int dg/p=\infty$, that
collapses the capacity). Explicit probes achieving the polarity and the
$L^2$-smallness simultaneously are the logarithmic cutoffs
$\psi_\varepsilon$ of \cref{prop:co-failure} (Step~3): capacity alone supplies
the pinning $g_n(0)=1$ and the energy smallness, while the simultaneous
$L^2(\bar p)$-smallness is an extra property, verified there by dominated
convergence; we take $g_n:=\psi_{\varepsilon_n}$. Set
$\Phi_n(\xi,\eta)=g_n(\eta-\xi)$. Being a pure difference-direction probe,
$\Phi_n$ is $\Gamma$-blind on the diagonal (its center-of-mass derivative
vanishes), so the diagonal (center-of-mass) term in \eqref{eq:Ey} vanishes, and
the open-stratum contribution to $\cE_{\mathbf y}(\Phi_n)$ is the weighted gap
energy ($(\partial_\xi\Phi_n)^2+(\partial_\eta\Phi_n)^2=2|g_n'|^2$ against the
weight $\tfrac12$)
\[
  \int(1-q)\,\bigl\langle|g_n'|^2\bigr\rangle_p\,d\mu_\xi
  \;=\;\int|g_n'|^2\,\bar p\,dg ,
\]
by the pushforward identity
$\int\mu_\xi(d\xi)\,(1-q(\xi))\,p(\xi,\xi+g)=\bar p(g)$ of
\cref{prop:co-failure} (Step~2), and this tends to $0$ by the probe
construction; thus $\cE_{\mathbf y}(\Phi_n)\to0$.

$q'\not\equiv0$, hence $\cE_{\mathbf x}(P\Phi_n)\not\to0$. By
\eqref{eq:kernel}, $P\Phi_n(\xi)=q(\xi)+(1-q(\xi))\langle g_n\rangle_p$, and
$P\Phi_n\to q$ in $L^2(\mu_\xi)$: by Jensen on the probability density
$p(\xi,\cdot)$, $(1-q)^2\le(1-q)$, and the pushforward identity of
\cref{prop:co-failure} (Step~2),
\[
  \int(1-q)^2\,\langle g_n\rangle_p^2\,d\mu_\xi
  \;\le\;\int(1-q)\,\langle g_n^2\rangle_p\,d\mu_\xi
  \;=\;\|g_n\|^2_{L^2(\bar p)}\;\longrightarrow\;0 .
\]
The
one-point form $\cE_{\mathbf x}(u)=\tfrac12\int(u')^2\,d\mu_\xi$ is closable (it is
the $m=1$ form, settled; (P1) at $m=1$), hence lower semicontinuous under
$L^2(\mu_\xi)$ convergence, and $q\in\Dom(\cE_{\mathbf x})$ with
$\cE_{\mathbf x}(q)=\tfrac12\|q'\|_{L^2(\mu_\xi)}^2$. Therefore
$\liminf_n\cE_{\mathbf x}(P\Phi_n)\ge\cE_{\mathbf x}(q)=\tfrac12\|q'\|^2>0$ by
\cref{lem:q-nonconstant}. (Lower semicontinuity replaces any convergence of
$\partial_\xi(P\Phi_n)$ to $q'$, which matters in phase (C), where
$I_p=\infty$ would defeat a direct derivative estimate.) Combined with
$\cE_{\mathbf y}(\Phi_n)\to0$, $\limsup_n D(\Phi_n)\le-\tfrac12\|q'\|^2<0$, so
$D(\Phi_n)<0$ for large $n$. Thus the deletion energy-contraction \eqref{eq:p3}
at $m=2$ fails (the inequality of \cref{thm:deletion} is violated), so the
sectorwise pairs are not deletion-compatible at the endpoint.
\end{proof}

The two phases can now be set side by side. Blackness is common to both, while
the boundary type alone toggles the obstruction.

\begin{theorem}[Boundary type discriminates the apex value-capacity obstruction]\label{thm:discriminant}
Both the reflecting and the coalescing Brownian sticky flows are black
(\cref{thm:black} and~\citep{LJR-circle-noise}, which also proves that the noise of
Arratia's coalescing flow is black). Yet the identified apex value-capacity
obstruction to the pairwise energy contraction \eqref{eq:p3} is discriminated by
the boundary type: in phase \textnormal{(R)} it is absent
(\cref{thm:unobstruction}), whereas in phase \textnormal{(C)}, for the
point-started time-$T$ law $\mu_\xi=m_{\theta,T,x}$, it is present and
the pairwise deletion energy-contraction (P3) fails
(\cref{thm:wall-endpoint}). Hence blackness is
shared but does not by itself decide It\^o-closability; the
reflecting/absorbing boundary type of \cref{def:trichotomy} is what toggles the
identified obstruction present/absent. This is not a reflecting-$\Leftrightarrow$-closable
equivalence: in phase \textnormal{(R)} the absence of the obstruction makes the
deletion energy-contraction (P3) hold under the stated first-order condition
$(\mathrm{CKC}_2^{\flat})$ on the diagonal-flat sub-core
(\cref{cor:neumann-subcore,thm:pairwise-ito}), so the sectorwise closed pairs
(each closed by (P1)/\cref{thm:sectorclos}) are deletion-compatible there; beyond
that sub-core, compatibility is open (\cref{prob:fullreflecting}).
\end{theorem}
\begin{proof}
Immediate from \cref{thm:black,thm:unobstruction,thm:wall-endpoint} and the
boundary classification \cref{def:trichotomy}.
\end{proof}

The discriminant lives on the point-started law. The energy datum and the
projective-limit criterion of \cref{sec:recap} are built on the point-started
consistent family $\mu_\xi=m_{\theta,T,x}$, with
$(\pi)_*\mu_{\mathbf y}=\mu_{\mathbf x}$, which is not a modelling choice but the
structure the closability apparatus runs on, so the discriminant of
\cref{thm:discriminant} lives intrinsically on that family. The invariant or
translation-started reference, under which $q'\equiv0$ (\cref{lem:checkA}) and the
phase-\textnormal{(C)} witness of \cref{thm:wall-endpoint} becomes borderline
($D(\Phi_n)\to0$ rather than $<0$), corresponds to a different operator,
the motion-semigroup object, which differs from the time-$T$-law operator by a
reference-measure conflation (cf.~\cref{rem:reference-measure}) rather than an
alternative reference within the present construction. The point-started reading
is thus the construction's own, and the discriminant is a property of it.

One objection deserves a direct answer in two parts: that the coalescing
failure is an artifact of conditioning on a non-stationary reference, which
the intrinsic noise never sees. First, why the point-started reading is
forced. The lifted form is $\Ecyl(F)=\E_{\Pth}[\GamN(F,F)]$ by definition
\eqref{eq:lifted-gamma}; on a fixed sector its reference measure is the
time-$T$ law $\mu_{\mathbf x}$ of the flow from its initial points, not one
choice among symmetrizing measures, but the only measure for which $\Ecyl$ is
the $L^2(\Pth)$-energy of the lifted observable, and the time-$T$ law
unavoidably charges the coincidence strata
(\cref{rem:reference-measure}). Substituting the translation-invariant measure
replaces the operator by the motion-semigroup object: it changes the question,
not the answer. Second, what fails where. The \textnormal{(P3)} witness of
\cref{thm:wall-endpoint} does use the point-started structure, through
$q'\not\equiv0$ (\cref{lem:q-nonconstant}), and degenerates to borderline in
the invariant reading. The representability failure of \cref{prop:co-failure}
does not: its mechanism uses only the linear vanishing of the diffuse
gap density at the absorbing boundary and a gap atom of mass in $(0,1)$, both
of which hold from any pair of distinct starts, and neither $q'$ nor any
conditional disintegration along $\xi$ enters. The coalescence indicator is a
nonconstant energy-null element of $L^2(\mu_{\mathbf y})$ for every
point-started pair, and the same probes defeat any reference whose gap
marginal vanishes linearly at the apex while charging the atom nontrivially.
At the endpoint, representability fails in every reference the construction
can carry; only the sharper contraction witness is reference-sensitive.

The two analyses live on opposite kinds of measure. The blackness
proof~\citep{LJR-circle-noise} controls the joint-minus-independent difference of
difference functionals integrated against an absolutely continuous time marginal,
which charges the diagonal with only $O(\sqrt\varepsilon)$ mass and sees the
diagonal carr\'e as $o(1)$ at the quadratic-variation scale; the energy
contraction \eqref{eq:p3} integrates the same diagonal carr\'e against the
time-$T$ law, whose singular part charges the diagonal
(\cref{rem:reference-measure}). Blackness misses the diagonal where closability
charges it, which is why the blackness machinery cannot supply the closability
input.

%==============================================================
\section{Main results II: the pairwise black Clark--Ocone representation}\label{sec:clarkocone}
%==============================================================

The identity \eqref{eq:black-ibp} is adjointness, not representation: it
determines $\divN=\DN^{*}$ against the energy pairing, but it does not produce,
for a given centered functional, an integrand whose divergence is that
functional. The classical upgrade is the Clark--Ocone
formula~\citep{Clark,Ocone,Nualart}
$F-\E F=\delta(\Pi\,DF)$ with $\Pi$ the predictable projection onto adapted
integrands, and both of its supports are gone here: blackness (\cref{thm:black})
removes the first chaos, hence the martingale-representation target onto which a
predictable projection could project, and the fixed sector is a single-time
object carrying no internal filtration in which ``predictable'' could be
pronounced. What can survive is an energy-module representation
\begin{equation}\label{eq:co-target}
  F-\E_{\mu_{\mathbf x}}F\;=\;\divN\bigl(\PiN\,\DN F\bigr),
  \qquad \E_{\mu_{\mathbf x}}F:=\textstyle\int F\,d\mu_{\mathbf x},
\end{equation}
with $\PiN$ built from the module $\cU_{\GamN}$ itself ($\E_{\mu_{\mathbf x}}$ is
the sector mean, the lift of $\E_{\Pth}$ to the fixed sector). Writing
$\Ltwoz(\mu_{\mathbf x}):=\{1\}^{\perp}$ for the centered subspace, the validity
of \eqref{eq:co-target} for all $F$ forces
$\mathrm{Ran}\,\divN\supseteq\Ltwoz(\mu_{\mathbf x})$, and this surjectivity is in
turn equivalent to a Poincar\'e inequality (spectral gap) for the sector energy
(\cref{cor:co-gap}). The content of this section is therefore a Poincar\'e
inequality. We prove it from the paper's own trace bounds plus one
uniformization, named $\COcap$; the deletion contraction
\textnormal{(P3)}/$(\mathrm{CKC}_2^{\flat})$ pays the two-point step, its second
appearance in the paper: the same inequality that makes the sectorwise pairs
deletion-compatible (\cref{thm:deletion}) also transports the one-point spectral
gap up the deletion tower $2\to1\to0$.

\begin{definition}[$\COcap$: uniform pairwise value capacity]\label{def:cocap}
In phase \textnormal{(R)}, with the conditional atom weight $q$, the value trace
$\tau_{\mathrm{wt}}$, and the conditional Poincar\'e constant $C^p_P$ of
\cref{sec:reduction}, set
\[
  K_{\mathrm{cap}}\;:=\;\operatorname*{ess\,sup}_{\xi\ (\mu_\xi)}
  \bigl[\,q(\xi)\,\tau_{\mathrm{wt}}(\xi)+C^p_P(\xi)\,\bigr].
\]
The condition $\COcap$ is: $K_{\mathrm{cap}}<\infty$.
\end{definition}

\begin{remark}[Status of $\COcap$]\label{rem:cocap-status}
Pointwise finiteness of both elements for $\mu_\xi$-a.e.\ $\xi$ is exactly the
phase-(R) statement of \cref{cor:dichotomy,thm:unobstruction}; $\COcap$ is its
uniform (essential-supremum) strengthening, of the same value-level, first-order
nature: no derivative trace ($\tau_{\mathrm{at}}$, \cref{hyp:dtrace}) enters. In
the translation-invariant reference the gap law does not depend on $\xi$
(\cref{lem:gap-sticky,lem:checkA}), the bracket is a constant, and $\COcap$ holds
trivially; under the point-started reference, the construction's own
(\cref{rem:reference-measure}), it is a genuine uniformity requirement
(\cref{prob:cocap}). Whether the reflecting phase alone yields the
representation below is open: \cref{thm:unobstruction} supplies pointwise
finiteness only, $\COcap$ asks for more, and \cref{thm:clarkocone} uses it.
Indeed we expect the essential-supremum form to fail under the point-started
reference (\cref{rem:two-well}); the operative point-started weakening is the
mixed-mean bound $\COmix$ of \cref{def:comix}, and on the autonomous gap sector
\cref{thm:gap-co} needs neither.
\end{remark}

\begin{lemma}[Sectorwise Poincar\'e inequality]\label{lem:sector-poincare}
Fix a sector $\mathbf x$ with $|\mathbf x|\le2$ under \cref{conv:LJR}, and let
$(\DN,\divN)$ be the fixed-sector closed pair of \cref{thm:pairwise-ito} (for
$|\mathbf x|=2$, the diagonal-flat pair), with closed form
$\cE(F)=\|\DN F\|^2_{\cU_{\GamN}}$ on $\Dom(\DN)\subset L^2(\mu_{\mathbf x})$.
\begin{enumerate}
\item[$(\mathrm i)$] ($m=1$; any phase.) $\mu_{\mathbf x}=N(x_1,T)$
and
\[
  \mathrm{Var}_{\mu_{\mathbf x}}(F)\;\le\;2T\,\cE(F),
  \qquad F\in\Dom(\DN).
\]
\item[$(\mathrm{ii})$] ($m=2$.) Assume phase \textnormal{(R)},
$(\mathrm{CKC}_2^{\flat})$ on the fixed sector, and $\COcap$. Then
\[
  \mathrm{Var}_{\mu_{\mathbf y}}(\Phi)\;\le\;C_N\,\cE(\Phi),
  \qquad C_N:=2K_{\mathrm{cap}}+2T,
  \qquad \Phi\in\Dom(\DN).
\]
\end{enumerate}
\end{lemma}
\begin{proof}
$(\mathrm i)$ The one-point motion is Brownian (\cref{conv:LJR}), so
$\mu_{\mathbf x}=N(x_1,T)$ and the classical Gaussian Poincar\'e
inequality~\citep{Bogachev} gives, for $u$ in the smooth core,
$\mathrm{Var}_{\mu_{\mathbf x}}(u)\le T\int(u')^2\,d\mu_{\mathbf x}
=2T\,\cE_{\mathbf x}(u)$; both sides are continuous along graph-norm convergent
sequences (the variance under $L^2$ convergence, the energy by closedness), so
the bound passes to $\Dom(\DN)$.

$(\mathrm{ii})$ We first prove the bound on the diagonal-flat smooth core, then
close. Disintegrate $\mu_{\mathbf y}(d\xi,d\eta)=\mu_\xi(d\xi)\,\kappa(\xi,d\eta)$
as in \eqref{eq:kernel} and use the total-variance decomposition
\[
  \mathrm{Var}_{\mu_{\mathbf y}}(\Phi)
  =\int\mu_\xi(d\xi)\,\mathrm{Var}_{\kappa(\xi,\cdot)}\bigl(\Phi(\xi,\cdot)\bigr)
   +\mathrm{Var}_{\mu_\xi}(P\Phi),
\]
with $P\Phi(\xi)=\langle\Phi\rangle_\xi$ the deletion conditional expectation.
Conditional term. For the atom-plus-density mixture
$\kappa=q\,\delta_\xi+(1-q)\,p\,d\eta$ the elementary mixture-variance identity is
\[
  \mathrm{Var}_{\kappa}\bigl(\Phi(\xi,\cdot)\bigr)
  =q(1-q)\bigl[\Phi(\xi,\xi)-\langle\Phi\rangle_p\bigr]^2
   +(1-q)\,\mathrm{Var}_{p}\bigl(\Phi(\xi,\cdot)\bigr).
\]
The value-trace inequality \eqref{eq:traces} bounds the first bracket by
$\tau_{\mathrm{wt}}(\xi)\,G_\eta(\xi)$ and the Poincar\'e constant of
$p(\xi,\cdot)$ bounds the second by $C^p_P(\xi)\,G_\eta(\xi)$, where
$G_\eta=\langle(\partial_\eta\Phi)^2\rangle_p$; hence
\[
  \mathrm{Var}_{\kappa}\bigl(\Phi(\xi,\cdot)\bigr)
  \le(1-q)\bigl[q\,\tau_{\mathrm{wt}}+C^p_P\bigr]\,G_\eta
  \le K_{\mathrm{cap}}\,(1-q)\,G_\eta
\]
by $\COcap$. Integrating against $\mu_\xi$ and comparing with the diffuse
$\eta$-budget of \eqref{eq:Ey},
$\cE_{\mathbf y}(\Phi)\ge\tfrac12\int\mu_\xi\,(1-q)\,G_\eta$, gives
\[
  \int\mu_\xi\,\mathrm{Var}_{\kappa}(\Phi(\xi,\cdot))
  \;\le\;2K_{\mathrm{cap}}\,\cE_{\mathbf y}(\Phi).
\]
Retained term. By part $(\mathrm i)$,
$\mathrm{Var}_{\mu_\xi}(P\Phi)\le2T\,\cE_{\mathbf x}(P\Phi)$ ($P\Phi$ is $C^1$
with finite one-point energy, \cref{prop:amended-star}); and on the diagonal-flat
sector, under $(\mathrm{CKC}_2^{\flat})$, the deletion contraction
\textnormal{(P3)} holds (\cref{thm:deletion,cor:neumann-subcore}), so
$\cE_{\mathbf x}(P\Phi)\le\cE_{\mathbf y}(\Phi)$. Summing the two terms gives the
bound with $C_N=2K_{\mathrm{cap}}+2T$ on the core; the passage to $\Dom(\DN)$ is
as in $(\mathrm i)$.
\end{proof}

The route through $\COcap$ disintegrates along the retained coordinate $\xi$.
There is a second route, along the gap path, whose first half holds outright
and which compresses the entire remaining difficulty into a single mixed-mean
bound.

\begin{proposition}[Center-of-mass reduction]\label{prop:com-half}
Assume \cref{conv:LJR} and phase \textnormal{(R)} on the fixed two-point sector.
Let $\mathcal G$ be the $\sigma$-field of the gap path,
$V=\tfrac12T+\tfrac12\mathrm{Occ}\in[\tfrac12T,T]$ and $\bar x$ as in
\cref{lem:gap-com}, and let $\beta_\theta>0$ be the diagonal-weight comparison
constant of \cref{rem:reference-measure} ($\omega_\Delta=\beta_\theta\,q\,\mu_\xi$).
Set $A_c(\theta):=\max\bigl(4,\,2/(\beta_\theta c_\theta)\bigr)<\infty$. Then for
every $\Phi$ in the diagonal-flat core:
\begin{enumerate}
\item[$(\mathrm i)$] $\E\bigl[\mathrm{Var}\bigl(\Phi(\xi,\eta)\mid\mathcal G\bigr)\bigr]
\;\le\;T\,A_c(\theta)\,\cE_{\mathbf y}(\Phi)$;
\item[$(\mathrm{ii})$] $\E[\Phi\mid\mathcal G]=H_\Phi(g_T,V)$, where
\[
  H_\Phi(g,v):=\int_\R\Phi\bigl(\bar x+\sqrt v\,z-\tfrac g2,\;
  \bar x+\sqrt v\,z+\tfrac g2\bigr)\,\varphi(z)\,dz,
  \qquad\varphi\ \text{the standard Gaussian};
\]
\item[$(\mathrm{iii})$] the $v$-sensitivity of the conditional mean is
first-order:
\[
  |\partial_vH_\Phi(g,v)|\;\le\;\tfrac12\,v^{-1/2}
  \Bigl(\E_{c\sim N(\bar x,v)}\bigl[\bigl((\partial_\xi+\partial_\eta)\Phi\bigr)^2
  (c-\tfrac g2,c+\tfrac g2)\bigr]\Bigr)^{1/2}
  \;\le\;(2T)^{-1/2}\,(\cdots)^{1/2}.
\]
\end{enumerate}
Consequently
\begin{equation}\label{eq:com-reduction}
  \mathrm{Var}_{\mu_{\mathbf y}}(\Phi)\;\le\;T\,A_c(\theta)\,\cE_{\mathbf y}(\Phi)
  \;+\;\mathrm{Var}\bigl(H_\Phi(g_T,V)\bigr).
\end{equation}
\end{proposition}
\begin{proof}
$(\mathrm i)$ By \cref{lem:gap-com}$(\mathrm{ii})$ (whose parts (i)--(ii) are
phase-agnostic), conditionally on $\mathcal G$ the center of mass is
$c\sim N(\bar x,V)$ with $V\le T$, and
$\tfrac{d}{dc}\Phi(c-\tfrac g2,c+\tfrac g2)=(\partial_\xi+\partial_\eta)\Phi$,
so the Gaussian Poincar\'e inequality gives
$\mathrm{Var}(\Phi\mid\mathcal G)\le T\,
\E[((\partial_\xi+\partial_\eta)\Phi)^2\mid\mathcal G]$. Taking expectations,
the left side of $(\mathrm i)$ is bounded by
$T\int((\partial_\xi+\partial_\eta)\Phi)^2\,d\mu_{\mathbf y}$. Split
$\mu_{\mathbf y}$ into its diffuse and diagonal parts: on the diffuse part,
$(a+b)^2\le2a^2+2b^2$ bounds the integral by $4$ times the open-stratum term of
\eqref{eq:Ey}; on the diagonal,
$\int((\partial_\xi+\partial_\eta)\Phi)^2(\xi,\xi)\,q\,d\mu_\xi
=\beta_\theta^{-1}\int((\partial_\xi+\partial_\eta)\Phi)^2\,d\omega_\Delta
=\tfrac{2}{\beta_\theta c_\theta}\,\cE^{\mathrm{diag}}_{\mathbf y}(\Phi)$.
Hence the total is $\le A_c(\theta)\,\cE_{\mathbf y}(\Phi)$.
$(\mathrm{ii})$ The conditional law of $c$ given $\mathcal G$ is $N(\bar x,V)$
and depends on $\mathcal G$ only through $V$, while the gap enters through
$g_T$.
$(\mathrm{iii})$ Differentiating under the integral,
$\partial_v\Phi(\bar x+\sqrt vz\mp\tfrac g2)=\tfrac{z}{2\sqrt v}\,
(\partial_\xi+\partial_\eta)\Phi$, and
$|\int z\,\psi(z)\varphi(z)\,dz|\le(\int z^2\varphi)^{1/2}(\int\psi^2\varphi)^{1/2}$
with $\int z^2\varphi=1$; finally $\tfrac12v^{-1/2}\le(2T)^{-1/2}$ for
$v\ge T/2$. The display \eqref{eq:com-reduction} is the total-variance
decomposition over $\mathcal G$ together with $(\mathrm i)$--$(\mathrm{ii})$.
\end{proof}

\begin{definition}[$\COmix$: mixed-mean bound]\label{def:comix}
$\COmix$ is the condition: there is $C_{\mathrm{mix}}<\infty$ with
\[
  \mathrm{Var}\bigl(H_\Phi(g_T,V)\bigr)\;\le\;C_{\mathrm{mix}}\,\cE_{\mathbf y}(\Phi)
  \qquad\text{for all }\Phi\text{ in the diagonal-flat core.}
\]
\end{definition}

\begin{corollary}[Sectorwise Poincar\'e under $\COmix$]\label{cor:comix-poincare}
In phase \textnormal{(R)}, under $\COmix$,
\[
  \mathrm{Var}_{\mu_{\mathbf y}}(\Phi)\;\le\;
  \bigl(T\,A_c(\theta)+C_{\mathrm{mix}}\bigr)\,\cE(\Phi),
  \qquad\Phi\in\Dom(\DN),
\]
on the diagonal-flat pair; neither $(\mathrm{CKC}_2^{\flat})$ nor $\COcap$ is
used. ($H_\Phi$ and both sides are continuous along graph-norm convergent
sequences, $\Phi\mapsto\E[\Phi\mid\mathcal G]$ being an $L^2$-contraction, so the
bound passes from the core to $\Dom(\DN)$.)
\end{corollary}

\begin{remark}[Numerical instantiation of $\COmix$, and the sector gap]\label{rem:comix-instantiation}
The closed-form joint law of \cref{rem:instantiation} makes $\COmix$ directly
testable: the conditioning data $(g_T,V)$ is exactly the tabulated
gap--occupation pair, and on Galerkin families the conditional mean reduces to
one-dimensional Gaussian-pair integrals; on the cubic-antisymmetric part of
the flat family the factor $(\eta-\xi)^3=g_T^3$ is $\mathcal G$-measurable and
factors out of $H_\Phi$. Against this law, at $g_0=T=1$,
$\gamma\in\{0.05,0.2,0.5\}$, over the exactly flat families of
\cref{rem:instantiation} deflated to the variance resolution the quadrature
supports (directions whose variance falls below the
quadrature noise floor report arbitrarily soft spurious modes, so the
deflation threshold is part of the protocol, and a fully resolved
one-dimensional computation of the dominant channel fixes the reference
value):
$(\mathrm a)$ the mixed-mean constant on the families is
$C_{\mathrm{mix}}\approx1.4$--$4.9$, against the exact gap-channel value
$2.00$--$2.05$ computed independently by resolved finite elements on the
closed-form gap--occupation marginal; the channel maximizer is the linear gap
mode $h(g)=g$, whose energy is the non-atomic gap mass and whose variance is
$\mathrm{Var}(g_T)\approx2T$, a ratio the autonomous gap theory
(\cref{thm:gap-co}) already controls, and at the stickier parameters
$\gamma\in\{0.2,0.5\}$ the planar $(g_T,V)$-content of $\COmix$ beyond this
pure gap channel is a correction of order fifteen percent;
$(\mathrm b)$ the conditional-variance half of \cref{prop:com-half} is
verified with a factor-two margin, the tested ratio
$\E[\mathrm{Var}(\Phi\mid\mathcal G)]/\cE_{\mathbf y}$ staying below $1.99$
against the proved $T\,A_c(\theta)=4$, and the contraction
$\mathrm{Var}(H_\Phi)\le\mathrm{Var}_{\mu_{\mathbf y}}(\Phi)$ holds strictly
on every family (largest tested eigenvalue $0.9997$);
$(\mathrm c)$ the resulting sector spectral gap is
$\lambda_N\approx0.19$--$0.53$ across the sweep, commensurate with the exact
one-point value $1/(2T)=0.5$: the point-started two-point flat sector does
not show the degeneration that \cref{prop:co-failure} exhibits at the
coalescing endpoint. In summary, at the same Galerkin (non-certified) grade
as \cref{rem:instantiation}: $\COmix$ holds at explicit parameters with
$C_{\mathrm{mix}}\approx2$, so the conditions of \cref{thm:clarkocone} are
simultaneously instantiated there in the point-started reference. Code and raw output are available at
\url{https://github.com/QuijoticResearch/sticky-flow-certificates}.
\end{remark}

\begin{remark}[Two-well tension: the essential-supremum form of $\COcap$ is likely too strong]\label{rem:two-well}
We record why we expect $\COcap$ to fail under the point-started
reference. The occupation-time element is exact
(\cref{prop:occupation}); the conditional-density steps remain heuristic. For
$w:=|\xi-x_1|$ large, the conditional gap law given $\xi$ is bimodal: a
never-met bulk near $g\approx x_2-\xi$, and a met-and-stuck
component near $g=0$ whose conditional mass is governed by the away time
$S:=T-\mathrm{Occ}$ of the gap. On $\{V=T-S/2\}$ the Gaussian cost ratio is
$\varphi_V(w)/\varphi_T(w)\approx e^{-w^2S/4T^2}$, and the exact tail
\eqref{eq:occupation-law} gives, by a Laplace saddle, with
$\kappa=|g_0|+2\gamma^{-1}T$, $g_0=x_2-x_1$,
\[
  -\log\E\bigl[e^{-w^2S/4T^2}\bigr]\;\sim\;
  \inf_{s>0}\Bigl(\frac{w^2s}{4T^2}+\frac{\kappa^2}{4s}\Bigr)
  \;=\;\frac{\kappa\,w}{2T}\qquad(w\to\infty),
\]
so the near-diagonal component carries conditional mass
$e^{-\kappa w/2T+o(w)}$: exponentially small of linear rate. The valley
between the two components, at intermediate gaps, has conditional density
exponentially small of quadratic rate $e^{-cw^2}$ (there both channels
are displaced Gaussians). A two-well test function then forces
\[
  C^p_P(\xi)\;\gtrsim\;e^{\,cw^2-\kappa w/2T}\;\longrightarrow\;\infty ,
\]
and the essential supremum in $\COcap$ diverges. What remains heuristic is the
valley upper bound and the transfer of \cref{prop:occupation} to the
conditional law given $\xi$; the occupation tail itself is now a closed-form
identity, and we note it is Gaussian-small, not polynomially small, in
$1/\epsilon$, the near-diagonal lump is far lighter than a polynomial
heuristic would suggest, but the divergence of $C^p_P(\xi)$ is only
strengthened, the valley being quadratically small against the lump's linear
rate. The essential-supremum formulation is thus the wrong uniformity, an
artifact of disintegrating along $\xi$. The gap-path route of
\cref{prop:com-half} avoids the two-well: the conditional well
structure never appears in the $(g_T,V)$ variables, and the operative remaining
condition is $\COmix$. In the translation-invariant reference $\COcap$ does
hold (\cref{rem:cocap-status}) and the issue is moot.
\end{remark}

\begin{corollary}[Spectral gap, kernel, closed range, Green operator]\label{cor:co-gap}
Assume a sectorwise Poincar\'e inequality
$\mathrm{Var}_{\mu_{\mathbf x}}\le C_N\,\cE$ holds on the fixed sector, as
supplied by \cref{lem:sector-poincare} or by \cref{cor:comix-poincare}
($C_N=2T$ at $m=1$), and set $\lambda_N:=C_N^{-1}>0$. Then on the fixed
sector:
\begin{enumerate}
\item[$(\mathrm i)$] $\ker\DN=\R\,1$: the constants are exactly the energy-null
functionals, and $\|F-\E_{\mu_{\mathbf x}}F\|^2\le C_N\,\|\DN F\|^2_{\cU_{\GamN}}$
for $F\in\Dom(\DN)$;
\item[$(\mathrm{ii})$] $\mathrm{Ran}\,\DN$ is closed in $\cU_{\GamN}$, and the
restriction $\DN\colon\Dom(\DN)\cap\Ltwoz\to\mathrm{Ran}\,\DN$ is a bijection
with bounded inverse $\iota$, $\|\iota\|\le C_N^{1/2}$;
\item[$(\mathrm{iii})$] the sector Dirichlet operator $L_{\cN}=-\divN\DN$ of
\cref{thm:pairwise-ito}, with domain
$\{F\in\Dom(\DN):\DN F\in\Dom(\divN)\}$, is self-adjoint with $-L_{\cN}\ge0$
(von Neumann's theorem for $\DN^{*}\DN$, \citep[Ch.~5]{Schmuedgen}); the centered
space $\Ltwoz$ reduces $L_{\cN}$, the spectrum of $-L_{\cN}$ on $\Ltwoz$ lies in
$[\lambda_N,\infty)$, and the Green operator
$G_N:=\bigl(-L_{\cN}|_{\Ltwoz}\bigr)^{-1}$ is bounded, $\|G_N\|\le C_N$, mapping
$\Ltwoz$ into $\Dom(L_{\cN})\cap\Ltwoz$;
\item[$(\mathrm{iv})$] $\ker\divN=(\mathrm{Ran}\,\DN)^{\perp}$ and
$\mathrm{Ran}\,\divN=\Ltwoz(\mu_{\mathbf x})$: the divergence maps onto the
centered space, and is injective on $\mathrm{Ran}\,\DN$ but not on its full
domain.
\end{enumerate}
\end{corollary}
\begin{proof}
$(\mathrm i)$ $\DN F=0$ forces $\mathrm{Var}(F)=0$ by
\cref{lem:sector-poincare}, so $F$ is constant; constants are energy-null by
\cref{thm:foldA}$(\mathrm{ii})$; the displayed bound is the lemma.
$(\mathrm{ii})$ If $F_n\in\Dom(\DN)\cap\Ltwoz$ and $\DN F_n\to v$, then by
$(\mathrm i)$ applied to differences $(F_n)$ is $L^2$-Cauchy, $F_n\to F\in\Ltwoz$,
and closedness of $\DN$ gives $F\in\Dom(\DN)$, $\DN F=v$; hence the range is
closed and $\|\iota v\|\le C_N^{1/2}\|v\|$.
$(\mathrm{iii})$ Self-adjointness of $\DN^{*}\DN$ for a closed densely defined
operator is von Neumann's theorem; $\langle-L_{\cN}F,F\rangle=\|\DN F\|^2\ge0$.
$L_{\cN}1=0$ (as $\DN1=0$), and $\langle L_{\cN}F,1\rangle=-\langle\DN F,\DN
1\rangle=0$ with $F-\E F\in\Dom(L_{\cN})$ whenever $F\in\Dom(L_{\cN})$, so
$\Ltwoz$ reduces $L_{\cN}$; on $\Ltwoz\cap\Dom(L_{\cN})$ the form bound
$\langle-L_{\cN}F,F\rangle=\|\DN F\|^2\ge\lambda_N\|F\|^2$ (part $(\mathrm i)$)
places the spectrum in $[\lambda_N,\infty)$, and the spectral theorem gives the
bounded inverse $G_N$.
$(\mathrm{iv})$ $\ker\divN=\ker\DN^{*}=(\mathrm{Ran}\,\DN)^{\perp}$ is standard;
$\langle\divN u,1\rangle=\langle u,\DN1\rangle_{\cU_{\GamN}}=0$ gives
$\mathrm{Ran}\,\divN\subseteq\Ltwoz$; conversely, for $F\in\Ltwoz$ put
$G:=G_NF\in\Dom(L_{\cN})\cap\Ltwoz$, so $\DN G\in\Dom(\divN)$ and
$\divN(\DN G)=-L_{\cN}G=F$. Injectivity on $\mathrm{Ran}\,\DN$ is
$(\mathrm{Ran}\,\DN)\cap\ker\divN=\{0\}$; the kernel
$(\mathrm{Ran}\,\DN)^{\perp}$ is in general nontrivial (the module contains
non-exact one-forms $U\cdot\DN F$).
\end{proof}

\begin{definition}[The minimal-norm Green representer $\PiN$]\label{def:PiN}
Under the Poincar\'e assumption of \cref{cor:co-gap}, let $P_{\nabla}$ be the
orthogonal projection of $\cU_{\GamN}$ onto the closed subspace
$\mathrm{Ran}\,\DN$ (the representable, or exact, one-forms), and define
\[
  \PiN\;:=\;\DN\,G_N\,\iota\,P_{\nabla}\;\colon\;\cU_{\GamN}\longrightarrow\cU_{\GamN},
\]
with $\iota,G_N$ as in \cref{cor:co-gap}. $\PiN$ vanishes on
$(\mathrm{Ran}\,\DN)^{\perp}=\ker\divN$, and on exact one-forms it inverts the
one-form (module) Laplacian: for $F\in\Dom(\DN)$,
$\bigl(\DN\divN\bigr)\,\PiN\,\DN F=\DN F$, so on $\mathrm{Ran}\,\DN$ the operator
$\PiN$ is the Green operator $(\DN\divN)^{-1}$.
\end{definition}

\begin{lemma}[Properties of $\PiN$]\label{lem:PiN-props}
$\PiN$ is everywhere defined and bounded with $\|\PiN\|\le C_N$, self-adjoint and
nonnegative on $\cU_{\GamN}$, and maps $\cU_{\GamN}$ into
$\mathrm{Ran}\,\DN\cap\Dom(\divN)$. It is not idempotent in general:
$\PiN^2=\PiN$ holds iff $\divN$ restricted to exact one-forms is a partial
isometry, which the energy setting does not provide (cf.\
\cref{rem:co-not-chaos}(2)).
\end{lemma}
\begin{proof}
Every element of $\mathrm{Ran}\,\DN$ is exact, $P_\nabla u=\DN F$ with
$F=\iota P_\nabla u\in\Dom(\DN)\cap\Ltwoz$ and $\|\DN F\|\le\|u\|$. Then
$\PiN u=\DN G_NF$, with
$\|\DN G_NF\|^2=\langle\divN\DN G_NF,G_NF\rangle=\langle F,G_NF\rangle
\le C_N\|F\|^2\le C_N^2\,\|\DN F\|^2$, giving $\|\PiN\|\le C_N$, and
$\DN G_NF\in\mathrm{Ran}\,\DN\cap\Dom(\divN)$ since
$G_NF\in\Dom(L_{\cN})$. Self-adjointness and nonnegativity: for $F,G\in
\Dom(\DN)\cap\Ltwoz$,
$\langle\PiN\DN F,\DN G\rangle=\langle\divN\DN G_NF,G\rangle=\langle F,G\rangle$,
symmetric in $F,G$ and $\ge0$ at $F=G$; both sides vanish against
$(\mathrm{Ran}\,\DN)^{\perp}$. For the displayed inversion property,
$\divN\PiN\DN F=F-\E F$ (\cref{thm:clarkocone}$(\mathrm i)$ below, whose proof
does not use this lemma) and $\DN(F-\E F)=\DN F$. Finally
$\PiN^2\,\DN F=\DN G_N^2F\neq\DN G_NF$ unless $G_N$ is the identity on the
relevant vectors, i.e.\ unless $\|\divN u\|=\|u\|$ on exact $u$.
\end{proof}

\begin{theorem}[Pairwise black Clark--Ocone representation]\label{thm:clarkocone}
Fix a sector $\mathbf x$ with $|\mathbf x|\le2$ under \cref{conv:LJR}. For
$|\mathbf x|=2$, assume phase \textnormal{(R)}, work with the diagonal-flat
closed pair of \cref{thm:pairwise-ito}, and assume either
$(\mathrm{CKC}_2^{\flat})$ on this fixed sector together with $\COcap$
(\cref{lem:sector-poincare}) or the mixed-mean bound $\COmix$
(\cref{cor:comix-poincare}); for $|\mathbf x|=1$ only \cref{conv:LJR}
is used, in any phase. Then, on $L^2(\mu_{\mathbf x})$:
\begin{enumerate}
\item[$(\mathrm i)$] (Representation.) For every $F\in\Dom(\DN)$,
$\PiN\DN F\in\Dom(\divN)$ and
\[
  F-\E_{\mu_{\mathbf x}}F\;=\;\divN\bigl(\PiN\,\DN F\bigr).
\]
\item[$(\mathrm{ii})$] (Uniqueness and minimality.) $\PiN\DN F$ is the
unique representer in $\mathrm{Ran}\,\DN=(\ker\divN)^{\perp}$; among all
$u\in\Dom(\divN)$ with $\divN u=F-\E F$ it is the one of minimal module norm.
\item[$(\mathrm{iii})$] (Variance identity and energy bound.)
\[
  \mathrm{Var}_{\mu_{\mathbf x}}(F)
  =\bigl\langle\PiN\DN F,\ \DN F\bigr\rangle_{\cU_{\GamN}},
  \qquad
  \|\PiN\DN F\|^2_{\cU_{\GamN}}\le C_N\,\mathrm{Var}_{\mu_{\mathbf x}}(F).
\]
\item[$(\mathrm{iv})$] (Surjectivity of the divergence.) Every centered
$F\in\Ltwoz(\mu_{\mathbf x})$, not necessarily in $\Dom(\DN)$, is a divergence:
$F=\divN(\DN G_NF)$; and $\DN G_N(F-\E F)=\PiN\DN F$ whenever $F\in\Dom(\DN)$,
so the two constructions agree on the common domain.
\end{enumerate}
\end{theorem}
\begin{proof}
$(\mathrm i)$ Write $F_0:=F-\E_{\mu_{\mathbf x}}F\in\Dom(\DN)\cap\Ltwoz$, so
$\DN F=\DN F_0\in\mathrm{Ran}\,\DN$, $P_\nabla\DN F=\DN F$, and
$\iota\,\DN F=F_0$; thus $\PiN\DN F=\DN G_NF_0$. Since
$G_NF_0\in\Dom(L_{\cN})$, the one-form $\DN G_NF_0$ lies in $\Dom(\divN)$, and
$\divN(\DN G_NF_0)=-L_{\cN}G_NF_0=F_0$.
$(\mathrm{ii})$ Two representers differ by an element of
$\ker\divN=(\mathrm{Ran}\,\DN)^{\perp}$ (\cref{cor:co-gap}$(\mathrm{iv})$);
$\PiN\DN F\in\mathrm{Ran}\,\DN$ is therefore the unique representer in that
subspace, and any other representer $u=\PiN\DN F+w$ with
$w\in(\mathrm{Ran}\,\DN)^{\perp}$ has
$\|u\|^2=\|\PiN\DN F\|^2+\|w\|^2$, so the norm is minimal.
$(\mathrm{iii})$ Pairing $(\mathrm i)$ against $F$:
$\langle\PiN\DN F,\DN F\rangle_{\cU_{\GamN}}
=\langle\divN\PiN\DN F,F\rangle=\langle F_0,F\rangle=\|F_0\|^2
=\mathrm{Var}_{\mu_{\mathbf x}}(F)$. For the bound,
$\|\PiN\DN F\|^2=\langle F_0,G_NF_0\rangle\le C_N\|F_0\|^2$ as in
\cref{lem:PiN-props}.
$(\mathrm{iv})$ Is \cref{cor:co-gap}$(\mathrm{iv})$; the coincidence on
$\Dom(\DN)$ is the computation in $(\mathrm i)$.
\end{proof}

\begin{remark}[Energy-module, not first-chaos; fixed-sector, not flow-space]\label{rem:co-not-chaos}
Three limitations delimit the scope of \cref{thm:clarkocone}.
\textnormal{(1)} It is not a martingale representation. The representer
$\PiN\DN F$ is a one-form in the sectorwise energy module $\cU_{\GamN}$, paired
through the diagonal carr\'e $\GamT$; it is not a stochastic integral against any
martingale of the flow, and by blackness (\cref{thm:black}) it could not be: the
first chaos a predictable projection would project onto is trivial. The
representation is the analytic shadow that survives blackness: a divergence of a
one-form, not an integral of a predictable process.
\textnormal{(2)} $\PiN$ is not a projection. In the Wiener
case~\citep{Clark,Ocone,Nualart} the
Clark--Ocone projection is idempotent because the It\^o isometry makes the
divergence a partial isometry on predictable integrands; here $\divN$ is not
partially isometric on exact forms, and the minimal-norm representer requires the
Green factor $G_N$, $\PiN=(\DN\divN)^{-1}P_\nabla$. We keep the symbol $\PiN$
only for the structural analogy $\delta\,\Pi\,D=I-\E$. Even at $|\mathbf x|=1$
the statement is the static Green representation for the Gaussian
measure $N(x_1,T)$, an Ornstein--Uhlenbeck-type identity, not the
filtration-indexed Brownian Clark--Ocone formula.
\textnormal{(3)} It is sectorwise. The theorem lives on a fixed one- or
two-point sector, on the diagonal-flat pair at $m=2$, with centering
$\E_{\mu_{\mathbf x}}$; the sectorwise family does not assemble into a single
operator on the flow $L^2$-space (\cref{rem:sectorwise}), and the dynamic
statement is the separate \cref{thm:dynamic} (cf.\ \cref{rem:ito-form}).
\end{remark}

\begin{proposition}[Failure at the coalescing endpoint: the coalescence indicator]\label{prop:co-failure}
In phase \textnormal{(C)} (the coalescing endpoint), on the fixed two-point
sector started from $x_1\neq x_2$, with the diagonal-flat closed pair (which
closes by \textnormal{(P1)}, \cref{thm:sectorclos}), the coalescence indicator
\[
  K:=\mathbf 1_{\Delta}\in L^2(\mu_{\mathbf y}),
  \qquad K=\mathbf 1_{\{X^{x_1}_T=X^{x_2}_T\}}\ \text{as a flow functional},
\]
satisfies $K\in\Dom(\DN)$ with $\DN K=0$, and $K$ is nonconstant in
$L^2(\mu_{\mathbf y})$. Consequently $\ker\DN\supsetneq\R\,1$, the sectorwise
Poincar\'e inequality fails, $\overline{\mathrm{Ran}\,\divN}\subseteq
\{1,K\}^{\perp}\subsetneq\Ltwoz(\mu_{\mathbf y})$, and the centered functional
$K-\E K$ admits no representation $\divN(u)$, for any
$u\in\Dom(\divN)$. The discriminant is therefore asymmetric: in phase
\textnormal{(R)} the representation holds under
$(\mathrm{CKC}_2^{\flat})$ and $\COcap$ (\cref{thm:clarkocone}), while at the
coalescing endpoint it fails outright, extending the
deletion-contraction failure of \cref{thm:wall-endpoint} from the inequality
\textnormal{(P3)} to representability itself.
\end{proposition}
\begin{proof}
Write $g_0:=x_2-x_1>0$ (relabelling the pair if necessary). All estimates
reduce to the autonomous gap, where everything is explicit.

Step 1: explicit gap law. At the coalescing endpoint the gap
$g_t$ is a Brownian motion of diffusivity $2$ absorbed at $0$
(\cref{lem:gap-sticky,lem:gap-com}; \cref{prop:sticky-bm}$(\mathrm{iii})$), and
coalescing one-point motions never cross: the two particles have continuous
paths, so they must meet before they cross, and at the endpoint meeting is
absorbing; hence the order is preserved and $g_T\ge0$. Its time-$T$ law is
$\bar q\,\delta_0+\bar p(g)\,dg$ with, by the method of images for the
absorbed kernel of $\partial_g^2$ and the reflection principle,
\[
  \bar p(g)=\frac{1}{\sqrt{4\pi T}}
  \Bigl[e^{-(g-g_0)^2/4T}-e^{-(g+g_0)^2/4T}\Bigr]\quad(g>0),
  \qquad
  \bar q=2\bigl(1-\Phi_{\mathrm{std}}(g_0/\sqrt{2T})\bigr),
\]
$\Phi_{\mathrm{std}}$ the standard normal distribution function. Hence
$0<\bar q<1$ for $0<g_0$ and $T<\infty$; $\bar p>0$ on $(0,\infty)$; and
$\bar p$ is $C^1$ with $\bar p(0)=0$ and
$\bar p'(0{+})=(4\pi T)^{-1/2}(g_0/T)\,e^{-g_0^2/4T}>0$, so there are
$\delta_0>0$ and $C<\infty$ with
\begin{equation}\label{eq:pbar-linear}
  \bar p(g)\;\le\;C\,g,\qquad 0<g\le\delta_0 .
\end{equation}
This is the linear vanishing of \cref{lem:checkB}(C), here in closed form.

Step 2: reduction of pure gap probes to the gap law. The diffuse part
of $\mu_{\mathbf y}$ is $\mu_\xi(d\xi)\,(1-q(\xi))\,p(\xi,\eta)\,d\eta$
\eqref{eq:kernel}, and its pushforward under $(\xi,\eta)\mapsto g=\eta-\xi$ is
the diffuse part of the gap law, i.e.
$\int\mu_\xi(d\xi)\,(1-q(\xi))\,p(\xi,\xi+g)=\bar p(g)$; likewise
$\int q\,d\mu_\xi=\bar q=\mu_{\mathbf y}(\Delta)$. For a probe
$\Phi=\psi(\eta-\xi)$ with $\psi\in C^\infty_b(\R)$, the center-of-mass
derivative $(\partial_\xi+\partial_\eta)\Phi$ vanishes identically, so the
diagonal term of \eqref{eq:Ey} is zero, and
$(\partial_\xi\Phi)^2+(\partial_\eta\Phi)^2=2\,\psi'(\eta-\xi)^2$ gives
\begin{equation}\label{eq:gap-energy}
  \cE_{\mathbf y}(\Phi)=\int_0^\infty\psi'(g)^2\,\bar p(g)\,dg,
  \qquad
  \|\Phi\|^2_{L^2(\mu_{\mathbf y})}
  =\bar q\,\psi(0)^2+\int_0^\infty\psi(g)^2\,\bar p(g)\,dg .
\end{equation}

Step 3: explicit logarithmic probes. Fix $\chi\in C^\infty(\R)$ with
$\chi\equiv0$ on $(-\infty,0]$, $\chi\equiv1$ on $[1,\infty)$, $0\le\chi\le1$,
$|\chi'|\le2$, and for $\varepsilon\in(0,\delta_0)$ set
\[
  \psi_\varepsilon(g):=\chi\!\Bigl(\frac{\ln(\delta_0/|g|)}{\ln(\delta_0/\varepsilon)}\Bigr)
  \quad(g\neq0),\qquad \psi_\varepsilon(0):=1 .
\]
Then $\psi_\varepsilon\in C^\infty_b(\R)$ is even, $\equiv1$ on
$[-\varepsilon,\varepsilon]$, and $\equiv0$ off $[-\delta_0,\delta_0]$
(smoothness across $|g|=\varepsilon,\delta_0$ because $\chi$ is flat at $1$ and
$0$); in particular $\Phi_\varepsilon:=\psi_\varepsilon(\eta-\xi)$ lies in the
diagonal-flat core ($\psi_\varepsilon'\equiv0$ near $0$), with
$\Phi_\varepsilon(\xi,\xi)=1$. On $[\varepsilon,\delta_0]$,
$|\psi_\varepsilon'(g)|\le2/\bigl(g\ln(\delta_0/\varepsilon)\bigr)$, and
$\psi_\varepsilon'=0$ elsewhere, so by \eqref{eq:gap-energy} and
\eqref{eq:pbar-linear}
\[
  \cE_{\mathbf y}(\Phi_\varepsilon)
  \;\le\;\frac{4}{\ln^2(\delta_0/\varepsilon)}
  \int_\varepsilon^{\delta_0}\frac{\bar p(g)}{g^{2}}\,dg
  \;\le\;\frac{4C}{\ln^2(\delta_0/\varepsilon)}
  \int_\varepsilon^{\delta_0}\frac{dg}{g}
  \;=\;\frac{4C}{\ln(\delta_0/\varepsilon)}
  \;\xrightarrow[\varepsilon\downarrow0]{}\;0 .
\]
Moreover $0\le\psi_\varepsilon\le1$ and $\psi_\varepsilon\to0$ pointwise on
$(0,\infty)$ as $\varepsilon\downarrow0$, so
$\int\psi_\varepsilon^2\,\bar p\,dg\to0$ by dominated convergence; with
\eqref{eq:gap-energy}, $\Phi_\varepsilon\to K$ in $L^2(\mu_{\mathbf y})$
(the atom values agree, $\Phi_\varepsilon(\xi,\xi)=1=K$ there) and
$\|\DN\Phi_\varepsilon\|^2_{\cU_{\GamN}}=\cE_{\mathbf y}(\Phi_\varepsilon)\to0$.

Step 4: the closure chain. We record each link separately; the only
structural input is that $\DN$ is the closure of the seed $D_0$ on the
diagonal-flat core (\textnormal{(P1)}, \cref{thm:sectorclos,thm:foldB}), valid
in every phase. (For Step~5 we also need $\divN=\DN^{*}$ to be well defined,
i.e.\ the flat core dense in $L^2(\mu_{\mathbf y})$: in gap/center-of-mass
coordinates, replace $\Phi(c,g)$ by $\Phi(c,\chi_\sigma(g))$ with
$\chi_\sigma$ smooth, odd, $\chi_\sigma'(0)=0$, and $\chi_\sigma(g)=g$ off
$[-\sigma,\sigma]$; this flattens any $\Phi\in C_b^\infty$ without changing
its diagonal values, and converges to $\Phi$ in $L^2(\mu_{\mathbf y})$ as
$\sigma\downarrow0$, the two functions differing only on the diffuse set
$\{0<|g|\le\sigma\}$, whose $\mu_{\mathbf y}$-mass vanishes. Since
$C_b^\infty(\R^2)$ is dense in $L^2(\mu_{\mathbf y})$, so is the flat core.)
\begin{enumerate}
\item[\textnormal{(a)}] ($L^2$ convergence.) $\Phi_\varepsilon\to K$ in
$L^2(\mu_{\mathbf y})$: the atom values agree exactly
($\Phi_\varepsilon(\xi,\xi)=\psi_\varepsilon(0)=1=K$ on $\Delta$), so by the
pushforward identity of Step~2 the difference is carried by the diffuse part
alone:
\[
  \|\Phi_\varepsilon-K\|^2_{L^2(\mu_{\mathbf y})}
  =\int_0^\infty\psi_\varepsilon(g)^2\,\bar p(g)\,dg
  \;\xrightarrow[\varepsilon\downarrow0]{}\;0 .
\]
\item[\textnormal{(b)}] (Energy-Cauchy.) $\Phi_\varepsilon-\Phi_{\varepsilon'}$
is again a pure gap probe, so \eqref{eq:gap-energy} applies to it directly and
\[
  \cE_{\mathbf y}(\Phi_\varepsilon-\Phi_{\varepsilon'})
  =\int_0^\infty\bigl(\psi_\varepsilon'-\psi_{\varepsilon'}'\bigr)^2\bar p\,dg
  \;\le\;2\,\cE_{\mathbf y}(\Phi_\varepsilon)+2\,\cE_{\mathbf y}(\Phi_{\varepsilon'})
  \;\xrightarrow[\varepsilon,\varepsilon'\downarrow0]{}\;0
\]
by Step~3. By the seed isometry
$\|D_0\Phi\|^2_{\cU_{\GamN}}=\cE_{\mathbf y}(\Phi)$
(\cref{thm:foldA}$(\mathrm{ii})$), the one-forms $(D_0\Phi_\varepsilon)$ are
Cauchy in $\cU_{\GamN}$, with limit $v=0$ since
$\|D_0\Phi_\varepsilon\|^2=\cE_{\mathbf y}(\Phi_\varepsilon)\to0$.
\item[\textnormal{(c)}] (Closedness.) The graph pairs
$(\Phi_\varepsilon,D_0\Phi_\varepsilon)\to(K,0)$ in
$L^2(\mu_{\mathbf y})\times\cU_{\GamN}$, and $\DN=\overline{D_0}$ is closed;
hence $K\in\Dom(\DN)$.
\item[\textnormal{(d)}] (Energy-null.) $\DN K=\lim_\varepsilon
D_0\Phi_\varepsilon=0$, equivalently $K$ lies in the domain of the closed form
with $\cE_{\mathbf y}(K)=\|\DN K\|^2=0$.
\item[\textnormal{(e)}] (Nonconstancy.) Under the point-started law
$\mu_{\mathbf y}=m_{\theta,T,\mathbf x}$ with $x_1\neq x_2$,
$\mathrm{Var}_{\mu_{\mathbf y}}(K)
=\mu_{\mathbf y}(\Delta)\bigl(1-\mu_{\mathbf y}(\Delta)\bigr)
=\bar q(1-\bar q)>0$ by the explicit reflection-principle value of Step~1, so
$K\notin\R\,1$ in $L^2(\mu_{\mathbf y})$.
\end{enumerate}
No deletion contraction, no $(\mathrm{CKC}_2^{\flat})$, and no capacity
uniformity enters: the chain uses only \cref{conv:LJR}, the
fixed-sector closability \textnormal{(P1)}, and the closed-form gap law of
Step~1. Note also that
\textnormal{(a)}--\textnormal{(d)} do not contradict closability: \textnormal{(P1)}
asserts that the closed form is well defined on graph limits, not that its
kernel is trivial; what fails at the endpoint is irreducibility.

Step 5: non-representability. For any
$u\in\Dom(\divN)$,
$\langle\divN u,K\rangle=\langle u,\DN K\rangle_{\cU_{\GamN}}=0$ and
$\langle\divN u,1\rangle=0$, so
$\mathrm{Ran}\,\divN\perp\mathrm{span}\{1,K\}$; if $\divN u=K-\E K$ then
$\|K-\E K\|^2=\langle\divN u,K\rangle-\E K\,\langle\divN u,1\rangle=0$,
contradicting $\bar q(1-\bar q)>0$.
\end{proof}

The dichotomy of \cref{thm:discriminant} therefore extends from the deletion
contraction to representability itself, asymmetrically: in phase (R), under
the named conditions of \cref{thm:clarkocone}, every centered sector
functional is a divergence with a canonical minimal-norm representer, while at
the coalescing endpoint the energy vanishes on the very
functional that records coalescence, and no divergence can reach it.

\subsection{The gap sector: the dichotomy in full}\label{subsec:gapco}

On one sub-sector the full dichotomy follows from \cref{conv:LJR} alone. The
diagonal-flat gap observables
\[
  \cA^{\flat}_{\mathrm{gap}}:=\bigl\{\,f(\eta-\xi)\;:\;f\in C_b^\infty(\R),\
  f'(0)=0\,\bigr\}
\]
form a genuine unital algebra ($(fg)'(0)=f'(0)g(0)+f(0)g'(0)=0$, and
$1\in\cA^{\flat}_{\mathrm{gap}}$), unlike the pairwise core
(\cref{rem:sectorwise}). Let $\bar\nu:=\bar q\,\delta_0+\bar p(g)\,dg$ be the
time-$T$ law of the autonomous gap (\cref{lem:gap-sticky,prop:sticky-bm}), with
$\int\bar p=1-\bar q$, and $\hat p:=\bar p/(1-\bar q)$ the normalized diffuse
density. Since $\|f(\eta-\xi)\|^2_{L^2(\mu_{\mathbf y})}=\int f^2\,d\bar\nu$,
the closure of $\cA^{\flat}_{\mathrm{gap}}$ in $L^2(\mu_{\mathbf y})$ is
$L^2(\bar\nu)$; the restricted $\Gamma$-tuple satisfies the five axioms of
\cref{ass:gamma-axioms} (the closability axiom being the one-dimensional
weighted statement, as in \cref{thm:sectorclos}), so
\cref{thm:foldA,thm:foldB} produce a closed gap pair
$(D_{\bar\nu},\delta_{\bar\nu}:=D_{\bar\nu}^{*})$ on $L^2(\bar\nu)$ with, on the
core (the two-sided version of \eqref{eq:gap-energy}),
\[
  \cE_{\bar\nu}(f)\;=\;\|D_{\bar\nu}f\|^2\;=\;\int_\R f'(g)^2\,\bar p(g)\,dg .
\]
This is a sub-pair of the fixed two-point sector, with its own adjoint; the
statements below are internal to it.

\begin{theorem}[Black Clark--Ocone dichotomy on the gap sector]\label{thm:gap-co}
Fix the two-point sector from $x_1\neq x_2$ under \cref{conv:LJR}, and consider
the gap pair $(D_{\bar\nu},\delta_{\bar\nu})$. Only \cref{conv:LJR} is
assumed; $(\mathrm{CKC}_2^{\flat})$, $\COcap$, and $\COmix$ do not enter.
\begin{enumerate}
\item[\textnormal{(R)}] In the reflecting phase, with
\[
  \bar K:=\bar q\,\hat\tau_{\mathrm{wt}}+C^{\hat p}_P,\qquad
  \hat\tau_{\mathrm{wt}}:=\Bigl(\int\sqrt{\hat r(g)}\,\hat p(g)\,dg\Bigr)^{2},\qquad
  \hat r(g):=\Bigl|\int_0^g\frac{ds}{\hat p(s)}\Bigr|,
\]
one has $\bar K<\infty$ and the Poincar\'e inequality
$\mathrm{Var}_{\bar\nu}(f)\le\bar K\,\cE_{\bar\nu}(f)$ for all
$f\in\Dom(D_{\bar\nu})$. Consequently the conclusions of
\cref{cor:co-gap,def:PiN,lem:PiN-props,thm:clarkocone} hold verbatim for the
gap pair with $C_N=\bar K$: every centered $F\in L^2(\bar\nu)$ is a divergence,
$F=\delta_{\bar\nu}(D_{\bar\nu}G_{\bar\nu}F)$, and for $F\in\Dom(D_{\bar\nu})$,
\[
  F-\E_{\bar\nu}F\;=\;\delta_{\bar\nu}\bigl(\Pi_{\bar\nu}D_{\bar\nu}F\bigr),
\]
with $\Pi_{\bar\nu}$ the minimal-norm Green representer of the gap pair.
\item[\textnormal{(C)}] At the coalescing endpoint the probes of
\cref{prop:co-failure} lie in $\cA^{\flat}_{\mathrm{gap}}$, so
$K=\mathbf 1_{\{g=0\}}\in\Dom(D_{\bar\nu})$ with $D_{\bar\nu}K=0$, $K$
nonconstant, and the centered functional $K-\E_{\bar\nu}K$ admits no
representation $\delta_{\bar\nu}(u)$.
\end{enumerate}
Thus the gap sector carries the full dichotomy:
the representation holds in the reflecting phase with explicit constant
$\bar K$, and fails at the coalescing endpoint.
\end{theorem}
\begin{proof}
\textnormal{(R)} Poincar\'e. For $f$ in the core, the mixture-variance
identity for $\bar\nu=\bar q\,\delta_0+(1-\bar q)\,\hat p\,dg$ gives
\[
  \mathrm{Var}_{\bar\nu}(f)
  =\bar q(1-\bar q)\bigl[f(0)-\langle f\rangle_{\hat p}\bigr]^2
   +(1-\bar q)\,\mathrm{Var}_{\hat p}(f).
\]
The value-trace bound
$[f(0)-\langle f\rangle_{\hat p}]^2\le\hat\tau_{\mathrm{wt}}\int f'^2\,\hat p$
is the one-dimensional instance of \eqref{eq:traces} (\cref{app:assembly}: write
$f(0)-\langle f\rangle_{\hat p}=-\int(\int_0^g f'\,ds)\hat p\,dg$ and apply
Cauchy--Schwarz with the resistance $\hat r$), and
$\mathrm{Var}_{\hat p}(f)\le C^{\hat p}_P\int f'^2\,\hat p$ is the Poincar\'e
inequality of the diffuse gap density. Collecting, and using
$(1-\bar q)\int f'^2\,\hat p=\int f'^2\,\bar p=\cE_{\bar\nu}(f)$,
\[
  \mathrm{Var}_{\bar\nu}(f)\le
  \bigl[\bar q\,\hat\tau_{\mathrm{wt}}+C^{\hat p}_P\bigr]\,\cE_{\bar\nu}(f)
  =\bar K\,\cE_{\bar\nu}(f),
\]
which passes to $\Dom(D_{\bar\nu})$ by graph-norm continuity of both sides.

Finiteness of $\bar K$. $C^{\hat p}_P<\infty$ is
\cref{prop:sticky-bm}$(\mathrm i)$. For $\hat\tau_{\mathrm{wt}}$: in phase (R),
$\hat p$ is continuous and positive at the wall
(\cref{prop:sticky-bm}$(\mathrm{ii})$), so
$\hat P_0:=\inf_{[-\delta,\delta]}\hat p>0$ for some $\delta>0$ and
$\hat r(g)\le|g|/\hat P_0$ there; for $|g|\ge\delta$ the two-sided Gaussian
envelopes of the sticky kernel (\cref{prop:sticky-bm}$(\mathrm i)$; stickiness
only delays the spreading, cf.\ the proof of \cref{lem:gap-com}$(\mathrm{iii})$)
give $\hat r(g)\le C\,e^{(|g|+|g_0|)^2/4T}$ and
$\hat p(g)\le C\,e^{-(|g|-|g_0|)^2/4T}$, whence
$\sqrt{\hat r(g)}\,\hat p(g)\le C'\,e^{-g^2/8T+C''|g|}$, integrable; so
$\hat\tau_{\mathrm{wt}}<\infty$ and $\bar K<\infty$.

Conclusions. \Cref{cor:co-gap,def:PiN,lem:PiN-props,thm:clarkocone} use
only the data ``closed pair with Poincar\'e inequality'' and apply verbatim to
$(D_{\bar\nu},\delta_{\bar\nu})$ with $C_N=\bar K$.

\textnormal{(C)} The probes $\Phi_\varepsilon=\psi_\varepsilon(\eta-\xi)$ of
\cref{prop:co-failure} are flat gap observables,
$\psi_\varepsilon\in\cA^{\flat}_{\mathrm{gap}}$ after restriction to the gap
variable, and all the estimates of that proof (Steps 1--4) are already expressed
in terms of $\bar\nu$ and $\cE_{\bar\nu}$; the closure chain and the
orthogonality argument run verbatim inside the gap pair, whose adjoint
$\delta_{\bar\nu}$ satisfies
$\langle\delta_{\bar\nu}u,K\rangle=\langle u,D_{\bar\nu}K\rangle=0$ and
$\langle\delta_{\bar\nu}u,1\rangle=0$.
\end{proof}

\Cref{thm:gap-co} gives a complete energy-space Clark--Ocone dichotomy,
toggled by the boundary type alone, on the sub-sector where the LJR structure
is one-dimensional and explicit. Its scope is also strictly limited: the gap
pair sees only difference observables of one pair of points, and the
center-of-mass directions, where the diagonal carr\'e lives, re-enter through
\cref{prop:com-half} and require $\COmix$ (or $\COcap$) for the full
two-point sector. The $|\mathbf x|=1$ case of \cref{thm:clarkocone} and
\cref{thm:gap-co} thus bracket the full sector from its two marginal
directions, and what is missing is exactly the joint $(g_T,V)$-mean control
named $\COmix$.

%==============================================================
\section{Main results III: the partitionwise dynamic black operator-It\^o formula}\label{sec:dynamic}
%==============================================================

The third main result produces the increment $f(Z_T)-f(Z_0)$ dynamically from
the energy calculus, in two stages: the Dynkin comparison, which identifies
exactly what a dynamic formula must disintegrate, and the adapted telescope,
which carries the disintegration out at the level of partitions, exactly. The
standing input is stated once: everything below rests on
the slicewise Poincar\'e/Clark--Ocone input
$(\mathrm{Poin}^{\mathrm{dyn}})$ of \cref{def:dynamic-div}, automatic at
$m=1$ and, in phase \textnormal{(R)}, on the gap sector, where the dynamic formula needs nothing further (\cref{cor:dynamic-gap}); at the full
two-point level it is the slicewise form of $\COmix$, and the formula is exactly as strong as that input.

\subsection{The Dynkin comparison: what a dynamic formula must disintegrate}

\begin{proposition}[Dynkin pair and the static-divergence comparison]\label{prop:dynkin}
Fix a sector $\mathbf x$, $|\mathbf x|\le2$, let $Z=(Z_t)_{t\le T}$ be the
$m$-point motion from $\mathbf x$ with Feller generator $L^{(m)}$
\citep{LJR-FCN}, and let $f,f^2\in\Dom(L^{(m)})$, $f$ bounded.
\begin{enumerate}
\item[$(\mathrm i)$] (Dynkin pair; classical.)
$M^f_t:=f(Z_t)-f(Z_0)-\int_0^tL^{(m)}f(Z_s)\,ds$ is a centered
$L^2(\Pth)$-martingale with
\[
  \langle M^f\rangle_t=\int_0^t\Gamma^{(m)}_\theta(f,f)(Z_s)\,ds,
  \qquad
  \E\bigl[(M^f_T)^2\bigr]=\E\!\int_0^T\Gamma^{(m)}_\theta(f,f)(Z_s)\,ds ,
\]
the time-integrated carr\'e, as opposed to the single-time datum
$\E[\GamN]=\E[\Gamma^{(m)}_\theta(f,f)(Z_T)]$ on which the energy pair is
built.
\item[$(\mathrm{ii})$] (Comparison.) Assume in addition the sectorwise
Poincar\'e inequality of \cref{cor:co-gap} on the fixed sector. With
$F:=f(Z_T)$ and $F_0:=F-\E F$, \cref{thm:clarkocone}$(\mathrm{iv})$ applies to
$F\in L^2(\mu_{\mathbf x})$ and, as an identity in $L^2(\Omega,\Pth)$,
\begin{equation}\label{eq:dynkin-compare}
  \divN\bigl(\DN G_NF_0\bigr)
  \;=\;M^f_T\;+\;\int_0^T\Bigl[L^{(m)}f(Z_s)-\E\,L^{(m)}f(Z_s)\Bigr]ds .
\end{equation}
\end{enumerate}
\end{proposition}
\begin{proof}
$(\mathrm i)$ is Dynkin's formula and the carr\'e-du-champ bracket identity for
a Feller process with $f,f^2\in\Dom(L^{(m)})$ (standard; the $m$-point
semigroups are Feller by~\citep{LJR-FCN}). For $(\mathrm{ii})$: centering
Dynkin's formula at $t=T$ (and noting $f(Z_0)$ is deterministic and
$\E M^f_T=0$) gives
$F-\E F=M^f_T+\int_0^T[L^{(m)}f(Z_s)-\E L^{(m)}f(Z_s)]\,ds$, while
\cref{thm:clarkocone}$(\mathrm{iv})$ gives $F-\E F=\divN(\DN G_NF_0)$; the two
right-hand sides agree as elements of $L^2(\Omega,\Pth)$.
\end{proof}

Neither summand on the right of \eqref{eq:dynkin-compare} is
$\sigma(Z_T)$-measurable; only the sum is. The static divergence therefore
aggregates the entire Dynkin pair, martingale plus centered
compensator, into a single one-form acting at the terminal time. A dynamic
black It\^o formula is exactly a filtration-compatible disintegration of the
representer $\DN G_NF_0$: a family of modules $\cU_{[0,t]}$ with compatible
divergences $\delta_t$ such that the aggregate splits into an adapted
$\delta$-integral matching $M^f$ and a drift matching the compensator, with
bracket $\int_0^\cdot\Gamma^{(m)}_\theta(f,f)(Z_s)\,ds$. By blackness
(\cref{thm:black}) the martingale $M^f$ admits no first-chaos representation,
so the disintegration, if it exists, must remain energy-spatial at every time;
nothing in the present single-time construction supplies the time-indexed
module family, and constructing it is \cref{prob:multitime}.

\subsection{The adapted telescope: a partitionwise dynamic formula}
\label{subsec:dynamic}

The disintegration called for by \eqref{eq:dynkin-compare} can be carried out
at the level of partitions, and exactly. The mechanism rests on one
observation about what blackness does and does not forbid. The first chaos
consists of limits of slice-additive functionals, sums
$\sum_kX_k$ with $X_k$ measurable with respect to the slice field
$\cF^\theta_{t_k,t_{k+1}}$ alone; \cref{thm:black} annihilates these.
The martingale differences $f(Z_{t_{k+1}})-\E[f(Z_{t_{k+1}})\mid\cF_{t_k}]$ are
not of this form: they depend on $Z_{t_k}$, i.e.\ on the past. Adapted
increments live exactly in the part of $L^2$ that blackness leaves available,
and the slicewise Clark--Ocone representation reaches them: conditionally on
$Z_{t_k}=z$, the increment is a centered functional on a fresh fixed sector
(start $z$, horizon $\Delta_k$), to which \cref{thm:clarkocone}$(\mathrm{iv})$
applies.

Fix a sector $\mathbf x$, $|\mathbf x|\le2$, and write $\cF_t$ for the
filtration of the $m$-point motion $Z$. For a start $z$ and horizon
$\Delta>0$ let $\mu^{\Delta}_z$ be the law of $Z_\Delta$ from $z$, and let
$(D^{z,\Delta},\delta^{z,\Delta})$ be the fixed-sector closed pair on
$L^2(\mu^{\Delta}_z)$ (the diagonal-flat pair at $m=2$), with Green operator
$G^{z,\Delta}$ when the sector Poincar\'e inequality holds; all the
fixed-sector results of \cref{sec:clarkocone} are uniform in $(z,\Delta)$ in
their assumptions, the time-$T$ sector being the case $(z,\Delta)=(\mathbf x,T)$.

\begin{definition}[Slicewise Poincar\'e; adapted slice one-forms; dynamic divergence]\label{def:dynamic-div}
Let $\pi:0=t_0<\dots<t_n=T$ be a partition, $\Delta_k:=t_{k+1}-t_k$,
$|\pi|:=\max_k\Delta_k$.
\textnormal{(a)} The slicewise Poincar\'e condition $(\mathrm{Poin}^{\mathrm{dyn}})$ is: for every
$\Delta\in(0,T]$ and (law of $Z_t$)-a.e.\ start $z$, the sector pair
$(D^{z,\Delta},\delta^{z,\Delta})$ satisfies
$\mathrm{Var}_{\mu^\Delta_z}\le C(z,\Delta)\,\cE$ with $C(z,\Delta)<\infty$;
no uniformity in $(z,\Delta)$ is required. It holds automatically at
$m=1$ ($C=2\Delta$, \cref{lem:sector-poincare}) and on the gap sector in phase
\textnormal{(R)} ($C=\bar K(z,\Delta)$, \cref{thm:gap-co}, every start
including the atom); at $m=2$ it is supplied by the slicewise form of
$\COmix$.
\textnormal{(b)} An adapted slice one-form is a family
$u=(u_k)_{k<n}$ where $u_k(z)\in\Dom(\delta^{z,\Delta_k})$ for a.e.\ $z$, with
the composites $W_k(u):=\delta^{z,\Delta_k}\bigl(u_k(z)\bigr)\big|_{z=Z_{t_k}}$
(functions of $(Z_{t_k},Z_{t_{k+1}})$) jointly measurable and square
integrable. Its dynamic divergence is
\[
  \delta^{\pi}(u)\;:=\;\sum_{k=0}^{n-1}W_k(u).
\]
Each $W_k(u)$ is a martingale difference,
$\E[W_k(u)\mid\cF_{t_k}]=\langle\delta^{z,\Delta_k}u_k(z),1\rangle\big|_{z=Z_{t_k}}=0$,
so the increments are orthogonal and
\begin{equation}\label{eq:orth-incr}
  \E\bigl[\delta^{\pi}(u)^2\bigr]\;=\;\sum_{k}\E\bigl[W_k(u)^2\bigr],
\end{equation}
the orthogonal-increment identity, the black substitute for the It\^o
isometry: it is supplied by adaptedness and the centering of the slice
divergences, not by any driving noise.
\end{definition}

\begin{definition}[Driven and driverless representations]\label{def:driverless}
An additive integrator for the noise $\cN_\theta$ is a family
$W=(W_{s,t})_{s\le t}$ with $W_{s,t}\in L^2(\cF^{\theta}_{s,t})$ centered,
additive ($W_{s,u}+W_{u,t}=W_{s,t}$ for $s\le u\le t$), and nondegenerate
($\E[W_{0,T}^2]>0$). A partitionwise dynamic representation is driven
if its martingale increments have the form
$W_k=a_k(Z_{t_k})\,W_{t_k,t_{k+1}}$, up to refinement-vanishing remainders,
for a single fixed additive integrator $W$ and predictable coefficients $a_k$, the Wiener case, with $W_{s,t}=B_t-B_s$, and driverless
otherwise. An additive integrator is exactly a nonzero element of the first
chaos of $\cN_\theta$ (a centered, additive-across-time, square-integrable
family), so \cref{thm:black} states that the sticky-flow noise admits
no additive integrator: every partitionwise dynamic representation over
$\cN_\theta$ is necessarily driverless. The term is used in this sense
throughout.
\end{definition}

\begin{theorem}[Partitionwise dynamic black operator-It\^o formula; under $(\mathrm{Poin}^{\mathrm{dyn}})$]\label{thm:dynamic}
Assume \cref{conv:LJR} and the slicewise condition
$(\mathrm{Poin}^{\mathrm{dyn}})$ of \cref{def:dynamic-div} on the fixed
sector (automatic at $m=1$ and, in phase
\textnormal{(R)}, on the gap sector; the slicewise form of $\COmix$ at the
full two-point level). Let $f\in\Dom(L^{(m)})$ be bounded. For a partition $\pi$ define the
canonical adapted integrand
\[
  u^{f,\pi}_k(z)\;:=\;D^{z,\Delta_k}\,G^{z,\Delta_k}\bigl(f-P_{\Delta_k}f(z)\bigr)
  \;\in\;\Dom(\delta^{z,\Delta_k}),
  \qquad P_\Delta f(z):=\E_z[f(Z_\Delta)] ,
\]
the conditional minimal-norm Green representer of the slice increment (no
flatness of $f$ is required: the representation of
\cref{thm:clarkocone}$(\mathrm{iv})$ covers every centered element of
$L^2(\mu^{\Delta_k}_z)$). Then:
\begin{enumerate}
\item[$(\mathrm i)$] (Exact telescope.) For every partition $\pi$,
as an identity in $L^2(\Omega,\Pth)$,
\begin{equation}\label{eq:dynamic-exact}
  f(Z_T)-f(Z_0)\;=\;\delta^{\pi}\bigl(u^{f,\pi}\bigr)
  \;+\;\sum_{k=0}^{n-1}\int_0^{\Delta_k}\bigl(P_sL^{(m)}f\bigr)(Z_{t_k})\,ds .
\end{equation}
\item[$(\mathrm{ii})$] (Bracket.) The dynamic divergence term obeys
\eqref{eq:orth-incr} with
$\E[W_k^2]=\E\bigl[\mathrm{Var}\bigl(f(Z_{t_{k+1}})\mid\cF_{t_k}\bigr)\bigr]$;
when moreover $f$ lies in the slice form domains, each summand is computed by
the energy pairing,
$\E[W_k^2]=\E\bigl[\langle\Pi^{z,\Delta_k}D^{z,\Delta_k}f,\
D^{z,\Delta_k}f\rangle\big|_{z=Z_{t_k}}\bigr]$
(\cref{thm:clarkocone}$(\mathrm{iii})$): no quadratic variation enters.
\item[$(\mathrm{iii})$] (Mesh limit.) As $|\pi|\to0$ the compensator in
\eqref{eq:dynamic-exact} converges in $L^2(\Pth)$ to
$\int_0^TL^{(m)}f(Z_s)\,ds$, hence $\delta^{\pi}(u^{f,\pi})$ converges in
$L^2(\Pth)$ to the Dynkin martingale value $M^f_T$ of \cref{prop:dynkin}, and
\begin{equation}\label{eq:dynamic-limit}
  f(Z_T)-f(Z_0)\;=\;\lim_{|\pi|\to0}\delta^{\pi}\bigl(u^{f,\pi}\bigr)
  \;+\;\int_0^TL^{(m)}f(Z_s)\,ds ,
\end{equation}
with, if also $f^2\in\Dom(L^{(m)})$,
\[
  \E\Bigl[\bigl(\lim_{|\pi|\to0}\delta^{\pi}(u^{f,\pi})\bigr)^{2}\Bigr]
  \;=\;\E\!\int_0^T\Gamma^{(m)}_\theta(f,f)(Z_s)\,ds .
\]
\end{enumerate}
\end{theorem}
\begin{proof}
$(\mathrm i)$ Telescope $f(Z_T)-f(Z_0)=\sum_k\bigl[f(Z_{t_{k+1}})-
\E[f(Z_{t_{k+1}})\mid\cF_{t_k}]\bigr]+\sum_k\bigl[\E[f(Z_{t_{k+1}})\mid\cF_{t_k}]
-f(Z_{t_k})\bigr]$. For the first sum: by the Markov property the conditional
law of $Z_{t_{k+1}}$ given $\cF_{t_k}$ is $\mu^{\Delta_k}_{Z_{t_k}}$, and by
$(\mathrm{Poin}^{\mathrm{dyn}})$ and \cref{thm:clarkocone}$(\mathrm{iv})$
applied on the slice sector with start $z=Z_{t_k}$,
\[
  f(Z_{t_{k+1}})-\E[f(Z_{t_{k+1}})\mid\cF_{t_k}]
  \;=\;\delta^{z,\Delta_k}\bigl(u^{f,\pi}_k(z)\bigr)\Big|_{z=Z_{t_k}}
  \;=\;W_k\bigl(u^{f,\pi}\bigr)
\]
(the composite $W_k=f(Z_{t_{k+1}})-(P_{\Delta_k}f)(Z_{t_k})$ is manifestly
jointly measurable). For the second sum: $f\in\Dom(L^{(m)})$ gives the Feller
identity $\E[f(Z_{t_k+\Delta})\mid\cF_{t_k}]-f(Z_{t_k})
=\int_0^{\Delta}(P_sL^{(m)}f)(Z_{t_k})\,ds$.
$(\mathrm{ii})$ Orthogonality is \eqref{eq:orth-incr};
$W_k=f(Z_{t_{k+1}})-\E[f(Z_{t_{k+1}})\mid\cF_{t_k}]$ gives the conditional
variance; the energy pairing is the slicewise instance of
\cref{thm:clarkocone}$(\mathrm{iii})$.
$(\mathrm{iii})$ Set $h:=L^{(m)}f$, which lies in $C_0$ since the Feller
generator maps its domain into $C_0$. The difference between the
compensator and $\int_0^Th(Z_s)\,ds$ is
$-\sum_k\int_0^{\Delta_k}\xi_{k,s}\,ds$ with
$\xi_{k,s}:=h(Z_{t_k+s})-\E[h(Z_{t_k+s})\mid\cF_{t_k}]$, and
\[
  \E[\xi_{k,s}^2]=\E\bigl[\mathrm{Var}(h(Z_{t_k+s})\mid\cF_{t_k})\bigr]
  \le\E\bigl[(h(Z_{t_k+s})-h(Z_{t_k}))^2\bigr]
  \le\sup_z\E_z\bigl[(h(Z_s)-h(z))^2\bigr]=:\eta(s),
\]
where $\eta(s)=\sup_z[(P_sh^2)(z)-2h(z)(P_sh)(z)+h^2(z)]
\le\|P_sh^2-h^2\|_\infty+2\|h\|_\infty\|P_sh-h\|_\infty\to0$ as $s\downarrow0$
by Feller strong continuity; set
$\bar\eta(\delta):=\sup_{s\le\delta}\eta(s)$, which also tends to $0$ as
$\delta\downarrow0$ ($\eta$ itself need not be monotone). By the triangle
inequality the $L^2$ norm of the
difference is at most $\sum_k\Delta_k\sqrt{\bar\eta(|\pi|)}=T\sqrt{\bar\eta(|\pi|)}\to0$
(cross-slice orthogonality of the $\xi_{k,s}$ would sharpen this to
$(T|\pi|\,\bar\eta(|\pi|))^{1/2}$; the simpler bound suffices).
Since the left side of \eqref{eq:dynamic-exact} does not depend on $\pi$,
$\delta^{\pi}(u^{f,\pi})\to f(Z_T)-f(Z_0)-\int_0^TL^{(m)}f(Z_s)\,ds=M^f_T$ in
$L^2$, and the bracket of $M^f_T$ is \cref{prop:dynkin}$(\mathrm i)$.
\end{proof}

\begin{corollary}[The dynamic formula on the gap sector]\label{cor:dynamic-gap}
In phase \textnormal{(R)}, for the autonomous gap $g_t$ and every bounded
$f\in\Dom(L_\gamma)$ (the sticky gap generator), the conclusions of
\cref{thm:dynamic} hold with no input beyond \cref{conv:LJR}: the slicewise Poincar\'e
constants are the $\bar K(z,\Delta)<\infty$ of \cref{thm:gap-co}, finite for
every start $z\in\R$ (including the atom $z=0$) and every horizon. Moreover
the drift is explicit, the Wentzell boundary term being clocked by the sticky
occupation time:
\[
  \int_0^TL_\gamma f(g_s)\,ds
  \;=\;\int_0^Tf''(g_s)\,\mathbf 1_{\{g_s\neq0\}}\,ds
  \;+\;\gamma^{-1}\bigl[f'(0{+})-f'(0{-})\bigr]\,\mathrm{Occ}_T ,
\]
with $\mathrm{Occ}_T$ the occupation time of \cref{prop:occupation}.
\end{corollary}
\begin{proof}
The gap is an autonomous Feller Markov process (\cref{lem:gap-sticky}), its
slice sectors are the gap pairs of \cref{thm:gap-co} with start $z$ and
horizon $\Delta$, and the proof of \cref{thm:gap-co} used only
\cref{prop:sticky-bm}, whose assumptions hold for every start and horizon in
phase \textnormal{(R)}; so $(\mathrm{Poin}^{\mathrm{dyn}})$ holds and
\cref{thm:dynamic} applies. The drift formula is the form of $L_\gamma$:
$L_\gamma f=f''$ off $0$ and
$L_\gamma f(0)=\gamma^{-1}[f'(0{+})-f'(0{-})]$ (\cref{rem:dtrace-gamma2}),
and $\int_0^T\mathbf 1_{\{g_s=0\}}\,ds=\mathrm{Occ}_T$.
\end{proof}

\begin{remark}[What is black about \cref{thm:dynamic}]\label{rem:driverless}
\textnormal{(1)} Driverless, in the sense of \cref{def:driverless}. In
the Wiener case the slicewise
representations collapse onto integration against the increments of a single
driving martingale: the predictable integrand varies, the driver does not, and
the driver is an additive integrator, i.e.\ first chaos. Here no additive
integrator exists, that is exactly \cref{thm:black} read through
\cref{def:driverless}, and the dynamic object is instead an orthogonal sum of
state-indexed conditional divergences $\delta^{z,\Delta_k}$, the
operator itself moving with $Z_{t_k}$. The orthogonal-increment identity
\eqref{eq:orth-incr} replaces the It\^o isometry, and the conditional Green
representers replace the predictable projection.
\textnormal{(2)} Compatibility with blackness. \Cref{thm:black} kills
slice-additive functionals; the increments $W_k$ are conditionally centered
but past-dependent, hence invisible to the first-chaos criterion. The dynamic
formula is built inside the part of $L^2$ that blackness spares.
\textnormal{(3)} Beyond the theorem. The continuum adapted module
$\bigcup_t\,\cU_{[0,t]}$ with a single divergence, the convergence of the
integrands $u^{f,\pi}$ themselves as $|\pi|\to0$ (only the divergences
converge, their limit identified through the compensator), pathwise chain
rules, and assembly across sectors all remain open
(\cref{prob:multitime}). For the gap
sector a classical semimartingale It\^o formula exists as well, the sticky gap
being a semimartingale; the content of \cref{cor:dynamic-gap} is that the
formula is reproduced by the intrinsic energy calculus without invoking the
driving Brownian motion, the only form that can survive at flow level,
where no driver exists.
\end{remark}

%==============================================================
\section{Open problems}\label{sec:openproblems}
%==============================================================

Four problems remain, each stated with the missing element identified.

\begin{problem}[Full reflecting case]\label{prob:fullreflecting}
Beyond the diagonal-flat sub-core the moving-atom trace is second-order
(\cref{rem:dtrace-gamma2}), and $(\mathrm{CKC}_2^{\flat})$ no longer governs it.
The enhanced energy $E_{H2}(\Phi)=E_{\mathrm{fo}}(\Phi)+\lambda\,\|L\Phi\|^2$
satisfies the five closed-pair axioms and produces a closed pair, so the
obstruction lies not in the construction but in the contraction:
\textnormal{(P3)} for the enhanced energy holds only if conditional expectation
intertwines with the $H^2$ structure, which is the time-$T$-law versus
motion-semigroup reference-measure question (\cref{rem:reference-measure}). Decide
this intertwining, and with it the full reflecting integration-by-parts identity.
(\Cref{prop:agmon} now discharges the trace condition \cref{hyp:dtrace} on the
second-order sub-core with the explicit constant
$C_{\mathrm{tr}}\sim1/\inf p$, and \cref{rem:order-mismatch} establishes the
mixed deletion inequality; the open content is thereby isolated to the
intertwining alone.)
\end{problem}

\begin{problem}[Compatibility at $m\ge3$]\label{prob:m3}
The pairwise sector is protected because the only stratum reached under $2\to1$
deletion is codimension one, hence non-polar in phase (R); at $m\ge3$ the
inclusion $\{1\}\subset\{1,2,3\}$ meets a multi-coalescence stratum of codimension
$\ge2$ with two $\Gamma$-blind difference directions, and \eqref{eq:p3} plausibly
fails, a difference-direction probe $g_n(\eta_1-\xi,\eta_2-\xi)$ sending
$\cE_{\mathbf y}(g_n)\to0$ while the atom-weight term keeps
$\cE_{\mathbf x}(Pg_n)\not\to0$. Deciding this requires the $m\ge3$ time-$T$
near-multi-coalescence law $m_{\theta,T,\mathbf x}$, not characterized
in~\citep{LJR-FCN,LJR-circle-noise}; the deciding quantity is whether the apex is polar
for the degenerate $\mu_y$-weighted gap energy, a transverse-density threshold.
(\Cref{prop:m3-blocked} now shows that, under the pairwise-stratum regularity
$(\mathrm R_3)$ of \cref{hyp:r3}, this polarity mechanism cannot close:
the one-dimensional partial-coincidence strata short-circuit the
codimension-two apex and give it positive capacity. The problem therefore
shifts: establish $(\mathrm R_3)$ from the LJR construction, and decide whether
a failure at $m\ge3$ can instead be driven by the score and atom-weight terms,
the only channels the capacity argument does not block.)
\end{problem}

\begin{problem}[The mixed-mean bound, and the sharp sector gap]\label{prob:cocap}
\Cref{thm:clarkocone} runs on either the uniformization $\COcap$
(\cref{def:cocap}), automatic in the translation-invariant reference, but
expected to fail under the point-started reference by the two-well
mechanism of \cref{rem:two-well}, or on the mixed-mean bound $\COmix$
(\cref{def:comix}), which the gap-path route of \cref{prop:com-half} isolates
as the operative condition. Prove $\COmix$ for the point-started two-point law
of \cref{conv:LJR}: by \cref{prop:com-half}$(\mathrm{iii})$ the
$V$-direction of the conditional mean $H_\Phi(g_T,V)$ is first-order, and by
\cref{thm:gap-co} the pure $g_T$-direction is controlled outright, so
what is missing is joint control of the planar $(g_T,V)$-law, e.g.\ quantitative
lower bounds on the conditional occupation density.
\Cref{rem:comix-instantiation} now tests the bound itself,
$C_{\mathrm{mix}}\approx2$ at $g_0=T=1$ across $\gamma\in[0.05,0.5]$, with the
dominant channel the linear gap mode already controlled by \cref{thm:gap-co}
and the planar excess of order fifteen percent at the stickier parameters, so
the missing planar control is quantitatively a perturbation of the
gap-sector bound; what remains is its proof, and the certified
(interval-arithmetic) form. Substantiate (or refute)
the expected failure of essential-supremum $\COcap$ via the two-well mechanism
of \cref{rem:two-well}, whose occupation-time element is now the exact law
\eqref{eq:occupation-law}; what is missing are the conditional-density valley
bounds. Determine
the sharp sector spectral gap $\lambda_N$ and its degeneration toward the
coalescing endpoint, where \cref{prop:co-failure} exhibits $\lambda_N=0$ with
an explicit kernel element; \cref{rem:comix-instantiation} locates
$\lambda_N\approx0.2$--$0.5$ at $g_0=T=1$, commensurate with the one-point
gap $1/(2T)$, so the open content is the proof and the rate of collapse
toward coalescence. Instantiate $(\mathrm{CKC}_2^{\flat})$: exhibit one
explicit $\theta>0$ with $c^{\flat}(\xi)\le1$ $\mu_x$-a.e.\ (a numerical
certificate against the sticky-gap law would suffice); by
\cref{prop:invariant-fails} the translation-invariant reference is closed
outright --- \textnormal{(P3)} fails there on the flat sector, the
center-of-mass channel being in deficit --- so the instantiation must run
through the point-started reference, where
\cref{rem:instantiation} now verifies the contraction itself numerically
($\sigma\in[0.89,0.97]$ across families and solvers at $g_0=T=1$,
$\gamma\in[0.05,0.5]$, in every case below $1$); what remains
is the certified (interval-arithmetic) form, and the $c^{\flat}$-version
with an explicit margin $\kappa_\xi>0$. Separately, decide whether the sectorwise
representers of \cref{thm:clarkocone} admit any consistent assembly across
sectors, which \cref{rem:sectorwise} leaves without a natural Hilbert frame.
\end{problem}

\begin{problem}[Dynamic extension]\label{prob:multitime}
The identity \eqref{eq:black-ito} is single-time, and its divergence term is
centered, so it cannot produce the increment $f(U_T)-f(U_0)$
(\cref{rem:ito-form}). A dynamic, pathwise black It\^o formula would require the
multi-time family $\{U_t\}_{t\le T}$ and its bracket, under a strictly stronger
closability condition. Establish or refute it. (Partially resolved:
\cref{thm:dynamic} constructs the disintegration at the level of partitions,
exactly, the increment is an adapted sum of conditional energy divergences
plus the generator drift, with the martingale part recovered in the mesh
limit and the bracket the time-integrated carr\'e; on the gap sector the
formula holds with no further input, \cref{cor:dynamic-gap}. What remains open is the
continuum object: an inductive-limit adapted module on which the integrands
$u^{f,\pi}$ themselves converge, a single divergence operator whose domain
contains their limit, and the pathwise form; \cref{rem:driverless}(3).)
\end{problem}

%==============================================================
\section{Conclusion}\label{sec:conclusion}
%==============================================================

The noise of the Le Jan--Raimond Brownian sticky flow is black, and
blackness removes the classical route to stochastic calculus at its first
step: there is no first chaos, hence no martingale representation, no
predictable projection, and no driver to integrate against. This paper
develops the calculus that survives the removal by reversing the usual
order of construction: the operator derivative $D$ on a carr\'e-du-champ
energy space comes first, the divergence is its adjoint $\delta=D^{*}$,
and everything downstream --- integration by parts, Clark--Ocone--Green
representation, dynamic formula --- is built from the adjoint pair on
fixed one- and two-point sectors of the flow.

Three structural conclusions emerge. First, what decides whether the
construction lives is not blackness but the Feller boundary type of the
two-point gap: at every reflecting-regular boundary the deletion
contraction is unobstructed and the diagonal-flat pairwise sector carries
closed operator pairs and a pairwise black integration-by-parts identity
under $(\mathrm{CKC}_2^{\flat})$, while at the coalescing endpoint a
difference-direction probe defeats the contraction outright; both phases
carry black noises, so the boundary type is the true discriminant
(\cref{thm:discriminant}). Second, representation survives as a Green
identity rather than a projection: $F-\E F=\divN(\PiN\DN F)$ with a
minimal-norm representer in the energy module, with explicit
constants on the autonomous gap sector, and failing at the coalescing
endpoint, where the coalescence indicator is a nonconstant energy-null
functional orthogonal to every divergence
(\cref{thm:gap-co,prop:co-failure}). Third, dynamics survives
partitionwise: the increment of an observable is an exact adapted sum of
conditional slicewise divergences plus the generator drift, with the
Dynkin martingale recovered in the mesh limit --- blackness forbids a
driving noise, not adapted increments, so the formula is driverless rather
than impossible (\cref{thm:dynamic}).

The boundary of what is proved is equally sharp, and
\cref{sec:openproblems} records it: the full reflecting case beyond the
diagonal-flat sub-core turns on a single semigroup-intertwining question;
the contraction at three and more points is protected against the polarity
route but not yet decided; the mixed-mean bound $\COmix$ awaits joint
control of the planar law of the point-started reference; and the
continuum limit of the partitionwise formula, a single divergence on an
inductive-limit adapted module, remains the natural next object. The
pairwise sector, where the reduction is complete, is small enough to be
decided and large enough to carry a calculus; that combination, rather
than any generality, is what the paper establishes.

\section*{Acknowledgements}
AI tools were used as interactive assistants for drafting, reorganizing,
and generating mathematical arguments, under the author's direction. The
author is responsible for all mathematical content and any errors.

%==============================================================
\appendix
%==============================================================

\section{Derivation of the amended gradient identity}\label{app:amended}

We prove \cref{prop:amended-star}. With $\kappa$ as in \eqref{eq:kernel} and
$\langle\Phi\rangle_p=\int\Phi(\xi,\eta)p(\xi,\eta)\,d\eta$,
\[
  P\Phi(\xi)=q(\xi)\Phi(\xi,\xi)+(1-q(\xi))\langle\Phi\rangle_p(\xi).
\]
Differentiate the atomic and diffuse parts separately. The atomic part:
since the atom location $\eta=\xi$ moves with $\xi$,
\[
  \partial_\xi\bigl[q\Phi(\xi,\xi)\bigr]
  =q'\Phi(\xi,\xi)+q\,\tfrac{d}{d\xi}\Phi(\xi,\xi)
  =q'\Phi(\xi,\xi)+q\,[(\partial_\xi\Phi)(\xi,\xi)+(\partial_\eta\Phi)(\xi,\xi)],
\]
where $\tfrac{d}{d\xi}\Phi(\xi,\xi)=(\partial_\xi+\partial_\eta)\Phi$ is the
center-of-mass derivative at the atom. The diffuse part:
\[
  \partial_\xi\bigl[(1-q)\langle\Phi\rangle_p\bigr]
  =-q'\langle\Phi\rangle_p+(1-q)\bigl[\langle\partial_\xi\Phi\rangle_p
   +\textstyle\int\Phi\,\partial_\xi p\,d\eta\bigr].
\]
Now $\langle\partial_\xi\Phi\rangle_\xi=q(\partial_\xi\Phi)(\xi,\xi)
+(1-q)\langle\partial_\xi\Phi\rangle_p$. Subtracting,
\[
  \partial_\xi(P\Phi)-\langle\partial_\xi\Phi\rangle_\xi
  =q'[\Phi(\xi,\xi)-\langle\Phi\rangle_p]
   +(1-q)\textstyle\int\Phi\,\partial_\xi p\,d\eta
   +q(\partial_\eta\Phi)(\xi,\xi),
\]
which is \eqref{eq:amended-star}. The bracket $\textnormal{(B)}+\textnormal{(C)}$
is centered: for $\Phi\equiv1$, $\textnormal{(B)}=q'[1-1]=0$ and
$\textnormal{(C)}=(1-q)\int\partial_\xi p=(1-q)\partial_\xi\!\int p=0$
(using $\int p\equiv1$), so it is a covariance functional; while
$\textnormal{(D)}=q\,\partial_\eta\Phi(\xi,\xi)$ is a derivative of $\Phi$,
not of covariance form, the extra moving-atom term. At
$\theta=0$, $q\equiv0$ and $p$ is $\xi$-independent, so
$\textnormal{(B)}=\textnormal{(C)}=\textnormal{(D)}=0$. \qed

\section{Term-by-term assembly of \texorpdfstring{$\CKC$}{CKC}}\label{app:assembly}

Trace inequalities \eqref{eq:traces}. For the value trace, write
$\Phi(\xi,\xi)-\langle\Phi\rangle_p=-\int(\int_\xi^\eta\partial_\eta\Phi\,ds)\,p\,d\eta$
and apply Cauchy--Schwarz in $s$ with resistance
$r(\xi,\eta)=|\int_\xi^\eta ds/p(\xi,s)|$:
$(\int_\xi^\eta\partial_\eta\Phi\,ds)^2\le r(\xi,\eta)\langle(\partial_\eta\Phi)^2\rangle_p$,
giving $|\Phi(\xi,\xi)-\langle\Phi\rangle_p|^2\le\tau_{\mathrm{wt}}G_\eta$
with $\tau_{\mathrm{wt}}=(\int\sqrt r\,p\,d\eta)^2$. The diffuse-score bound
is $\int\Phi\,\partial_\xi p=\mathrm{Cov}_p(\Phi,\partial_\xi\log p)$ (since
$\int\partial_\xi p=0$) followed by Cauchy--Schwarz and the Poincar\'e
inequality $\mathrm{Var}_p(\Phi)\le C^p_P\,G_\eta$ (these three are the
value/first-order pieces, finite under \cref{lem:checkB}(R)). The
point-derivative trace $\tau_{\mathrm{at}}$ is not of this type: point
evaluation of $\partial_\eta\Phi$ is discontinuous for the first-order energy
$G_\eta$, so the bound $|\partial_\eta\Phi(\xi,\xi)|^2\le\tau_{\mathrm{at}}\,
\mathcal E_{\eta,\xi}(\Phi)$ requires the enhanced ($H^2$/Agmon) norm
$\mathcal E_{\eta,\xi}$ and finiteness of $\tau_{\mathrm{at}}$ is exactly
\cref{hyp:dtrace}, not implied by \cref{lem:checkB}(R). The
assembly below bounds the $\tau_{\mathrm{at}}$ terms against
$\mathcal E_{\eta,\xi}$ (under \cref{hyp:dtrace}) and the remaining terms
against $G_\eta$; the finite/infinite verdict is unchanged, since in phase (C)
even the value pieces diverge (\cref{cor:dichotomy}(C)).

Assembly. By $(a+b+c)^2\le3(a^2+b^2+c^2)$, $R^2\le3S_0G_\eta$ with
$S_0=q'^2\tau_{\mathrm{wt}}+q^2\tau_{\mathrm{at}}+(1-q)^2I_pC^p_P$. The
cross term in \eqref{eq:obstruction} is split by Young,
$-2\langle\partial_\xi\Phi\rangle_\xi R\ge-\delta\langle\partial_\xi\Phi\rangle_\xi^2-\delta^{-1}R^2$;
the $\delta\langle\partial_\xi\Phi\rangle_\xi^2$ is absorbed by the
center-of-mass/diagonal margin provided $\delta\le\kappa_\xi$, and the
$\delta^{-1}R^2$ inflates the $R^2$-coefficient to $1+\delta^{-1}$. Optimizing
$\delta=\kappa_\xi$ gives the factor $1+1/\kappa_\xi$. The diffuse paying
budget is $\tfrac12(1-q)G_\eta$, so $D(\Phi)\ge0$ follows from
$3(1+1/\kappa_\xi)S_0\le\tfrac12(1-q)$, i.e.
$c(\xi)=\tfrac{6(1+1/\kappa_\xi)}{1-q}S_0\le1$. The strict positivity of
the margin, $\kappa_\xi>0$, is part of the condition and is itself a
statement about the law: the absorption charges the center-of-mass
channel, whose budget the conditional expectation can exhaust, and
\cref{prop:invariant-fails} exhibits the translation-invariant reference
as a case where that channel is in outright deficit and no choice of
$\delta$ closes the assembly. Finally, the
difference-direction half of the atom's $\xi$-energy deficit not covered by the
Jensen reference is, after the same trace bound
\eqref{eq:traces}, of the form $\mathrm{const}\cdot q(\xi)\,\tau_{\mathrm{at}}(\xi)\,
G_\eta$; absorbing this constant into a single universal $A>0$ yields the extra
summand $A\,q\,\tau_{\mathrm{at}}$, so that $S_A=S_0+A\,q\,\tau_{\mathrm{at}}$
with $S_0$ the three principal terms. The constants $3,6,A$ are not
optimized. \qed

\section{Sticky-BM boundary facts used}\label{app:stickyBM}

The body uses a small set of one-dimensional sticky-Brownian-motion facts about
the gap process. We isolate them in a single proposition so that no statement of
the paper rests on an unlabelled ``standard'' marker. The facts are classical;
we cite rather than re-derive, and the exact closed forms of the atom mass and
diffuse density are not used anywhere.

\begin{proposition}[Sticky-BM gap law and boundary behaviour]\label{prop:sticky-bm}
Let $g$ be two-sided sticky Brownian motion on $\R$ (diffusivity $2$, stickiness
$\gamma\in[0,\infty]$) with speed measure $m(dg)=dg+\gamma\,\delta_0$ and
symmetric transition kernel $u(T,\cdot,\cdot)$, the gap process of
\cref{lem:gap-com}. (The atom of the speed measure grows with the stickiness:
$\gamma=0$ is the independent endpoint, $\gamma=\infty$ the absorbing one,
consistently with \cref{lem:gap-sticky}.) For each $T>0$ and each start $g_0$:
\begin{enumerate}
\item[$(\mathrm i)$] (atom $+$ diffuse decomposition). The time-$T$ law is
$\mathcal L(g_T)=q\,\delta_0+p(\cdot)\,dg$, with atom mass
$q=\gamma\,u(T,g_0,0)\in[0,1)$ and continuous diffuse density $p(g)=u(T,g_0,g)$
of Gaussian tail; in particular $C^p_P$ and $I_p$ (the weighted Poincar\'e and
information constants of \cref{lem:checkA}) are finite.
\item[$(\mathrm{ii})$] (phase R, reflecting regular,
$\gamma\in(0,\infty)$). The boundary $0$ is regular-reflecting, the atom is
genuine ($q>0$), and the diffuse density does not vanish at the wall:
$p(0\pm)>0$.
\item[$(\mathrm{iii})$] (phase C, exit/absorbing, $\gamma=\infty$
coalescing). The boundary $0$ is an exit point, $p(0\pm)=0$ with linear vanishing
$p(g)\sim C|g|$ as $g\to0$, and the apex carries zero weighted capacity in the
sense used in \cref{lem:checkB}.
\end{enumerate}
Only the boundary behaviour in $(\mathrm{ii})$--$(\mathrm{iii})$ is used in the
body; the assignment of $\gamma\in(0,\infty)$ to phase R and $\gamma=\infty$ to
phase C is the LJR reflecting/exit classification (\cref{def:trichotomy} and the
reflecting-regular and exit/absorbing cases of~\citep{LJR-FCN}).
\end{proposition}

\begin{proof}[Provenance]
The speed-measure form $m(dg)=dg+\gamma\,\delta_0$ and the resulting
atom-plus-density structure of the one-dimensional resolvent are the standard
generalized-diffusion description of sticky Brownian
motion~\citep[\S4.11]{ItoMcKean},~\citep{EngelbertPeskir}; the Gaussian tail of $p$
and finiteness of $C^p_P,I_p$ follow from the heat-kernel bounds there. The
regular/exit dichotomy of the endpoint $0$, and the corresponding non-vanishing
($p(0\pm)>0$) versus linear vanishing ($p(g)\sim C|g|$) of the diffuse density,
are Feller's boundary classification applied to $m$; the identification of the
two relevant cases with the LJR phases is~\citep[\S9]{LJR-FCN}; the sticky-flow
realization is~\citep{HowittWarren}. At finite $T$ the atom mass stays bounded
away from $1$, $q=\Prob[\text{two Brownian points met by }T]<1$ uniformly in
$\gamma$ (\cref{rem:struct-vs-cert}).
\end{proof}

One further fact, the exact law of near-total occupation, is used by
\cref{rem:two-well}. Occupation-time laws for one-dimensional diffusions are
classical territory~\citep{SalminenStenlund}, and the time-change realization
below is Warren's construction of sticky Brownian motion~\citep{Warren97}. The
identity is surely known in some form, but we could not locate the exact
normalization needed, and the derivation is short.

\begin{proposition}[Exact law of near-total occupation]\label{prop:occupation}
Let $g$ be the two-sided sticky Brownian motion of \cref{prop:sticky-bm}
(diffusivity $2$, stickiness $\gamma\in(0,\infty)$, speed measure
$m=dg+\gamma\,\delta_0$), started at $g_0\in\R$, and let
$\mathrm{Occ}_T=|\{s\le T:g_s=0\}|$. Then for every $0<\varepsilon<T$,
\begin{equation}\label{eq:occupation-law}
  \Prob_{g_0}\bigl[\mathrm{Occ}_T\ge T-\varepsilon\bigr]
  \;=\;2\,\overline\Phi\!\left(\frac{|g_0|+2\gamma^{-1}(T-\varepsilon)}
  {\sqrt{2\varepsilon}}\right),
\end{equation}
$\overline\Phi$ the standard normal tail. In particular, with
$\kappa:=|g_0|+2\gamma^{-1}T$, as $\varepsilon\downarrow0$
\begin{equation}\label{eq:occupation-asymp}
  \Prob_{g_0}\bigl[\mathrm{Occ}_T\ge T-\varepsilon\bigr]
  \;=\;\frac{2\sqrt\varepsilon}{\sqrt{\pi}\,\kappa}\;
  e^{\kappa/\gamma}\;e^{-\kappa^{2}/4\varepsilon}\,\bigl(1+o(1)\bigr):
\end{equation}
the near-total-occupation tail is Gaussian-small in $1/\varepsilon$ (with a
$\sqrt\varepsilon$ prefactor), not of polynomial order. The identity
\eqref{eq:occupation-law} is exact for all $\varepsilon\in(0,T)$, is increasing
in $\gamma$, and degenerates correctly at both endpoints: as
$\gamma\to\infty$ (absorbing) it converges to
$2\overline\Phi(|g_0|/\sqrt{2\varepsilon})=\Prob[\tau_0\le\varepsilon]$, the
absorbed answer, and as $\gamma\to0$ (independent) it vanishes.
\end{proposition}
\begin{proof}
Step 0: time-change realization. Let $W_t=g_0+\sqrt2\,B_t$ ($B$ a
standard Brownian motion), let $\ell_t$ be the occupation density of $W$ at $0$
with respect to Lebesgue measure
($\int_0^tf(W_u)\,du=\int f(y)\,\ell_t(y)\,dy$), and set
\[
  A_t:=t+\gamma\,\ell_t(0),\qquad \alpha:=A^{-1},\qquad g_s:=W_{\alpha_s}.
\]
$A$ is continuous and strictly increasing with $A_0=0$, $A_\infty=\infty$, so
$\alpha$ is well defined; $g$ is the quasidiffusion with scale $x$ and speed
$m=dg+\gamma\delta_0$, i.e.\ the sticky Brownian motion of
\cref{prop:sticky-bm} (Warren's construction~\citep{Warren97};
see also~\citep[\S4.11]{ItoMcKean},~\citep{EngelbertPeskir}).

Step 1: occupation identity. $\mathrm{Occ}_s=\gamma\,\ell_{\alpha_s}(0)$.
Indeed, substituting $u=A_t$ ($du=dt+\gamma\,d\ell_t$),
\[
  \mathrm{Occ}_s=\int_0^s\mathbf 1_{\{W_{\alpha_u}=0\}}\,du
  =\int_0^{\alpha_s}\mathbf 1_{\{W_t=0\}}\,dt
  +\gamma\int_0^{\alpha_s}\mathbf 1_{\{W_t=0\}}\,d\ell_t
  =0+\gamma\,\ell_{\alpha_s}(0),
\]
since the zero set of $W$ is Lebesgue-null a.s.\ while $d\ell$ is carried by it.

Step 2: deterministic equivalence. From $T=A_{\alpha_T}=
\alpha_T+\gamma\ell_{\alpha_T}(0)=\alpha_T+\mathrm{Occ}_T$ and the monotonicity
of $A$,
\[
  \{\mathrm{Occ}_T\ge T-\varepsilon\}
  =\{\alpha_T\le\varepsilon\}
  =\{A_\varepsilon\ge T\}
  =\bigl\{\ell_\varepsilon(0)\ge\gamma^{-1}(T-\varepsilon)\bigr\}:
\]
near-total occupation is exactly a local-time tail event for the
underlying Brownian motion.

Step 3: the local-time tail. Scaling: $\ell_t(0)=
\tfrac1{\sqrt2}L_t(c)$, where $L(c)$ is the occupation density of the standard
$B$ at the level $c:=-g_0/\sqrt2$. Let $\tau_c$ be the hitting time of $c$ and,
for the post-$\tau_c$ Brownian motion (a standard Brownian motion at its
starting level, independent of $\mathcal F_{\tau_c}$), let $\sigma_a$ be the
inverse local time at that level; then
$\{L_\varepsilon(c)\ge a\}=\{\tau_c+\sigma_a\le\varepsilon\}$. By L\'evy's
identity $(L_t(0))_{t}\overset{d}{=}(\max_{s\le t}B_s)_t$, the inverse local
time has the law of the first-passage process, $\sigma_a\overset d=T_a$, and
the first-passage process $(T_y)_{y\ge0}$ is a subordinator (strong Markov),
so
\[
  \tau_c+\sigma_a\;\overset d=\;T_{|c|}+T'_a\;\overset d=\;T_{|c|+a},
  \qquad
  \Prob\bigl[L_\varepsilon(c)\ge a\bigr]
  =\Prob\bigl[T_{|c|+a}\le\varepsilon\bigr]
  =2\,\overline\Phi\!\left(\frac{|c|+a}{\sqrt\varepsilon}\right)
\]
by the reflection principle. With $a=\sqrt2\,\gamma^{-1}(T-\varepsilon)$ and
$|c|=|g_0|/\sqrt2$,
\[
  \Prob_{g_0}\bigl[\mathrm{Occ}_T\ge T-\varepsilon\bigr]
  =2\,\overline\Phi\!\left(\frac{|g_0|/\sqrt2+\sqrt2\,\gamma^{-1}(T-\varepsilon)}
  {\sqrt\varepsilon}\right),
\]
which is \eqref{eq:occupation-law}. The asymptotics
\eqref{eq:occupation-asymp} is $\overline\Phi(x)=(1+O(x^{-2}))\,\varphi(x)/x$
applied at $x=\kappa_\varepsilon/\sqrt{2\varepsilon}$,
$\kappa_\varepsilon=\kappa-2\gamma^{-1}\varepsilon$, using
$\kappa_\varepsilon^2/4\varepsilon=\kappa^2/4\varepsilon-\kappa/\gamma
+O(\varepsilon)$. The endpoint degenerations are immediate from
\eqref{eq:occupation-law}.
\end{proof}

%==============================================================
\section{Energy-first closed-pair construction}\label{app:closedpair}
%==============================================================

For self-containedness we give here the closed-pair construction used by
\cref{thm:pairwise-ito}: a Hilbert module of one-forms and a closed
derivative/divergence pair, in the abelian carr\'e-du-champ setting. The
construction below is the abelian projection of the Cipriani--Sauvageot
Hilbert-bimodule derivation~\citep{CiprianiSauvageot}; for the Dirichlet-space
and energy-module background see also~\citep{BouleauHirsch,FOT}; the
functional-analytic facts (closed operators, adjoints, $T^{**}=T$) are
standard~\citep{Rudin-FA,Schmuedgen}.

\begin{assumption}[The carr\'e-du-champ axioms]\label{ass:gamma-axioms}
$\cA$ is a unital subalgebra of $L^\infty(\Omega)\cap L^2(\Omega)$ and
$\Gamma\colon\cA\times\cA\to L^1(\Omega)$ a bilinear map with:
\textnormal{(Sym)} $\Gamma(F,G)=\Gamma(G,F)$;
\textnormal{(Vanish)} $\Gamma(c,\cdot)=0$ for constants $c$;
\textnormal{(MatPos)} $\sum_{i,j}\E[U_iU_j\,\Gamma(F_i,F_j)]\ge0$ for all finite
$(F_i,U_i)$ in $\cA$;
\textnormal{(WMatPos)} the same with a nonnegative multiplier $Q\in\cA_+$
inserted, $\sum_{i,j}\E[Q\,U_iU_j\,\Gamma(F_i,F_j)]\ge0$;
\textnormal{(FormClos)} the form $\cE(F,G)=\E[\Gamma(F,G)]$ is closable on
$L^2(\Omega)$: $F_n\to0$ in $L^2$ and $\cE(F_n-F_m)\to0$ imply $\cE(F_n)\to0$.
\end{assumption}

\subsection*{Theorem A: the Hilbert module of one-forms.}
The algebraic one-forms are $\Omega^1_{\mathrm{alg}}:=\cA\otimes_\R\cA$
($F\otimes U$ standing for $U\,dF$), with the semi-inner product
\begin{equation}\label{eq:fold-semiIP}
  \Bigl\langle\textstyle\sum_iF_i\otimes U_i,\ \sum_jG_j\otimes V_j\Bigr\rangle
  :=\sum_{i,j}\E\bigl[U_iV_j\,\Gamma(F_i,G_j)\bigr].
\end{equation}
By \textnormal{(Sym)}+\textnormal{(MatPos)} this is symmetric, bilinear, and
positive semidefinite; let $\mathcal N_0=\{\omega:\langle\omega,\omega\rangle=0\}$
and let $\cU_\Gamma$ be the completion of
$\Omega^1_{\mathrm{alg}}/\mathcal N_0$ (a real Hilbert space), with $[F\otimes U]$
the class of $F\otimes U$. The seed is $D_0F:=[F\otimes1]$.

\begin{theorem}[Hilbert module of one-forms]\label{thm:foldA}
Under \cref{ass:gamma-axioms}:
\begin{enumerate}
\item[$(\mathrm i)$] the action $W\cdot[F\otimes U]:=[F\otimes(WU)]$ ($W\in\cA$)
is bounded, $\|W\cdot\omega\|_{\cU_\Gamma}\le\|W\|_\infty\|\omega\|_{\cU_\Gamma}$,
so $\cU_\Gamma$ is a bounded $\cA$-bimodule;
\item[$(\mathrm{ii})$] $D_0\colon\cA\to\cU_\Gamma$ is linear with the isometry
$\langle D_0F,D_0G\rangle_{\cU_\Gamma}=\E[\Gamma(F,G)]$, and $D_0c=0$ for
constants $c$;
\item[$(\mathrm{iii})$] $\{U\cdot D_0F:F,U\in\cA\}=\{[F\otimes U]\}$ is dense in
$\cU_\Gamma$;
\item[$(\mathrm{iv})$] $\langle W\cdot D_0F,D_0G\rangle_{\cU_\Gamma}
=\E[W\,\Gamma(F,G)]$ for $F,G,W\in\cA$.
\end{enumerate}
\end{theorem}
\begin{proof}
$(\mathrm{ii}),(\mathrm{iv})$ are \eqref{eq:fold-semiIP} with $U=V=1$
(resp.\ $U=W,V=1$); $(\mathrm{vanish})$ gives $\|D_0c\|^2=\E[\Gamma(c,c)]=0$.
$(\mathrm{iii})$: $\sum_iF_i\otimes U_i=\sum_iU_i\cdot D_0F_i$ in $\cU_\Gamma$.
For $(\mathrm i)$, with $Q:=\|W\|_\infty^2-W^2\in\cA_+$,
\[
  \|W\cdot\omega\|^2_{\cU_\Gamma}
  =\sum_{i,j}\E[W^2U_iU_j\Gamma(F_i,F_j)]
  =\|W\|_\infty^2\|\omega\|^2_{\cU_\Gamma}-\sum_{i,j}\E[Q\,U_iU_j\Gamma(F_i,F_j)]
  \le\|W\|_\infty^2\|\omega\|^2_{\cU_\Gamma}
\]
by \textnormal{(WMatPos)}; the action descends to the quotient since
$\omega\in\mathcal N_0\Rightarrow W\cdot\omega\in\mathcal N_0$.
\end{proof}

\subsection*{Theorem B: the closed derivative and its divergence.}
\begin{lemma}[\textnormal{(FormClos)} implies closability]\label{lem:foldclos}
Under \cref{ass:gamma-axioms} the seed $D_0\colon L^2(\Omega)\supset\cA\to\cU_\Gamma$
is closable.
\end{lemma}
\begin{proof}
If $F_n\to0$ in $L^2$ and $D_0F_n\to v$ in $\cU_\Gamma$, then $(D_0F_n)$ is
Cauchy, i.e.\ $\E[\Gamma(F_n-F_m,F_n-F_m)]=\|D_0(F_n-F_m)\|^2_{\cU_\Gamma}\to0$
by the isometry. \textnormal{(FormClos)} gives
$\|D_0F_n\|^2_{\cU_\Gamma}=\E[\Gamma(F_n,F_n)]\to0$, so $v=0$.
\end{proof}

\begin{theorem}[Construction of $(\DN,\divN)$ and integration by parts]\label{thm:foldB}
Under \cref{ass:gamma-axioms}, let $\DN:=\overline{D_0}$ be the closure of the
seed and $\divN:=\DN^{*}$ its Hilbert-space adjoint. Then:
\begin{enumerate}
\item[$(\mathrm i)$] $\DN\colon\Dom(\DN)\subset L^2(\Omega)\to\cU_\Gamma$ is
closed and densely defined (\cref{lem:foldclos});
\item[$(\mathrm{ii})$] $\divN=\DN^{*}\colon\Dom(\divN)\subset\cU_\Gamma\to
L^2(\Omega)$ is closed and densely defined, with
$\Dom(\divN)=\{u:F\mapsto\langle\DN F,u\rangle_{\cU_\Gamma}\text{ is }L^2\text{-continuous}\}$
(\citealp[Thm.~13.9]{Rudin-FA}; \citealp[Ch.~5]{Schmuedgen}; using $\DN^{**}=\DN$);
\item[$(\mathrm{iii})$] integration by parts holds:
\begin{equation}\label{eq:foldIBP}
  \E\bigl[F\,\divN(u)\bigr]=\langle\DN F,u\rangle_{\cU_\Gamma},
  \qquad F\in\Dom(\DN),\ u\in\Dom(\divN);
\end{equation}
\item[$(\mathrm{iv})$] for $u=D_0U$ with $U\in\cA$ and $F\in\cA$,
$\E[F\,\divN(D_0U)]=\E[\Gamma(F,U)]$ whenever $D_0U\in\Dom(\divN)$.
\end{enumerate}
\end{theorem}
\begin{proof}
$(\mathrm i)$ is \cref{lem:foldclos} and the definition of the closure;
$(\mathrm{ii})$ is the standard closed-operator/adjoint theory
\citep[Thm.~13.9]{Rudin-FA} with $\DN^{**}=\DN$ \citep[Ch.~5]{Schmuedgen};
$(\mathrm{iii})$ is the defining property of the adjoint; $(\mathrm{iv})$ uses
$\langle\DN F,D_0U\rangle_{\cU_\Gamma}=\langle D_0F,D_0U\rangle_{\cU_\Gamma}
=\E[\Gamma(F,U)]$ by the seed isometry of \cref{thm:foldA}$(\mathrm{ii})$.
\end{proof}

In the application of \cref{sec:main}, the construction is applied
sectorwise (\cref{rem:sectorwise}): on each fixed one- or two-point sector
$\cA=C_b^\infty(\R)$ resp.\ $C_b^\infty(\R^2)$ (a genuine unital algebra),
$\Gamma=\GamN$ restricted to that sector, and $\cU_\Gamma=\cU_{\GamN}$; the four
algebraic axioms are immediate for $\GamN$ from the square-of-sum structure
\eqref{eq:sos} (\cref{sec:recap}), and the global pairwise object is the
sectorwise family of these constructions (\cref{rem:sectorwise}), not a
single global operator. \textnormal{(FormClos)} on each fixed $\le2$-point sector
is supplied by \cref{app:pairclos} (\cref{thm:sectorclos}), which uses no algebra
structure.

%==============================================================
\section{Pairwise closability criterion}\label{app:pairclos}
%==============================================================

We give the pairwise ($m\le2$) instance of the projective-limit closability
criterion; the full multi-level machinery is not needed here. Fix
$T>0$ and recall the pairwise sub-core
$\Apair=\{F_{\mathbf x,T,\Phi}:|\mathbf x|\le2\}$, stable under coordinate
deletion $2\to1\to0$. For a tuple $\mathbf x$ ($|\mathbf x|\le2$) write
$\mu_{\mathbf x}$ for the time-$T$ law of $Z_{\mathbf x}$,
$\cE_{\mathbf x}(\Phi)=\int\Gamma^{(|\mathbf x|)}_\theta(\Phi,\Phi)\,d\mu_{\mathbf x}$,
and for $\mathbf x\subseteq\mathbf y$ ($|\mathbf y|\le2$) let
$P^{\mathbf y}_{\mathbf x}=\E[\,\cdot\mid\sigma(Z_{\mathbf x})]$ be the
coordinate-deletion conditional expectation. The three conditions are:
\begin{itemize}
\item[\textnormal{(P1)}] $(\cE_{\mathbf x},\cD_{\mathbf x})$ is closable on
$L^2(\mu_{\mathbf x})$ for $|\mathbf x|=1,2$;
\item[\textnormal{(P2)}] $\cE_{\mathbf y}(\Phi\circ\pi_{\mathbf y,\mathbf x})
=\cE_{\mathbf x}(\Phi)$ for $\mathbf x\subseteq\mathbf y$;
\item[\textnormal{(P3)}] $\cE_{\mathbf x}(P^{\mathbf y}_{\mathbf x}\Phi)
\le\cE_{\mathbf y}(\Phi)$ for the $2\to1$ deletion.
\end{itemize}

\begin{lemma}[Conditional expectation contracts energy; factor lemma, pairwise instance]\label{lem:foldfactor}
Let $(\cE_{\mathbf y},\cD_{\mathbf y})$, $(\cE_{\mathbf x},\cD_{\mathbf x})$ be
closed Markovian symmetric Dirichlet forms on $L^2(\mu_{\mathbf y})$,
$L^2(\mu_{\mathbf x})$ ($\mathbf x\subset\mathbf y$ the $2\to1$ deletion,
$\mu_{\mathbf x}=(\pi_{\mathbf y,\mathbf x})_*\mu_{\mathbf y}$) whose semigroups
intertwine, $T^{\mathbf y}_t(g\circ\pi)=(T^{\mathbf x}_t g)\circ\pi$. Then
$\cE_{\mathbf x}(P^{\mathbf y}_{\mathbf x}\Phi)\le\cE_{\mathbf y}(\Phi)$ for all
$\Phi\in\cD_{\mathbf y}$; i.e.\ \textnormal{(P3)} holds. (This is the
conditional-expectation contraction of a Dirichlet form, in the manner
of~\citep[Thm.~1.6.1]{FOT}.)
\end{lemma}
\begin{proof}
Identifying $L^2(\mu_{\mathbf x})$ with $L^2(\mu_{\mathbf x})\circ\pi\subset
L^2(\mu_{\mathbf y})$, the conditional expectation is the orthogonal projection
$\Pi$ onto this subspace. Intertwining makes $\Pi$ commute with $T^{\mathbf y}_t$,
hence with $Q_t:=t^{-1}(I-T^{\mathbf y}_t)$; orthogonal-projection contraction
gives $\langle Q_t\Pi\Phi,\Pi\Phi\rangle\le\langle Q_t\Phi,\Phi\rangle$, and
letting $t\downarrow0$ (with $\langle Q_t u,u\rangle\nearrow\cE(u)$,
\citep[Lem.~1.3.4]{FOT}) yields $\cE_{\mathbf x}(P^{\mathbf y}_{\mathbf x}\Phi)
\le\cE_{\mathbf y}(\Phi)$.
\end{proof}

The lemma is the abstract template, recorded for orientation; its symmetry
and intertwining assumptions are \emph{not} claimed for the time-$T$-law
forms of \cref{rem:reference-measure}, which is exactly the
reference-measure subtlety, and the operative route to \textnormal{(P3)} in
this paper is the explicit assembly under $(\mathrm{CKC}_2^{\flat})$
(\cref{thm:deletion}).

We separate the two roles that the projective conditions play, because they are
often conflated. Fixed-sector closability is supplied by (P1) alone; the
contraction (P3)/$(\mathrm{CKC}_2^{\flat})$ supplies the $2\to1$
deletion-compatibility of the sectorwise pairs, not their closability.

\begin{theorem}[Fixed-sector closability, from \textnormal{(P1)} alone]\label{thm:sectorclos}
Under \cref{conv:LJR}, fix a tuple $\mathbf x$ with $|\mathbf x|\le2$. Under
\textnormal{(P1)} at level $|\mathbf x|$, the lifted single-time energy form
$\Ecyl$ restricted to the fixed sector $C_b^\infty(\R^{|\mathbf x|})$ is closable
on $L^2(\mu_{\mathbf x})$ (equivalently, \cref{hyp:beta1} holds on that sector).
No \textnormal{(P3)} is required: within a fixed sector the deletion projection is
the identity, so closability is exactly the sector-form closability
\textnormal{(P1)}.
\end{theorem}
\begin{proof}
Work within the fixed sector: let $(F_n)\subset C_b^\infty(\R^{|\mathbf x|})$, all
on the tuple $\mathbf x$, with $F_n\to0$ in $L^2(\mu_{\mathbf x})$ and
$\Ecyl(F_n-F_m)\to0$. Every element lies on the same fixed tuple $\mathbf x$, so
the projection onto that tuple acts as the identity, $P_{\mathbf x}F_n=F_n$, and
the deletion contraction (P3) plays no role here (it would hold with
equality). Hence $\Ecyl(F_n)=\cE_{\mathbf x}(F_n)$ throughout, with
$\cE_{\mathbf x}(F_n-F_m)\to0$ and $F_n\to0$ in $L^2(\mu_{\mathbf x})$; by (P1) at
level $|\mathbf x|$, $\cE_{\mathbf x}(F_n)\to0$, i.e.\ $\Ecyl(F_n)\to0$. \qed
\end{proof}

\begin{theorem}[Pairwise deletion-compatibility, from $(\mathrm{CKC}_2^{\flat})$]\label{thm:deletion}
Under $(\mathrm{CKC}_2^{\flat})$, the conditional expectation deleting one
coordinate is energy-contracting on the diagonal-flat sector,
\[
  \cE_{\mathbf x}(P^{\mathbf y}_{\mathbf x}\Phi)\;\le\;\cE_{\mathbf y}(\Phi),
  \qquad \Phi\in\cD^{\flat}_{\mathbf y}\quad(\mathbf x\subset\mathbf y\text{ the }2\to1\text{ deletion}),
\]
which is \textnormal{(P3)} there, so the one-point and two-point
fixed-sector closed pairs of \cref{thm:sectorclos} are compatible under the
$2\to1$ deletion on that sector. This is a statement about
deletion-compatibility, not about closability of either sector.
\end{theorem}
\begin{proof}
Two routes exist and carry different assumptions, which we keep separate. If
the deletion projection intertwined the relevant semigroups in the sense of
\cref{lem:foldfactor}, the contraction would be the abstract
conditional-expectation contraction of a Dirichlet form, in the manner
of~\citep[Thm.~1.6.1]{FOT}; but the forms here are integrated against the
time-$T$ laws, not against symmetrizing measures, and the symmetry and
intertwining assumptions of that route are exactly what
\cref{rem:reference-measure} flags as unavailable. The operative route is
therefore the explicit reduction of the body: by
\cref{prop:amended-star} and the trace inequalities \eqref{eq:traces}, the
deletion deficit $D(\Phi)=\cE_{\mathbf y}(\Phi)-\cE_{\mathbf x}(P\Phi)$ is
bounded below by the assembly of \cref{app:assembly}, restricted to the
diagonal-flat sector where the moving-atom term is rotated into the
center-of-mass direction and absorbed by the diagonal energy
(\cref{cor:neumann-subcore}); under $(\mathrm{CKC}_2^{\flat})$, whose
margin requirement $\kappa_\xi>0$ is part of the stated condition (and which, by
\cref{prop:invariant-fails}, genuinely excludes the translation-invariant
reference), the assembled lower bound is nonnegative, $D(\Phi)\ge0$ for
all $\Phi$ in the sector, which is \textnormal{(P3)} there. \qed
\end{proof}

Thus fixed-sector closability is \textnormal{(P1)} (\cref{thm:sectorclos}), while
$(\mathrm{CKC}_2^{\flat})$/\textnormal{(P3)} supplies the $2\to1$
deletion-compatibility of the sectorwise closed pairs (\cref{thm:deletion}), not
closability. The full linear span $\Apair$ of all pairwise sectors lies outside the
criterion: across sectors, $P_{\mathbf x}$ of a functional on a different
pair $\mathbf y$ is a conditional expectation involving the joint law of
$\mathbf x\cup\mathbf y$, up to three and four points ($m\ge3$), neither assumed
nor proved here. The closed-pair operators $(\DN,\divN)$ of \cref{thm:pairwise-ito}
are correspondingly built per fixed sector (\cref{rem:sectorwise}); there is
no single global operator and no global pairwise algebra.

%==============================================================
\section{Sharpness and remaining obstructions}\label{sec:sharpness}
%==============================================================

This appendix collects the supporting results that sharpen, without closing,
the two remaining structural fronts: the full reflecting case beyond the
diagonal-flat sector, and the $m\ge3$ compatibility. Together with the coalescing failure (\cref{prop:co-failure}) and the two-well
tension (\cref{rem:two-well}), they delimit how far the three main results can
be pushed with the present tools.

\subsection{The moving-atom trace is a theorem on the second-order sub-core}

\begin{proposition}[Agmon discharge of \cref{hyp:dtrace} in phase (R)]\label{prop:agmon}
Assume phase \textnormal{(R)}. Fix $\xi$ with
$P_0(\xi):=\inf_{|\eta-\xi|\le\delta}p(\xi,\eta)>0$ (such $\delta>0$ exists for
$\mu_\xi$-a.e.\ $\xi$ by \cref{lem:checkB}\textnormal{(R)}). Then for every
$\Phi\in C_b^\infty(\R^2)$,
\[
  |\partial_\eta\Phi(\xi,\xi)|^2\;\le\;
  C_{\mathrm{tr}}(\xi)\,\Bigl[\,G_\eta(\xi)
  +\bigl\langle(\partial_\eta^2\Phi)^2\bigr\rangle_p(\xi)\Bigr],
  \qquad
  C_{\mathrm{tr}}(\xi)\;=\;\frac{1}{P_0(\xi)}\max\Bigl(\frac2\delta,\;\delta\Bigr).
\]
In particular \cref{hyp:dtrace} holds in phase \textnormal{(R)} on the
second-order sub-core, with the enhanced norm
$\mathcal E_{\eta,\xi}(\Phi):=G_\eta(\xi)+\langle(\partial_\eta^2\Phi)^2\rangle_p(\xi)$
and the explicit constant above: $\tau_{\mathrm{at}}(\xi)\le C_{\mathrm{tr}}(\xi)<\infty$.
\end{proposition}
\begin{proof}
One-dimensional Agmon bound: for $h\in C^1([0,\delta])$ and $s\in[0,\delta]$,
$h(0)^2=h(s)^2-\int_0^s(h^2)'\le h(s)^2+2\int_0^\delta|h\,h'|$; averaging over
$s\in[0,\delta]$ and using $2AB\le\delta^{-1}A^2+\delta B^2$,
\[
  h(0)^2\;\le\;\frac1\delta\int_0^\delta h^2
  +2\Bigl(\int_0^\delta h^2\Bigr)^{1/2}\Bigl(\int_0^\delta h'^2\Bigr)^{1/2}
  \;\le\;\frac2\delta\int_0^\delta h^2+\delta\int_0^\delta h'^2 .
\]
Apply this to $h(s):=\partial_\eta\Phi(\xi,\xi+s)$, so $h(0)=\partial_\eta\Phi(\xi,\xi)$,
$h'=\partial_\eta^2\Phi$, and convert Lebesgue to weighted integrals by
$\int_0^\delta(\cdot)\,ds\le P_0(\xi)^{-1}\int(\cdot)\,p(\xi,\xi+s)\,ds
\le P_0(\xi)^{-1}\langle\cdot\rangle_p$.
\end{proof}

\begin{remark}[Why this does not yet give the full reflecting theorem: the order mismatch]\label{rem:order-mismatch}
\Cref{prop:agmon} converts \cref{hyp:dtrace} from condition to theorem
on the $H^2$ sub-core, with $C_{\mathrm{tr}}\sim1/\inf p$ exactly as
anticipated in \cref{rem:dtrace-gamma2}, and it degrades as $p\to0$ at the
coalescing endpoint, consistently with \cref{thm:wall-endpoint}. What it does
not give is the deletion contraction \textnormal{(P3)} for the
first-order forms beyond the flat sector: in the assembly
(\cref{app:assembly}), the $\tau_{\mathrm{at}}$-payments now draw on the
second-order budget $\langle(\partial_\eta^2\Phi)^2\rangle_p$, which the
first-order energy $\cE_{\mathbf y}$ does not contain. Under the value-part of
$\CKC$ and $\Lambda:=\operatorname*{ess\,sup}_\xi
q(\xi)\bigl(q(\xi)+A\bigr)C_{\mathrm{tr}}(\xi)<\infty$ one obtains, by the
existing assembly, only the mixed deletion inequality
\[
  \cE_{\mathbf x}(P\Phi)\;\le\;\cE_{\mathbf y}(\Phi)
  +\Lambda\!\int\!\bigl\langle(\partial_\eta^2\Phi)^2\bigr\rangle_p\,d\mu_\xi ,
  \qquad\Phi\in C_b^\infty(\R^2),
\]
first-order out, second-order in. A single-form \textnormal{(P3)}, for either
the first-order or the enhanced pair, requires the conditional expectation to
intertwine with the second-order structure; this is exactly the open content
of \cref{prob:fullreflecting}, now isolated as the only missing
element (the trace itself no longer being one).
\end{remark}

\subsection{The polarity route to a three-point failure is blocked}

The conjecture of \cref{prob:m3} rested on codimension-two polarity: at $m=3$
the triple apex has two $\Gamma$-blind difference directions, and a point in
two dimensions is polar for an elliptic Dirichlet integral, suggesting a probe
defeating \textnormal{(P3)}. This mechanism in fact fails, because the energy at
$m=3$ is not purely two-dimensional: the partial-coincidence strata carry
one-dimensional energies, and one-dimensional points have positive
capacity. We isolate the needed property of the (uncharacterized) three-point
law as a named condition.

\begin{hypothesis}[$(\mathrm R_3)$: pairwise-stratum regularity of the
three-point law]\label{hyp:r3}
For the three-point time-$T$ law from pairwise-distinct starts, in phase
\textnormal{(R)}, the partial-coincidence stratum
$\Delta_{23}=\{\eta_1=\eta_2\}$ carries the singular part of
$\mu^{(3)}_{\mathbf y}$ with cluster carr\'e coefficient $c_{\theta,2}>0$
\eqref{eq:sos}, and the pushforward of that stratum weight under the surviving
gap coordinate $w=\eta-\xi$ ($\eta$ the common position) has density bounded
below, $\rho_S(w)\ge\rho_0>0$ for $|w|\le\delta$.
\end{hypothesis}

This is the three-point analogue of \cref{lem:checkB}\textnormal{(R)}
(pairwise stickiness should populate the pairwise strata with positive density
near the triple apex); it concerns the $m=3$ law that \citep{LJR-FCN,LJR-circle-noise}
do not characterize, which is why it enters as a stated condition.

\begin{proposition}[The difference-direction witness cannot close at $m=3$]\label{prop:m3-blocked}
Assume phase \textnormal{(R)} and $(\mathrm R_3)$. Let
$\Phi=g(\eta_1-\xi,\eta_2-\xi)$ with $g\in C_b^\infty(\R^2)$, $g(0,0)=1$, be a
probe pinned at the triple apex, and suppose its restriction to the stratum,
$h(w):=g(w,w)$, satisfies $\int_{\delta/2\le|w|\le\delta}h^2\rho_S\,dw\le
\tfrac{1}{16}\rho_0\delta$ (as any \textnormal{(P3)}-defeating sequence must,
$\Phi_n\to0$ in $L^2(\mu^{(3)}_{\mathbf y})$ forcing $h_n\to0$ in
$L^2(\rho_S)$). Then
\[
  \cE^{(3)}_{\mathbf y}(\Phi)\;\ge\;\frac{c_{\theta,2}\,\rho_0}{8\,\delta}\;>\;0 .
\]
In particular no sequence of such probes can send
$\cE^{(3)}_{\mathbf y}(\Phi_n)\to0$ while remaining pinned at the apex: the
witness mechanism conjectured in \cref{prob:m3} is blocked by the
one-dimensional capacity of the pairwise strata.
\end{proposition}
\begin{proof}
On $\Delta_{23}$ the surviving cluster derivative of \eqref{eq:sos} acts on
$\Phi$ by $(\partial_{\eta_1}+\partial_{\eta_2})\Phi=
(\partial_1g+\partial_2g)(w,w)=h'(w)$, so
\[
  \cE^{(3)}_{\mathbf y}(\Phi)\;\ge\;\frac{c_{\theta,2}}2
  \int h'(w)^2\,\rho_S(w)\,dw .
\]
By the $L^2$-smallness assumption there is $w^*\in[\delta/2,\delta]$ with
$|h(w^*)|\le\tfrac12$: otherwise $|h|>\tfrac12$ on all of $[\delta/2,\delta]$
and $\int_{\delta/2\le|w|\le\delta}h^2\rho_S\ge\tfrac14\,\rho_0\,\tfrac\delta2
=\tfrac18\rho_0\delta>\tfrac1{16}\rho_0\delta$, a contradiction. Then, by
Cauchy--Schwarz with the resistance of $\rho_S$ on $[0,\delta]$,
\[
  \tfrac12\le|h(0)-h(w^*)|\le\int_0^{\delta}|h'|\,dw
  \le\Bigl(\int h'^2\rho_S\Bigr)^{1/2}\Bigl(\int_0^{\delta}\frac{dw}{\rho_S}\Bigr)^{1/2}
  \le\Bigl(\int h'^2\rho_S\Bigr)^{1/2}\Bigl(\frac{\delta}{\rho_0}\Bigr)^{1/2},
\]
so $\int h'^2\rho_S\ge\rho_0/(4\delta)$, and the stratum lower bound gives
$\cE^{(3)}_{\mathbf y}(\Phi)\ge\tfrac{c_{\theta,2}}2\cdot\rho_0/(4\delta)
=c_{\theta,2}\rho_0/(8\delta)$.
\end{proof}

%==============================================================
\bibliographystyle{alea3}
\bibliography{blackito}

\end{document}
